\chapter{从上电到程序运行：搭机、全链路 Trace 与故障定位}

\begin{keyidea}
最后一章把全书重新连成一台可观察的机器：从供电复位到取指执行，从主存和设备访问到输出与排错。目标不是只让程序运行，而是能解释每一个边界。
\end{keyidea}

前面二十五章把每个部件都单独讲过：逻辑门、触发器、加法器、寄存器堆、控制字、总线、存储器和调试方法。本章只做一件事：把这些部件按一张确定的总图接成一台真的能执行程序的 8 位教学计算机。机器很小——8 位数据总线、16 字节地址空间、11 条指令（含 NOP）——但小是刻意的：小到每个控制信号都能用 74HC 器件数出来，小到任何故障都能用 LED 和单步时钟定位。这样的机器在教学文献里常被称为 SAP-1 类结构，它的价值不在性能，而在于“没有一处不可观察”。

\begin{longtable}{L{0.2\linewidth}L{0.34\linewidth}L{0.36\linewidth}}
\toprule
对象 & 规格 & 实现要点 \\
\midrule
数据总线 & 8 位，三态共享 & 任何时刻只有一个驱动者 \\\tablerowrule
地址空间 & 16 字节（4-bit 地址） & 程序和数据同处一块 RAM \\\tablerowrule
寄存器 & A、B、IR、MAR、OUT、PC、Flags & 74HC273/173，带装载使能 \\\tablerowrule
ALU & 加法和减法 & 74HC283 加法器，减法是取反加一 \\\tablerowrule
标志 & Zero、Carry & 供条件跳转使用 \\\tablerowrule
控制 & 微码 ROM 译码 & opcode 加 T 状态索引微码 \\\tablerowrule
时钟 & 自动（555 振荡）和单步（按钮） & 调试时逐拍推进 \\
\bottomrule
\end{longtable}

指令集只有 11 条（含 NOP），高 4 位是 opcode，低 4 位是地址或立即数。每条指令占一个字节，程序最长 16 条（不含数据时；剩余字节放数据）。这套编码与本章微码 ROM 的内容一一对应，后面“控制单元：指令译码与微码”一节会逐条展开微操作。

\tablechunkintro{1}{2}{\code{NOP}}{\code{JMP addr}}
\begin{longtable}{L{0.12\linewidth}L{0.16\linewidth}L{0.3\linewidth}L{0.32\linewidth}}
\toprule
指令 & opcode & 含义 & 用到的主要路径 \\
\midrule
\code{NOP} & \code{0x0} & 空操作 & 只占周期 \\\tablerowrule
\code{LDA addr} & \code{0x1} & A ← RAM[addr] & MAR、RAM 读、A 装载 \\\tablerowrule
\code{ADD addr} & \code{0x2} & A ← A + RAM[addr] & B 装载、ALU、标志 \\\tablerowrule
\code{SUB addr} & \code{0x3} & A ← A − RAM[addr] & B 取反加一、ALU \\\tablerowrule
\code{STA addr} & \code{0x4} & RAM[addr] ← A & MAR、RAM 写 \\\tablerowrule
\code{LDI imm} & \code{0x5} & A ← imm & 总线低 4 位、A 装载 \\\tablerowrule
\code{JMP addr} & \code{0x6} & PC ← addr & PC 装载 \\
\bottomrule
\end{longtable}

\tablechunkintro{2}{2}{\code{JC addr}}{\code{HLT}}
\begin{longtable}{L{0.12\linewidth}L{0.16\linewidth}L{0.3\linewidth}L{0.32\linewidth}}
\toprule
指令 & opcode & 含义 & 用到的主要路径 \\
\midrule
\code{JC addr} & \code{0x7} & Carry=1 时 PC ← addr & 标志、PC 装载 \\\tablerowrule
\code{JZ addr} & \code{0x8} & Zero=1 时 PC ← addr & 标志、PC 装载 \\\tablerowrule
\code{OUT} & \code{0xE} & OUT ← A & 输出寄存器 \\\tablerowrule
\code{HLT} & \code{0xF} & 停机 & 时钟停止 \\
\bottomrule
\end{longtable}

\begin{keyidea}
教学计算机的全部秘密在“每一步可观察”：总线值用 LED 显示，时钟可以单步，控制线可以逐根核对。搭建顺序和调试顺序都围绕这个原则组织。
\end{keyidea}

\section{一、总体架构：总线、时钟与复位}

\topic{面包板布局按总线居中组织}整机里流量最大的结构是总线：取指时
PC 的值经它去 MAR，取回的字节经它去 IR，运算时数据经它在 RAM、A、B 与
ALU 之间流动，停机前最后一个动作还是经它把 A 的值送进 OUT。布局的第一
原则因此不是省面包板，而是让总线这条中轴最短、最直、最容易测量。具体
做法是把 8 根数据线排在面包板阵列的中央：可以用一整块面包板专走总线，
各功能模块分列上下两侧；也可以把几块面包板纵向拼成长条，让总线沿中央
沟槽两侧贯通到底。模块大体按数据流向排位：PC、MAR、RAM 居一侧，A、B、
ALU 居另一侧，IR 与控制单元靠一端，OUT 与显示排靠另一端。这样每个模块
到总线的距离相近，飞线短而且彼此平行，出了故障能顺着总线一根一根查，
而不是在一团交叉的杜邦线里猜。

先把整机的连接关系画成一张总图：总线居中，模块分列两侧，控制单元居
下。

\begin{pdfdiagram}
  % 中央 8-bit 总线（竖条）
  \draw[draw=BookMuted, line width=1.8pt] (0,3.25) -- (0,-2.5);
  \node[hw label] at (0,3.55) {8-bit 三态总线（空闲下拉读 0x00）};
  % 左列：PC、MAR→RAM、OUT
  \node[hw state, minimum width=1.6cm] (pc) at (-3.0,2.6) {PC\\（4 位）};
  \node[hw state, minimum width=1.3cm] (mar) at (-3.0,1.1) {MAR};
  \node[hw memory, minimum width=1.8cm] (ram) at (-3.0,-0.2) {RAM\\16×8};
  \node[hw state, minimum width=1.6cm] (out) at (-3.0,-1.7) {OUT\\显示};
  % 右列：IR、A、ALU、B
  \node[hw state, minimum width=1.4cm] (ir) at (3.0,2.75) {IR};
  \node[hw state, minimum width=1.4cm] (acc) at (3.0,1.7) {A};
  \node[hw compute, minimum width=1.7cm] (alu) at (3.0,0.45) {ALU\\加/减};
  \node[hw state, minimum width=1.4cm] (regb) at (3.0,-0.7) {B};
  % 控制单元居下
  \node[hw control, minimum width=6.4cm] (ctrl) at (0,-3.9) {控制单元\\IR 高 4 位 + T 计数器 + 微码 ROM};
  % 左列总线接口（五个驱动来源各标在出口）
  \draw[hw bus] ($(pc.east)+(0,0.12)$) -- (0,2.72) node[midway, above=1pt, hw label, fill=white, inner sep=1pt]{CO};
  \draw[hw bus] (0,2.48) -- ($(pc.east)+(0,-0.12)$) node[midway, below=1pt, hw label, fill=white, inner sep=1pt]{J};
  \draw[hw bus] (0,1.1) -- (mar.east) node[midway, above=1pt, hw label, fill=white, inner sep=1pt]{MI};
  \draw[hw bus] (mar.south) -- (ram.north) node[midway, right=1pt, hw label, fill=white, inner sep=1pt]{4 位地址};
  \draw[hw bus] ($(ram.east)+(0,0.12)$) -- (0,-0.08) node[midway, above=1pt, hw label, fill=white, inner sep=1pt]{RO};
  \draw[hw bus] (0,-0.32) -- ($(ram.east)+(0,-0.12)$) node[midway, below=1pt, hw label, fill=white, inner sep=1pt]{RI};
  \draw[hw bus] (0,-1.7) -- (out.east) node[midway, above=1pt, hw label, fill=white, inner sep=1pt]{OI};
  % 右列总线接口
  \draw[hw bus] (0,2.87) -- ($(ir.west)+(0,0.12)$) node[midway, above=1pt, hw label, fill=white, inner sep=1pt]{II};
  \draw[hw bus] ($(ir.west)+(0,-0.12)$) -- (0,2.63) node[midway, below=1pt, hw label, fill=white, inner sep=1pt]{IO 低 4 位};
  \draw[hw bus] ($(acc.west)+(0,0.12)$) -- (0,1.82) node[midway, above=1pt, hw label, fill=white, inner sep=1pt]{AO};
  \draw[hw bus] (0,1.58) -- ($(acc.west)+(0,-0.12)$) node[midway, below=1pt, hw label, fill=white, inner sep=1pt]{AI};
  \draw[hw bus] (alu.west) -- (0,0.45) node[midway, above=1pt, hw label, fill=white, inner sep=1pt]{$\Sigma$O};
  \draw[hw bus] (0,-0.7) -- (regb.west) node[midway, above=1pt, hw label, fill=white, inner sep=1pt]{BI};
  % A、B 到 ALU 的专线（不经总线）
  \draw[hw bus] (acc.south) -- (alu.north) node[midway, right=1pt, hw label, fill=white, inner sep=1pt]{A 专线};
  \draw[hw bus] (regb.north) -- (alu.south);
  % opcode 专线直下控制单元，CF/ZF 回报
  \draw[hw return] ($(ir.east)+(0,-0.15)$) -- (4.15,2.6) -- (4.15,-3.9) -- (3.2,-3.9);
  \node[hw label, fill=white, inner sep=1pt] at (4.0,2.15) {opcode};
  \draw[hw return] ($(alu.east)+(0,0.15)$) -- (4.6,0.6) -- (4.6,-3.6) -- (3.2,-3.6);
  \node[hw label, anchor=east, fill=white, inner sep=1pt] at (4.55,-1.0) {CF、ZF};
  % 16 根控制线：左右两条干线分到各模块
  \draw[hw return] ($(ctrl.west)+(0,0.15)$) -- (-4.6,-3.75) -- (-4.6,2.6);
  \draw[hw return] (-4.6,2.6) -- (pc.west);
  \draw[hw return] (-4.6,1.1) -- (mar.west);
  \draw[hw return] (-4.6,-0.2) -- (ram.west);
  \draw[hw return] (-4.6,-1.7) -- (out.west);
  \draw[hw return] ($(ctrl.east)+(0,0.15)$) -- (4.75,-3.75) -- (4.75,2.85);
  \draw[hw return] (4.75,2.85) -- ($(ir.east)+(0,0.15)$);
  \draw[hw return] (4.75,1.7) -- (acc.east);
  \draw[hw return] (4.75,0.3) -- ($(alu.east)+(0,-0.15)$);
  \draw[hw return] (4.75,-0.7) -- (regb.east);
  \fill[BookTeal] (-4.6,1.1) circle (0.05);
  \fill[BookTeal] (-4.6,-0.2) circle (0.05);
  \fill[BookTeal] (-4.6,-1.7) circle (0.05);
  \fill[BookTeal] (4.75,2.85) circle (0.05);
  \fill[BookTeal] (4.75,1.7) circle (0.05);
  \fill[BookTeal] (4.75,0.3) circle (0.05);
  \fill[BookTeal] (4.75,-0.7) circle (0.05);
  \node[hw label, rotate=90] at (-4.95,0.4) {16 根控制线};
  \node[hw label, rotate=90] at (5.1,0.4) {16 根控制线};
\end{pdfdiagram}
\diagramnote{整机总框图。中央竖条是 8 位三态总线；左列挂 PC、MAR→RAM、OUT，右列挂 IR、A、ALU、B，控制单元居下。实线是数据路径：五个总线驱动来源各标在出口处（CO、RO、AO、IO、$\Sigma$O），装载类控制线标在入口一侧；虚线是控制与状态线，16 根控制线经左右两条干线分到各模块。A、B 到 ALU 走不经总线的专线，IR 高 4 位的 opcode 专线直下控制单元，CF、ZF 沿虚线回报给条件跳转。任何一拍的“谁在说、谁在听”，都应能在这张图上逐根指出来。}

以常见的 830 孔面包板估量，整机大约要三到四块：一块专走总线与 LED 排，
一块放 A、B 与 ALU，一块放 RAM、MAR 与 OUT，控制单元与微码 ROM 单独占
一块或半块。时钟与复位电路最小，放在面板开关手边即可。宁可多留空排也
不要挤：模块之间留出的空隙，后面插测试钩、挪飞线、加扩展时都会用到。

8 根线之间还有一个位序约定：全机统一，例如 D7 在上、D0 在下，每个模块
按同一方向插接。位序接反不会烧坏任何东西，但字节会被镜像，\code{0x01}
变成 \code{0x80}——这种故障不冒烟、不发热，只表现为数值诡异，排查起
来格外费时，值得在插第一根线之前就用标签把方向定下来。顺手再给导线分
个颜色：红色电源、黑色地、黄色总线、蓝色控制线，后面任何一个模块都可
以凭颜色先猜出每根线的身份，猜错了再量，也比逐根追线快得多。

电源轨也要先规划。面包板两侧的长条电源轨常常在板子中间断开，上下两段
之间要用跳线跨接，全机所有板块的轨还要连成同一个 5V 和同一个地——共
地不是可选项，所有三态约定、所有电平判断都以同一参考地为准。每块面包
板的电源入口先放一只 10$\mu$F 电容，后面每插一片芯片再补一只 100nF。
电源走线的优先级高于信号线：先把轨连通、测过电压，再插第一片芯片。

每往总线上挂一个模块，都按同一组三问接线：它的数据口接总线哪几位，它
的装载由哪根控制线在什么沿触发，它的输出由哪根控制线使能。三个问题都
有确定答案，模块才算真正挂上了总线；有一个含糊，就先别插。A、B、IR、
RAM、ALU 的这三问答案各不相同，“寄存器与 ALU 模块”一节、“内存与输出模块”一节会逐个给出，但问法是全机统一
的。

总线是共享资源，共享的接线约定只有一条：任何时刻至多一个驱动者。全机
能驱动总线的来源只有五个，各对应一根输出类控制线：CO（PC 输出）、RO
（RAM 输出）、AO（A 输出）、IO（IR 低 4 位输出）、$\Sigma$O（ALU 输
出）。微码保证同一拍这五根线至多一根有效，“控制单元：指令译码与微码”一节会逐拍列出；但只靠微码
正确还不够，每个驱动者的输出都必须经过三态门：输出使能由对应控制线驱
动，控制线一撤，输出立即回到高阻，把总线让出来。74HC173 自带三态输出，
74HC273 则要外接 74HC244 之类的三态缓冲，“寄存器与 ALU 模块”一节会给出两种接法。

把“谁能说、说什么”先列成一张表，接线时照表核对：

\begin{longtable}{L{0.22\linewidth}L{0.14\linewidth}L{0.54\linewidth}}
\toprule
驱动来源 & 控制线 & 有效时总线上出现的内容 \\
\midrule
PC（低 4 位） & CO & 当前指令地址，总线高 4 位无人驱动 \\\tablerowrule
RAM & RO & MAR 选中单元的字节 \\\tablerowrule
A 寄存器 & AO & 累加器当前值 \\\tablerowrule
IR 低 4 位 & IO & 指令的地址或立即数字段，高 4 位无人驱动 \\\tablerowrule
ALU & $\Sigma$O & 加法器当前结果（A+B 或 A$-$B） \\
\bottomrule
\end{longtable}

约定的另一半是“没轮到你说话时必须沉默”。上电瞬间所有输出类控制线还
没有意义，因此每个三态使能的默认电平都必须是无效——这要靠接线保证，
不能指望程序小心驾驶。还有一个容易漏掉的角落：五个来源都不驱动时，总
线本身悬空，LED 可能半亮，挂在总线上的 CMOS 输入长期停在中间电平还会
额外耗电发热。给 8 根总线各接一只 100k$\Omega$ 下拉电阻，空闲拍总线稳
定读出 \code{0x00}。这既给高阻节点一个默认归宿，也顺手送了一个调试提
示：某拍本该有驱动者，总线却读 \code{0x00}，先查那根输出控制线到场了
没有。

总线正中央还值得留一整排位置给 8 只 LED：每根数据线经一只限流电阻
（330$\Omega$ 到 1k$\Omega$，按亮度调整）接一只 LED，常年点亮。这排灯
占的地方不小，换来的是整机唯一的实时读数：单步执行时，每一拍总线上是
谁的值、是不是预期的值，低头就能看到。后面调试的大多数问题——该 RO
的拍没人驱动、取指取回了错字节、ALU 结果没上总线——第一反应都是看这
排灯。本章开头说教学机的目标是“没有一处不可观察”，总线 LED 排就是这
条原则最大的一笔投资：一排面包板位置，换全机最繁忙结构的全程可见。

\begin{warningbox}
两只推挽输出同时驱动同一根线、且一高一低时，电流只受两级导通电阻限制：
轻则逻辑电平被夹在中间、芯片发热，重则损坏输出级。争用不会自动报错，
它表现为“偶尔算错”“某片芯片发烫”这类离病根很远的怪现象。所以“同
一时刻至多一个驱动者”这条约定要用接线来保证——不用的三态使能固定在
无效电平——而不是只靠微码写得对。
\end{warningbox}

\topic{时钟分自动和单步两路}全机所有触发器共用一棵时钟树，树根只有一
个。但运行和调试对时钟的要求正好相反：运行时希望它自己往前走，调试时
希望它一步一停、每步都可检查。教学机的解决办法是做两路时钟源头，用一
个开关选一路：自动挡由 555 定时器产生连续方波，单步挡由一个经过消抖的
按钮产生人工脉冲，两路经选择器会合后再送往时钟树。

两路时钟从振荡源到时钟树一共四级，先画出结构再逐件展开。

\begin{pdfdiagram}
  % 两路源头
  \node[hw compute, minimum width=2.4cm] (osc) at (0,1.2) {555 非稳态\\0.7--48Hz};
  \node[hw compute, minimum width=2.4cm] (step) at (0,-0.8) {按钮 + RC 消抖\\74HC14 整形};
  % 选择器
  \node[hw compute, minimum width=1.8cm] (mux) at (3.4,0.2) {74HC157\\二选一};
  \node[hw label] (sel) at (3.4,1.95) {面板开关：自动/单步};
  \draw[hw return] (sel.south) -- (mux.north);
  % HLT 钳停
  \node[hw compute, minimum width=1.8cm] (clamp) at (6.4,0.2) {与门\\HLT 钳停};
  \node[hw label] (hlt) at (6.4,-1.65) {HLT 有效 → 输出钳 0};
  \draw[hw return] (hlt.north) -- (clamp.south);
  % 时钟树
  \node[hw state, minimum width=2.0cm] (tree) at (9.5,0.2) {时钟树\\星型分送全机};
  % 走线
  \draw[hw bus] (osc.east) -- ++(0.6,0) |- ($(mux.west)+(0,0.15)$);
  \draw[hw bus] (step.east) -- ++(0.4,0) |- ($(mux.west)+(0,-0.15)$);
  \draw[hw bus] (mux.east) -- node[midway, above=1pt, hw label, fill=white, inner sep=1pt]{clk\_src} (clamp.west);
  \draw[hw bus] (clamp.east) -- node[midway, above=1pt, hw label, fill=white, inner sep=1pt]{clk\_tree} (tree.west);
\end{pdfdiagram}
\diagramnote{时钟两路。自动挡 555 非稳态输出 0.7--48Hz 连续方波，单步挡按钮经 RC 消抖和 74HC14 施密特整形给出干净脉冲；74HC157 按面板开关二选一，选中一路再经与门——HLT 有效时输出被钳在低电平，时钟树冻结，机器停在当前拍。时钟树从与门输出星型分到全机所有触发器的时钟脚，不菊花链。}

\begin{longtable}{L{0.14\linewidth}L{0.32\linewidth}L{0.22\linewidth}L{0.22\linewidth}}
\toprule
路径 & 核心器件 & 输出特征 & 适用场合 \\
\midrule
自动挡 & 555 非稳态，R1、R2、C 定时 & 0.7--48Hz 连续方波（C=10$\mu$F） & 连续运行程序 \\\tablerowrule
单步挡 & 按钮、RC 消抖、74HC14 整形 & 按一下一个干净脉冲 & 逐拍调试 \\\tablerowrule
选择器 & 74HC157 一路二选一 & 面板开关选源，HLT 可钳停 & 两路会合进时钟树 \\
\bottomrule
\end{longtable}

自动挡用 555 的非稳态接法：电阻 R1 接在 VCC 与放电脚之间，R2 接在放电
脚与阈值脚之间，电容 C 从触发脚到地，输出脚直接出方波，频率
$f \approx 1.44/((R1+2R2)\cdot C)$。取 R1=1k$\Omega$、R2 用一只
100k$\Omega$ 电位器串联 1k$\Omega$ 固定电阻、C=10$\mu$F：电位器拧到最
小（R2$\approx$1k$\Omega$）时 $f \approx 48$Hz，拧到最大
（R2$\approx$101k$\Omega$）时 $f \approx 0.7$Hz——一只旋钮覆盖 1 到
50Hz，常用调试区间都在里面；把 C 换成 1$\mu$F，整个范围上移十倍。占空
比为 $(R1+R2)/(R1+2R2)$：低速端接近 1/2，高速端约 2/3。这个不对称对本
机无碍：最窄的半拍也是毫秒量级，而 74HC 器件要求的最小时钟脉宽只有几
十纳秒，余量有几十万倍。

为什么不干脆跑快些？面包板加人眼给时钟频率划了上限：LED 每拍都要看得
清，排查时每拍都要停得住，1 到 10Hz 是观察最舒服的区间。教学机追求的
不是快，而是慢到每一步都可观察；想验证整机在更高频率下是否仍然正确，
等程序全部调通以后再把电位器拧上去也不迟。

单步挡只有一个按钮：按一下，机器走一拍。麻烦在于机械触点会在几毫秒内
反复通断，把触点信号直接当时钟用，“按一下”会被计数器看成好几拍。消
抖用 RC 加施密特整形：按钮一端接地，另一端经 10k$\Omega$ 电阻上拉到
VCC，同一节点再经 1$\mu$F 电容到地——时间常数 10ms，足以盖过典型的抖
动窗口；该节点送 74HC14 施密特反相器整形，输出才是干净的单步脉冲。
“按一下、只跳变一次”这条性质，后面每一步调试都要依赖，所以消抖不是
可选项。

两路时钟经一片 74HC157（四路二选一，只用其中一路）选择，选择端接面板
上的拨动开关；选中一路再与停机逻辑会合——HLT 有效时把时钟树输入钳在
低电平，机器停在当前拍。用伪码写就是：

\begin{lstlisting}
clk_src  = sel ? clk_555 : clk_step   -- 面板开关选一路
clk_tree = halt ? 0 : clk_src         -- HLT 有效时钳住时钟树
\end{lstlisting}

切换选择开关的瞬间可能切出窄毛刺，所以约定只在停机或按住复位按钮时换挡。时钟树从选择器输出这一点出发，星型分到各模块的时钟脚，不要菊花链
式地一个模块串一个模块——每根时钟线末端的边沿质量，决定那个模块的状
态边界是否真实存在。时钟模块自身的检查靠眼睛：自动挡拧到 1Hz 左右，总
线 LED 应约每秒整体变一次；单步挡按一下，LED 变一次且只变一次。按一下
变了两次，查消抖；自动挡完全不动，先量 555 的输出脚。

两路时钟的分工要养成一个硬习惯：任何程序第一次上板，一律从单步开始。
逐拍走完第一条指令——看 T0 拍 PC 的值上总线、T1 拍指令字节进 IR、T2
拍以后各条控制线按预期起落——确认无误后再拨到自动挡连续运行。自动挡
是“跑”，单步挡是“看”；没看过的程序直接跑，出了错连错误发生在第几
拍都不知道，只能从头再来。调试的全部增量信息都来自单步，自动挡只负责
最后确认。

\begin{classicnote}
经典教材讲数据通路时，时钟和复位总是作为第一批元件画进图里，而不是最
后补上的配角：没有确定的时钟边界就没有“状态”，没有确定的复位就没有
“初态”。本章的两路时钟加复位链，不过是把这两条经典要求做进面包板
的电源轨之间。
\end{classicnote}

\topic{复位把 PC、T 状态和控制状态清到已知值}上电瞬间，每一片触发器
的输出都是随机的：PC 可能是 0 也可能是 13，T 状态计数器可能停在 T5，
IR 里是一个说不上名字的字节。不复位就直接运行，第一条被执行的“指令”
来自随机地址、从随机 T 状态进入微码，同一块板子上电十次可能有十种开局。
程序的第一条性质不是正确而是可复现，可复现的前提是每次上电都从同一个
已知状态出发。复位链的任务，就是把这个前提实实在在地做出来。

必须清掉的状态只有三类。一是 PC：程序计数器清 0，第一条指令永远从地址
0 取。二是 T 状态计数器：清 0 后机器从 T0 开始，而 T0、T1 是公共取指
拍，微码 ROM 里所有 opcode 对应的 T0 行都填同一个字 \code{CO+MI}，T1 行
都填 \code{RO+II+CE}（“控制单元：指令译码与微码”一节会展开）——所以 IR 的随机初值无害：T1 拍 II
有效，RAM[0] 的字节会把 IR 整体覆盖。三是 Flags：程序并不保证第一次
JC/JZ 之前一定执行过 ADD 或 SUB，标志寄存器必须接复位，否则开头几次条
件跳转读的是随机值。A、B 的内容不影响执行顺序，程序第一条 LDA 或
LDI 自然会给 A 一个确定值，把它们接进复位只是让上电后的读数确定；OUT
用的 74HC377 没有清除端，不接复位，上电读数任意，第一条 OUT 执行后自
然被覆盖，不影响执行。

复位信号由复位按钮经 RC 消抖与施密特整形产生，链路与单步时钟相同。使
用时遵守“异步拉入、同步释放”。按下按钮：复位直接进各片的异步清零脚
（74HC273 的 $\overline{CLR}$、74HC161 的 $\overline{CLR}$ 等），不等时
钟沿、立即生效——上电过程中时钟还不干净，复位不能反过来依赖时钟。松
开按钮：释放沿先经过一级触发器与时钟对齐，再撤掉各片的清零。具体做法
是把消抖后的复位信号接一片 74HC74 的 D 端、系统时钟接它的时钟脚，用 Q
去驱动各片的清零脚，全机就在同一个时钟沿之后一起离开复位，不会出现一
半模块先走、一半还停着的局面。释放后第一个上升沿执行的就是 T0：
\code{CO+MI} 把 PC=0 送进 MAR。机器的第一拍从此永远相同，这是后面所有
调试能成立的地基。

把从按下到松开的全过程画成时序，拉入与释放的不对称直接可见。

\begin{waveform}[x=0.85cm,y=0.9cm]
  % 参考虚线：拉入时刻（x=2.5，不在时钟沿上）与释放时刻（x=9，上升沿）
  \draw[BookMuted, dashed, line width=0.5pt] (2.5,1.4) -- (2.5,9.2);
  \draw[BookMuted, dashed, line width=0.5pt] (9,1.4) -- (9,9.2);
  \node[hw label, anchor=west] at (2.65,6.0) {拉入：不等时钟沿};
  \node[hw label, anchor=east] at (8.85,4.4) {释放：与上升沿对齐};
  % CLK：上升沿在 1/3/5/7/9
  \draw[BookAccent, line width=0.9pt] (0,8.2) -- (1,8.2) -- (1,9.0) -- (2,9.0) -- (2,8.2) -- (3,8.2) -- (3,9.0) -- (4,9.0) -- (4,8.2) -- (5,8.2) -- (5,9.0) -- (6,9.0) -- (6,8.2) -- (7,8.2) -- (7,9.0) -- (8,9.0) -- (8,8.2) -- (9,8.2) -- (9,9.0) -- (10,9.0) -- (10,8.2) -- (11,8.2);
  \node[hw label, anchor=east] at (-0.2,8.6) {CLK};
  % 复位按钮：按下与松开瞬间有弹跳
  \draw[BookAccent, line width=0.9pt] (0,7.2) -- (2,7.2) -- (2,6.4) -- (2.12,6.4) -- (2.12,7.2) -- (2.24,7.2) -- (2.24,6.4) -- (8,6.4) -- (8,7.2) -- (8.12,7.2) -- (8.12,6.4) -- (8.24,6.4) -- (8.24,7.2) -- (11,7.2);
  \node[hw label, anchor=east] at (-0.2,6.8) {复位按钮};
  % 消抖后复位：干净电平，2.5 拉低、8.5 释放
  \draw[BookAccent, line width=0.9pt] (0,5.6) -- (2.5,5.6) -- (2.5,4.8) -- (8.5,4.8) -- (8.5,5.6) -- (11,5.6);
  \node[hw label, anchor=east] at (-0.2,5.2) {消抖后复位};
  % 74HC74 的 Q：异步清零立即拉低，释放等 x=9 的上升沿
  \draw[BookAccent, line width=0.9pt] (0,4.0) -- (2.5,4.0) -- (2.5,3.2) -- (9,3.2) -- (9,4.0) -- (11,4.0);
  \node[hw label, anchor=east] at (-0.2,3.6) {同步器 Q};
  % 各片清零脚：跟随 Q
  \draw[BookAccent, line width=0.9pt] (0,2.4) -- (2.5,2.4) -- (2.5,1.6) -- (9,1.6) -- (9,2.4) -- (11,2.4);
  \node[hw label, anchor=east] at (-0.2,2.0) {各片清零脚};
\end{waveform}
\diagramnote{复位链时序：异步拉入、同步释放。按钮按下（电平有弹跳）后，消抖复位接在 74HC74 的异步清零端，Q 与各片清零脚立即变低，不等待时钟沿；按钮松开后，消抖复位早已接在 74HC74 的 D 端，Q 要等下一个时钟上升沿才变高，各片在同一沿之后一起离开复位，第一个装载沿执行的就是 T0 的 CO+MI。}

复位按钮不只在上电时用。调试中途程序走乱了、想换一个初始条件、或者只
是想再看一遍开头几拍，按一下复位就回到已知起点：PC 归 0、T 归 0、标
志清零，而 RAM 的内容不受复位影响——程序还在，随时可以重新单步。复位
因此是面板上使用频率最高的键，它的地位相当于调试器里的“回到开头”；
把它装在顺手的位置，比把显示排摆得好看更实际。

正式产品里常用 RC 加电压监控做成上电自动复位，让机器一得电就自己进入已
知状态。教学机用一个手动按钮已经足够：上电后先按一下复位，效果相同。
要知道的只是差别——自动复位保证的是“没有人按也确定”，手动按钮保证
的是“按过就确定”，调试场景里后者反而更方便，因为复位时机由你掌握。

\begin{warningbox}
“程序很短，不复位多半也跑得对”是一种错觉。随机状态意味着十次里可能
九次碰巧正常，第十次 PC 从 7 起跳、T 计数器从 T3 进入微码，程序行为完
全变样，而排查者很容易误判为板子接错了。复位链是整机最便宜的确定性来
源：一个按钮、两只电阻、一只电容，再加一片施密特反相器。
\end{warningbox}

\topic{上电检查顺序先于任何程序}机器接完、或每次改动接线之后，都不要
急着把程序拨进 RAM。上电检查有一套固定顺序，每一步都有明确的预期现象，
前一步不符合就不做下一步。顺序按被依赖的程度排：电源不依赖任何东西；
时钟只依赖电源；复位的效果要靠时钟才能观察；总线空闲态则依赖前三者再
加上全部三态使能。

把这套顺序画成流程图，四步的依赖与各自的出口一并标出：

\begin{pdfdiagram}
  \node[hw state, align=center, minimum width=2.6cm] (s1) at (0,1.0) {1 电源\\5 V $\pm$ 0.25 V\\无芯片发热};
  \node[hw compute, align=center, minimum width=2.6cm] (s2) at (3.3,1.0) {2 时钟\\自动约 1 Hz\\单步一下一变};
  \node[hw control, align=center, minimum width=2.6cm] (s3) at (6.6,1.0) {3 复位\\PC = 0、回 T0\\仅 CO+MI 有效};
  \node[hw state, align=center, minimum width=2.6cm] (s4) at (9.9,1.0) {4 总线空闲\\读 0x00\\至多一个驱动者};
  \draw[hw bus] (s1.east) -- node[midway, above=1pt, hw label, fill=white, inner sep=1pt]{通过} (s2.west);
  \draw[hw bus] (s2.east) -- node[midway, above=1pt, hw label, fill=white, inner sep=1pt]{通过} (s3.west);
  \draw[hw bus] (s3.east) -- node[midway, above=1pt, hw label, fill=white, inner sep=1pt]{通过} (s4.west);
  % 不符出口：四步共用一条干线汇入
  \draw[draw=BookMuted, dashed, line width=0.5pt] (s1.south) -- (0,-0.1);
  \draw[draw=BookMuted, dashed, line width=0.5pt] (s2.south) -- (3.3,-0.1);
  \draw[draw=BookMuted, dashed, line width=0.5pt] (s3.south) -- (6.6,-0.1);
  \draw[draw=BookMuted, dashed, line width=0.5pt] (s4.south) -- (9.9,-0.1);
  \draw[draw=BookMuted, dashed, line width=0.5pt] (0,-0.1) -- (9.9,-0.1);
  \fill[BookMuted] (0,-0.1) circle (0.05);
  \fill[BookMuted] (3.3,-0.1) circle (0.05);
  \fill[BookMuted] (6.6,-0.1) circle (0.05);
  \fill[BookMuted] (9.9,-0.1) circle (0.05);
  \node[hw memory, align=center, minimum width=3.6cm] (stop) at (4.95,-1.1) {任一步不符：停在该步修好\\再从该步继续，不跳级};
  \draw[hw return] (4.95,-0.1) -- (stop.north);
  \node[hw label, anchor=east] at (-0.08,-0.55) {预期不符};
\end{pdfdiagram}
\diagramnote{上电检查流程（与正文顺序表一致）。顺序按依赖排：电源不依赖任何东西，时钟只依赖电源，复位的现象要靠时钟观察，总线空闲态依赖前三者加全部三态使能。每步先写下预期再通电：符合才前进；任一步不符就停在该步，修好后从该步继续——换一片芯片、挪一组飞线这类局部改动，至少回到被改动模块所属的那一步重查。}

\begin{longtable}{L{0.13\linewidth}L{0.27\linewidth}L{0.30\linewidth}L{0.20\linewidth}}
\toprule
步骤 & 检查内容 & 预期现象 & 异常先查何处 \\
\midrule
1 电源 & 各片 VCC/GND、去耦电容 & 每片 VCC 脚对地 5V$\pm$0.25V，无芯片发热 & 电源轨断点、极性接反 \\\tablerowrule
2 时钟 & 自动挡约 1Hz，再试单步挡 & 总线 LED 每秒约变一次；按一下只变一次 & 555 外围 R/C、消抖电容 \\\tablerowrule
3 复位 & 按住复位按钮再松开 & PC 显示 0、回到 T0、仅 CO 与 MI 有效 & 复位链路、各片清零脚 \\\tablerowrule
4 总线 & 复位后逐拍单步空走 & 空闲拍读 \code{0x00}，每拍至多一个驱动者 & 三态使能默认电平、控制线接反 \\
\bottomrule
\end{longtable}

第一步查电源与去耦：每片 74HC 的电源脚到位，每片旁边一只 100nF 去耦电
容，每块面包板的电源轨入口再加一只 10$\mu$F。第二步查时钟，该现象前文
已经写过。第三步按住复位按钮：PC 的 4 位 LED 应全灭，T 状态显示应回到
T0，控制线上只剩 CO 与 MI 有效，其余全灭；松开后再按一次，现象应完全
重复。第四步看总线空闲态：复位释放后、程序装入前逐拍单步，凡是没有输
出控制线有效的拍，总线 LED 都应因下拉电阻而读 \code{0x00}；任何两片芯
片微热，都说明某处有两个驱动者在同时说话。

如果只是局部改动——换了一片 244、挪了一组飞线——不必每次都从第一步
重来，但至少要从被改动模块所属的那一步查起：动过电源轨就回第一步，动
过时钟或复位就回第二、三步，动过任何三态使能就回到第四步逐拍空走。四
步全部符合预期，机器才具备装入程序的条件。这个顺序看起来慢，其实是最
快路径：程序运行不对时，故障一定在四步之一或程序本身，先把四步重走一
遍，就把排查范围切掉一半。后面各节搭每个模块时沿用同一思想的展开：先
静态、再单步、后程序。

最后补一条纪律：每一步的预期现象要先写在纸上，再通电。先写预期再看现
象，是防止“看着差不多就放过”的唯一办法——人眼有一种替电路找借口的
天赋，灯大致在闪就当时钟正常，字大致在变就当总线正常。纸上的预期是通
电前写下的，不会被现场说服；实测与它不符就停下来查，这一步省掉的时间，
最后都会在排查里加倍还回去。

\begin{keyidea}
总线居中、两路时钟、确定复位、固定上电顺序，这四件事的共同目的只有一
个：让“机器此刻在干什么”永远是一个可以当场回答的问题。教学机的性能
无关紧要，可观察性就是它的全部价值。
\end{keyidea}

\section{二、寄存器与 ALU 模块}

\topic{寄存器模块 = 触发器 + 三态输出}一个挂在总线上的寄存器要做三件
事：平时保持自己的 8 位值；装载控制线有效的那拍，在时钟沿把总线上的值
锁进来；轮到输出时把自己的值驱动到总线上，其余时刻保持沉默。三件事对
应三种器件资源：一组 D 触发器、一条受控的装载通路、一排三态输出门。
74 系列里现成的组合有两档：74HC273 一片集成 8 个 D 触发器，带公共时钟
和异步清零，但输出直出不带三态；74HC173 一片集成 4 个 D 触发器，自带
装载使能和三态输出，两片拼成一个 8 位寄存器。

\begin{longtable}{L{0.18\linewidth}L{0.36\linewidth}L{0.36\linewidth}}
\toprule
 & 74HC273 & 74HC173 \\
\midrule
集成规模 & 一片 8 位，正好一个寄存器 & 一片 4 位，两片拼 8 位 \\\tablerowrule
装载使能 & 无，需门控时钟 & 有（低有效），时钟直连 \\\tablerowrule
输出形式 & 直出，上总线需外加三态缓冲 & 自带三态（低有效使能），可直挂总线 \\\tablerowrule
异步清零 & 有（$\overline{CLR}$，低有效） & 有（MR，高有效） \\
\bottomrule
\end{longtable}

用 273 的接法：8 个 D 端接总线，$\overline{CLR}$ 接复位线。273 没有装
载使能脚，“AI 有效才装载”用门控时钟实现——系统时钟与 AI 进一片
74HC08 与门，输出接 273 的时钟脚：AI=0 时时钟被挡住，寄存器保持；
AI=1 时的那个上升沿，总线值进入寄存器。门控时钟能安全使用有一个前提：
控制线必须在时钟低电平期间就已经稳定。本机的 T 状态在时钟下降沿推进
（“控制单元：指令译码与微码”一节），AI 在装载上升沿到来前半个拍就已就位，与门输出因此只有一个干
净的装载沿，不会切出多余边沿。输出侧补一片 74HC244 三态缓冲：输入接
273 的 Q，输出接总线；注意 244 的输出使能是低有效，AO 是高有效，中间
要加一级 74HC14 反相。AO 一撤，244 回到高阻，A 寄存器对总线沉默。

这段接法画成图就是“一进一出两条路、一根门控时钟”。

\begin{pdfdiagram}
  % 寄存器本体与三态输出
  \node[hw state, minimum width=2.4cm] (reg) at (3.2,1.4) {74HC273\\8×D 触发器};
  \node[hw compute, minimum width=2.2cm] (buf) at (7.0,1.4) {74HC244\\三态缓冲};
  \draw[hw bus] (0.4,1.4) -- node[midway, above=1pt, hw label, fill=white, inner sep=1pt]{总线 D[7:0]} (reg.west);
  \draw[hw bus] (reg.east) -- node[midway, above=1pt, hw label, fill=white, inner sep=1pt]{Q[7:0]} (buf.west);
  \draw[hw bus] (buf.east) -- node[midway, above=1pt, hw label, fill=white, inner sep=1pt]{回总线（三态）} (9.8,1.4);
  % 门控时钟：听 = 时钟沿
  \node[hw compute, minimum width=1.6cm] (gate) at (3.2,-0.9) {与门\\74HC08};
  \draw[hw return] (0.6,-0.75) -- node[midway, above=1pt, hw label, fill=white, inner sep=1pt]{系统时钟} ($(gate.west)+(0,0.15)$);
  \draw[hw return] (0.6,-1.05) -- node[midway, above=1pt, hw label, fill=white, inner sep=1pt]{AI} ($(gate.west)+(0,-0.15)$);
  \draw[hw return] (gate.north) -- node[midway, right=1pt, hw label, fill=white, inner sep=1pt]{门控时钟} (reg.south);
  % 复位与输出使能：说 = 组合电平
  \draw[hw return] (3.2,2.95) -- node[midway, right=1pt, hw label, fill=white, inner sep=1pt]{复位（/CLR）} (reg.north);
  \draw[hw return] (7.0,2.95) -- node[midway, right=1pt, hw label, fill=white, inner sep=1pt]{AO 经反相接 /OE} (buf.north);
  \node[hw label] at (5.9,-1.9) {听 = 时钟沿（沿触发）；说 = 组合（电平使能）};
\end{pdfdiagram}
\diagramnote{寄存器模块（273 方案）。左侧总线值进 D 口，右侧 Q 口经 74HC244 三态缓冲回总线。装载是沿触发的：系统时钟与 AI 相与产生门控时钟，只有 AI=1 的拍才出现装载沿；输出是组合的：AO 经反相打开 244 的使能，整拍驱动，AO 一撤立即高阻。173 方案把装载使能和三态都收进片内，时钟直连，但信号语义与此图完全相同。}

用 173 的接法可以省掉外挂缓冲：两片分管低 4 位与高 4 位，D 端接总线，
Q 端直接挂总线。173 的装载使能和输出使能同样都是低有效，各加一级反相
再接 AI、AO；时钟则直接接系统时钟，不再需要门控时钟——这是 173 比
273 干净的地方，代价是芯片从一片变两片再加反相器。两种接法在教学机里
都常见，本章后续一律只说“AI 有效沿装载、AO 有效驱动总线”，具体器件
任选一种即可。无论选哪种，每个寄存器的 Q 端都配一排 LED：寄存器存的
是什么，应当和总线值一样随时可读。两种方案的核心接线可以各浓缩成三行，
插线时对照：

\begin{lstlisting}
-- 273 方案（以 A 寄存器为例）
D[7:0] <- 总线；CLK <- 系统时钟 AND AI；/CLR <- 复位
Q[7:0] -> 74HC244 -> 总线；244 的 /OE <- NOT AO
-- 173 方案（两片拼 8 位）
D[7:0] <- 总线；装载使能 <- NOT AI；CLK <- 系统时钟
Q[7:0] -> 总线；输出使能 <- NOT AO
\end{lstlisting}

\begin{warningbox}
“状态、时钟与时序逻辑”一章把“用普通与门切时钟”列为禁令，针对的是有同步时序预算的设计：
时钟一快，门控引入的偏斜和使能毛刺就会被当成真实时钟沿，建立与保持
无从保证。这台教学机里门控时钟却是标准做法，差别不在原则而在条件：
只有低速、控制字提前半拍稳定时门控时钟才安全——本机时钟慢到毫秒量
级，T 状态在下降沿推进（“控制单元：指令译码与微码”一节），控制线在装载上升沿前半拍就已就位，
与门输出不可能切出多余边沿，“状态、时钟与时序逻辑”一章担心的那条失效路径在这里不成立。
条件变了结论就变——若哪天把这台机器移植到同步时序严格的器件上，第
一件事就是把门控时钟改回时钟使能。
\end{warningbox}

A、B 两个寄存器比一般模块多一捆线：除总线接口外，A 的 Q 还有 8 根专用
线直连 ALU 的 A 输入，B 的 Q 同样直连 ALU 的 B 输入，不走总线、不经三
态、常年有效。原因在 ALU 是组合逻辑：它要求两个操作数在同一拍里稳定
在场，而总线一拍只能传一个值——若操作数也走总线，一拍之内先送 A 再送
B，结果还没地方回来。专用线的代价是 16 根飞线，换来加法那拍的极简格
局：A、B 站好，ALU 出结果，$\Sigma$O 把结果送上总线，AI 把它收回 A。
另外，B 寄存器根本不给总线配输出：没有任何指令需要把 B 单独送上总线，
它的输出三态可以直接省掉，273 方案里因此少接一片 244。B 的值从哪来？
仍然从总线来：ADD、SUB 的第一个执行拍用 RO+BI 把 RAM 里的操作数抄进
B，专用线再把这份稳定的 B 呈给 ALU。装载走总线、输出走专线，两个方向
各用各的路，这正是 A、B 与一般模块的差别所在。

A 在全机寄存器里活得最累：LDA、LDI 从总线装它，ADD、SUB 的结果写回它，
STA 要它说、OUT 要它说、JC/JZ 的判断间接来自它。它既是累加器又是事实
上的通用暂存，驱动总线的次数全机最多。所以 A 模块的接线要格外扎实：
AI、AO 两根控制线、总线接口、ALU 专用线三路交汇，任何一处虚接都会在
最频繁的通路上发作。查整机怪故障时，A 模块永远值得第一个怀疑。

寄存器模块的插线顺序也固定：先插芯片、接电源与去耦，量过 VCC；再接 D
口到总线的 8 根线；然后接 Q 端到 LED 排；最后才接控制线——装载、输出
使能、清零。控制线最后接的理由是：没接控制线时模块应当完全安静，不装
载也不驱动，此时总线其他部分的调试不受它干扰。每接好一组就停下来核对
一次，比一次全接完再回头找错快得多。

寄存器模块不进程序就能检查。静态：用跳线在一排空排上摆出已知字节（如
\code{0x5A}），手动给一个装载沿，Q 端 LED 应显示同一字节。单步：撤掉
装载使能再按几拍，值应保持不变；让输出使能有效，总线 LED 应重现该字
节。存、装、说三个动作分别正确，这片寄存器才算搭好，才可以往总线上挂。

\topic{IR 和 OUT 是两类特殊寄存器}IR 也是 8 位寄存器，装载与 A 完全一
样：II 有效沿把总线上的指令字节锁进来。特殊在输出分两半、各走各路。
高 4 位是 opcode：不经总线，专用线直连微码 ROM 的地址高位（“控制单元：指令译码与微码”一节），常
年有效——译码逻辑要随时看得到“当前是哪条指令”。低 4 位是地址或立即
数字段：经三态上总线，由 IO 控制，服务于两条路径：把地址字段送 MAR
（IO+MI，见于 LDA、ADD、STA 等带地址的指令），或把立即数送 A
（IO+AI，即 LDI 的装载拍）。用接线列表写下来就是：

\begin{lstlisting}
IR.D[7:0] <- 总线            -- II 有效沿锁存
IR.Q[7:4] -> 微码 ROM 地址   -- 直连，不经三态
IR.Q[3:0] -> 三态 -> 总线    -- 仅 IO 有效时驱动低 4 位
\end{lstlisting}

接法上的要点是：低 4 位的三态缓冲只接 Q0--Q3，高 4 位绝不允许跟着上总
线。用两片 173 拼 IR 时，低 4 位那片的输出使能接 IO，高 4 位那片的输
出使能直接固定在无效电平——从接线上删掉高 4 位驱动总线的能力，比指
望微码永远不犯错可靠。CO 与 IO 都只驱动总线低 4 位，高 4 位悬空：MAR
本就只锁存低 4 位，不受影响；LDI 拍 AI 锁存整个字节时，高 4 位正是
“总体架构：总线、时钟与复位”一节那排下拉电阻读出的 0，于是 A 收到的是干净的 \code{0x0} 拼接立即数。
一个前面看似多余的布局决定，到这里显出用途。

高 4 位走专用线而不是总线，也是同一套时刻关系决定的。译码逻辑在每个
T 状态都要看到 opcode：T2 起的每一拍，微码 ROM 的地址都是 opcode 与 T
的拼接（“控制单元：指令译码与微码”一节展开），opcode 必须常年在场。若让 opcode 也经总线送译码，
就得专门安排一拍把 IR 高 4 位驱动出去，而那一拍总线往往正被取指或运算
占用。寄存器输出分两路、各走各的，译码与数据传送就互不排队。

OUT 是另一种特殊：只进不出。OI 有效沿从总线锁存——OUT 指令的装载拍
是 AO+OI，A 说、OUT 听；Q 端直驱显示，8 只 LED 一排，或经译码驱动数码
管（接法“内存与输出模块”一节展开），输出侧没有任何通往总线的三态门。输出设备不该驱动
总线，理由有三层。功能上，没有任何指令需要把 OUT 的值读回来，配了输出
门也是浪费。电气上，总线上每多一个驱动者，就多一根可能插反的使能线、
多一处争用隐患。原则上，“只进不出”应当由接线来保证：即使微码写错、
杜邦线插错排，OUT 在物理上也不可能与别人抢总线。把约定做进硬件而不是
写进纪律，是这台小机器反复使用的思路——“总体架构：总线、时钟与复位”一节的三态使能默认无效、本段
IR 高 4 位不上总线、OUT 不配输出门，是同一条思路的三次应用。

OUT 还是整台机器的脸面。总线 LED 显示的是“此刻正在流动”的值，随拍就
变；OUT 显示的是“程序想让人看到”的值，装载一次就一直亮着。机器执行
HLT 之后，时钟停了、总线静了，最后留在面板上的就是 OUT 的那一排灯——
程序的对错，最终由它宣布。所以 OUT 的显示排值得放在面板最显眼的位置，
朝向观看者，而不是藏在模块背后。

\topic{ALU = 74HC283 加法器加一路取反}本机 ALU 只做两种运算：A+B 与
A$-$B。两种运算共用一条加法路径，差别全在 B 进入加法器之前怎样处理。
依据是补码恒等式 $A-B = A+\bar{B}+1$，其中 $\bar{B}$ 表示 B 的逐位取
反。于是 ALU 的全部家当是：两片 74HC283 拼成 8 位加法器，两片 74HC86
拼成 8 位可控取反器，再加几门标志逻辑。这正是算术与数据通路相关章节“一条加法路径加输
入处理”层次链的末端：门组成全加器，全加器连成 283，283 配上可控取反
就是加减法 ALU。

两片 283 级联：低位片的进位输出接高位片的进位输入，高位片的进位输出
就是全 8 位的 CF 来源。A 输入 8 位直连 A 寄存器的 Q，即上一段的专用线。
B 输入不直连，先过 86：每片 86 含 4 个异或门，每个门一端接 B 的对应位，
另一端 8 个门并在一起接 SU 控制线。异或的性质 $X\oplus0=X$、
$X\oplus1=\bar{X}$ 正好构成可控取反：SU=0 时 B 直通，SU=1 时 B 逐位取
反。同一根 SU 还接到最低位加法器的进位输入：加法时 CI=0，结果是 A+B；
减法时 CI=1，结果是 $A+\bar{B}+1$，即 A$-$B。一根控制线同时完成取反与
加一，这是补码减法最省器件的形态。ALU 的结果经一片 74HC244 三态缓冲挂
总线，使能经反相接 $\Sigma$O。

这套组合画成结构图，就是一条加法路径加一路可控取反。

\begin{pdfdiagram}
  % 操作数：A、B 专线常年在场
  \node[hw state, minimum width=1.9cm] (areg) at (-0.4,1.5) {A 寄存器};
  \node[hw state, minimum width=1.9cm] (breg) at (-0.4,-1.0) {B 寄存器};
  \node[hw compute, minimum width=2.2cm] (xor) at (2.9,-1.0) {74HC86 ×2\\可控取反};
  \node[hw compute, minimum width=3.0cm] (add) at (6.3,-0.3) {74HC283 ×2\\8 位加法器};
  \node[hw compute, minimum width=1.7cm] (buf) at (9.1,-0.3) {74HC244};
  \draw[hw bus] (areg.east) -- (5.7,1.5) -- (5.7,0.24);
  \node[hw label, above=1pt, fill=white, inner sep=1pt] at (3.3,1.5) {A[7:0] 专线};
  \draw[hw bus] (breg.east) -- node[midway, above=1pt, hw label, fill=white, inner sep=1pt]{B[7:0]} (xor.west);
  \draw[hw bus] (xor.east) -- (4.4,-1.0) -- (4.4,-0.45) -- (4.8,-0.45);
  % SU：一路控取反，一路做最低位 CI
  \draw[hw return] (2.9,-2.5) -- (xor.south);
  \node[hw label, left=1pt, fill=white, inner sep=1pt] at (2.9,-2.0) {SU};
  \draw[hw return] (2.9,-2.2) -- (6.1,-2.2) -- (6.1,-0.84);
  \fill[BookTeal] (2.9,-2.2) circle (0.05);
  \node[hw label, above=1pt, fill=white, inner sep=1pt] at (4.5,-2.2) {SU 兼作最低位 CI（减法 +1）};
  % 结果上总线
  \draw[hw bus] (add.east) -- node[midway, above=1pt, hw label, fill=white, inner sep=1pt]{$\Sigma$[7:0]} (buf.west);
  \draw[hw bus] (buf.east) -- node[midway, above=1pt, hw label, fill=white, inner sep=1pt]{上总线} (10.5,-0.3);
  \draw[hw return] (9.55,-1.7) -- ($(buf.south)+(0.45,0)$);
  \node[hw label, anchor=west, fill=white, inner sep=1pt] at (9.65,-1.2) {$\Sigma$O};
  % 标志：CF 取自高位片 C4，ZF 来自 8 输入或非
  \draw[hw bus] (7.0,0.24) -- (7.0,1.5);
  \node[hw label, anchor=west, fill=white, inner sep=1pt] at (7.1,1.5) {CF（高位 C4）};
  \node[hw compute, minimum width=2.4cm] (nor) at (7.5,-2.1) {74HC4078\\8 输入或非};
  \draw[hw bus] (8.0,-0.3) -- (8.0,-1.58);
  \fill[BookMuted] (8.0,-0.3) circle (0.05);
  \draw[hw bus] (nor.east) -- node[midway, above=1pt, hw label, fill=white, inner sep=1pt]{ZF → Flags} (10.0,-2.1);
\end{pdfdiagram}
\diagramnote{ALU 结构。A、B 经专线常年到场：A 直连两片级联 74HC283 的 A 输入，B 先过两片 74HC86 可控取反——SU=0 直通、SU=1 逐位取反，同一根 SU 分叉接最低位进位输入 CI，取反加一一步完成，两片 283 之间的片间进位（低位 C4 接高位 CI）在框内完成。结果经 74HC244 三态上总线，使能由 $\Sigma$O 反相驱动；CF 取自高位片进位输出，ZF 由 74HC4078 对 8 个和位取或非得到。两个标志此刻还是组合信号，要等 FI 那拍才锁存。}

接线之前先手工演算两组数，确认这套组合真的等于减法。算 7$-$5：
\code{00000111+11111010+1}，和为 \code{100000010}——8 位结果是
\code{00000010} 即 2，进位输出为 1。再算 5$-$7：
\code{00000101+11111000+1}，和为 \code{011111110}，8 位结果
\code{11111110} 正是 $-$2 的补码，进位输出为 0。两组数合在一起说明：
减法时进位输出为 1 表示没有借位，即 A$\geq$B；为 0 表示 A$<$B、结果是
补码负数。所以 JC 接在 SUB 后面时，语义是“A 大于等于 B 则跳”，这个约
定“编程、调试与扩展”一节写程序时会反复用到。加法的进位输出就是普通溢出，例如 200+100：
\code{11001000+01100100}，结果 \code{00101100} 即 44，CF=1。

两个标志都从加法器输出组合产生。CF 最直接：高位 283 的进位输出引出即
是。ZF 要求“结果全 0”：8 个和位送进一片 74HC4078（8 输入或/或非），
取它的或非输出——任何一位为 1 输出即 0，全为 0 时输出 1，这正是 ZF。
要注意此刻两个标志都还是组合信号：它们只对当前这一拍的 A、B、SU 有效，
下一拍 A 或 B 一变，标志跟着变。组合标志没法直接给跳转用，因为 JC/JZ
的判断拍在加法拍之后，中间还隔着若干拍——标志必须先被记住。

顺带有两个接线细节。其一，244 之前、加法器输出之后是一个值得利用的观
察点：在这里分一排 LED，就能看到“当下 A、B 与 SU 的函数值”——A、B
的专用线常年有效，这组组合结果也就常年有效，即使 ΣO 没让结果上总线，
低头也能读出它。调试 ALU 时看的正是这组灯。其二，SU 一根线要分叉到 9
个终点：8 个异或门加最低位进位输入。74HC 输出的扇出能力带这 9 个负载
毫无压力，但走线要短而整齐，因为 SU 一切换，整条 ALU 的函数就变了。

这条 ALU 不进整机就能先做静态检查，它本质上就是算术与数据通路相关章节那个 4 位加减法
实验台加宽到 8 位：A、B 各用 8 只拨码开关代替寄存器直接摆值，SU 用一
只开关手控，结果与 CF、ZF 各接 LED。加法摆几组、减法摆几组，其中包括
7$-$5、5$-$7、200+100 这三组已经手工算过的数，LED 读数与手算一一对上，
ALU 这片电路才算可信，A、B 寄存器的专用线这时再接过来。若等到装进整机
才发现减法不对，就分不清是 ALU 错、专用线错还是微码错。

进位链的速度也不用操心：两片 283 级联，进位从最低位传到最高位经过两级
片内进位加一次片间传递，总延迟在几十纳秒量级；本机一拍是几十毫秒，进
位有充足时间走完，不需要先行进位之类的加速结构。慢，在这里不是缺陷而
是设计目标的一部分。

\begin{classicnote}
“减法就是加上负数”是经典教材反复强调的一点：有了补码，数据通路里只
需要一条加法路径，减法、比较、条件跳全都复用它。本机的 ALU 把这一点做
到了极简——一片加法器、一路可控取反、一根兼作加一的 SU 线，再没有别
的运算硬件。
\end{classicnote}

\topic{标志寄存器在加法完成那拍装载}记住标志的办法是在它还新鲜的那拍
锁存。Flags 是一片小寄存器：两个 D 触发器（一片 74HC74 正好，或从
74HC273 借两位），D 端分别接组合的 CF 与 ZF，装载时钟由 FI 控制，与
A、B 相同的门控时钟结构。微码只在 ADD、SUB 的求和拍让 FI 有效——与
$\Sigma$O、AI 同一拍，这一拍的上升沿同时做两件事：A 锁存和或差，
Flags 锁存这一对 CF、ZF。此后 FI 一直保持无效，锁存值原样保持，直到下
一次加法或减法把它们刷新。

清零接法随器件而定：从 273 借两位时，那片 273 的 $\overline{CLR}$ 直接
接复位线；用 74HC74 时，两个触发器各自的异步清零脚并起来接复位线。两
种接法的效果相同——复位期间 CF、ZF 都为 0，与“总体架构：总线、时钟与复位”一节“标志必须接复位”
的要求一一对应。Flags 的 Q 端同样各接一只 LED：锁存的 CF、ZF 是条件跳
转的依据，应当像寄存器值一样常年可见。

跳转指令读的是锁存值：JZ 的判断拍看锁存的 ZF 是否为 1，JC 看锁存的
CF。因为读的是寄存器输出而不是组合输出，判断拍总线上走什么、ALU 输入
变成什么，都与这次判断无关——这正是要锁存的原因。由此也看得清微码的
一条硬约定：全机只有 ADD、SUB 的求和拍有 FI，LDA、LDI、OUT 等一概不刷
新标志。程序里“先 ADD、隔几条指令再 JZ”是安全的，“先 LDA 再 JZ”判
断的是更早一次运算留下的标志。复位时 Flags 清零（“总体架构：总线、时钟与复位”一节），保证第一次
ADD/SUB 之前 JC/JZ 读到的是确定值。

“运算顺便刷标志”还塑造了本机的程序风格。循环计数不必单独比较：把循
环变量放在 RAM 里，每轮用 SUB 减一，减到零的那次 ZF 自然置起，紧跟一
条 JZ 就跳出循环——一条 SUB 同时完成算数和判断，这正是把标志锁存在求
和拍带来的便利。“编程、调试与扩展”一节的完整程序会用到这个惯用法。

FI 这一拍也可以单步核对。让机器停在 ADD 的求和拍：$\Sigma$O 应有效、
总线应显示和、AI 与 FI 应同时有效；按一下单步，A 的 LED 更新为和，标
志的 LED 同步更新。再按几拍让总线和 ALU 输入随意变化，标志 LED 应纹丝
不动——动的只有下一次 ADD 或 SUB 的求和拍。这两段观察合起来，锁存的
意义就不是定义而是亲眼所见。

模块按依赖关系从简单到复杂检查：先 ALU、后标志——标志锁存的只是 ALU
的组合输出，ALU 不对，标志锁得再准也是错的。每个模块内部则统一按“先
静态、再单步、后程序”三步走：静态检查不上时钟，用跳线摆输入、看组合
输出，它不过，后面两步都没有意义；单步检查用按钮时钟逐拍推进，核对每
一拍的装载与输出，它不过，程序层面的任何异常都无法定位；程序检查才让
模块在整机里跑指令。各模块的检查点汇总如下：

\begin{longtable}{L{0.13\linewidth}L{0.29\linewidth}L{0.27\linewidth}L{0.21\linewidth}}
\toprule
模块 & 静态检查 & 单步检查 & 程序检查 \\
\midrule
A/B 寄存器 & 总线摆 \code{0x5A}，手动装载沿后 Q 显示 \code{0x5A} & 使能有效拍装载、无效拍保持；AO 拍总线重现该值 & LDI/LDA 后 A 的 LED 为预期值 \\\tablerowrule
IR & 装载后 Q 直通显示，高低 4 位分路正确 & II 拍锁存；IO 拍仅低 4 位上总线 & 取指派后 IR 与 RAM[0] 一致 \\\tablerowrule
OUT & OI 装载后显示保持不闪 & AO+OI 拍把 A 的值抄到显示 & OUT 指令后显示当前 A \\\tablerowrule
ALU & A/B 摆定值，SU 切换看和、差与 CF/ZF & $\Sigma$O 拍结果上总线，其余拍高阻 & ADD/SUB 后 A 与标志同步更新 \\\tablerowrule
Flags & 翻动 ALU 输入，锁存值不随动 & FI 拍锁存新值，此后保持不变 & ADD 后接 JC/JZ，走向与手工预算一致 \\
\bottomrule
\end{longtable}

至此数据通路的一半已经立在板上：寄存器能存能说，ALU 能算，标志能记。
“内存与输出模块”一节给它配上记忆——MAR 与 RAM——和脸面，“控制单元：指令译码与微码”一节再给它配上发号施令的控
制单元。

\section{三、内存与输出模块}

本节用到的整机信号先集中登记，读到后文任何一根都可以回来查：

\begin{longtable}{L{0.14\linewidth}L{0.16\linewidth}L{0.52\linewidth}}
\toprule
信号 & 方向 & 在本节的作用 \\
\midrule
MI & 入 & MAR 的装载使能，拍末把总线低 4 位锁存为地址 \\\tablerowrule
RI & 入 & RAM 写请求，经门控成为 189 的写窗口 \\\tablerowrule
RO & 入 & RAM 读请求，经反相打开 240 的三态输出 \\\tablerowrule
OI & 入 & OUT 寄存器的装载使能，拍末锁存总线值 \\\tablerowrule
总线 D0–D7 & 双向 & 地址、数据、显示值都经它中转 \\\tablerowrule
时钟、复位 & 入 & 节拍与清零，“总体架构：总线、时钟与复位”一节的约定原样沿用 \\
\bottomrule
\end{longtable}

\topic{RAM 模块的两种实现路线}
16 字节内存用两片 74HC189 拼成。74HC189 是一片 16 字乘 4 位的静态 RAM：4 根地址线选中 16 个字之一，4 路数据输入、4 路数据输出，控制线只有两根——片选 S（低有效）和写使能 W（低有效）。读的条件是 S 有效且 W 无效，输出脚驱动读出值；写的条件是 S 有效且 W 有效，输入脚的值进入选中字。两片并排放置，一片接总线低 4 位、一片接总线高 4 位，地址线并联，就得到 16 乘 8 的内存。位宽扩展在这里没有任何译码技巧，纯粹是每片管一半位，与存储与设备事务相关章节讲的存储器位宽扩展是同一种做法，只是规模小到能放在手掌上。

接线之前，先把 189 的引脚与整机信号一一对应列成表。照表核对比对着原理图找线可靠得多，尤其是两片并联的模块，任何一根跨片的错接都会表现为“高 4 位或低 4 位整体出错”。

\begin{longtable}{L{0.14\linewidth}L{0.20\linewidth}L{0.12\linewidth}L{0.42\linewidth}}
\toprule
引脚 & 名称 & 方向 & 在整机中的接向 \\
\midrule
A0–A3 & 地址输入 & 入 & MAR 的 Q0–Q3，两片并联 \\\tablerowrule
D1–D4 & 数据输入 & 入 & 数据总线对应的 4 位 \\\tablerowrule
S & 片选，低有效 & 入 & 接地，恒选中 \\\tablerowrule
W & 写使能，低有效 & 入 & RI 与时钟低半拍相与后取反，两片并联 \\\tablerowrule
O1–O4 & 数据输出，反相三态 & 出 & 74HC240 的输入 \\\tablerowrule
Vcc、GND & 电源 & — & 5 V，每片电源脚旁并 0.1$\mu$F 去耦电容 \\
\bottomrule
\end{longtable}

两片的分工按总线位次划：一片的 D、O 接总线 D0–D3，另一片接 D4–D7，两片之间没有任何信号交叉。先焊低 4 位那片、单独测通，再焊高 4 位那片，是最省心的顺序。去耦电容要贴着每片的电源脚放：RAM 在读写切换瞬间的电流突变全靠它就地吸收，省掉这一步，后面会出现“时而对时而错”的软故障，是所有故障里最难定位的一类。

时序余量也值得算一次。5 V 下 74HC189 的读取延迟约三四十纳秒，74HC240 的打开延迟十几纳秒，合起来从 RO 有效到总线数据稳定不超过一百纳秒；而本机时钟最激进的挡位也不过 48Hz，一拍是毫秒级——最快挡约 20 毫秒。一百纳秒对 20 毫秒，余量约五个数量级，所以这台机器不需要任何等待状态——读周期里“地址稳定、打开 RO、等访问时间、拍末采样”这四个动作，前三个在一拍的开头瞬间就完成了，存储与设备事务相关章节讨论的那种等存储器的无奈，在这里一次都不会出现。

顺手把读数据手册的方法在这片芯片上练一遍。存储器手册里真正要读的参数只有三类：读路径的延迟、写路径的窗口、各信号的先后约束。把 74HC189 的典型量级列在这里（确切数值以手边器件的手册为准），并对照本机给出结论：

\begin{longtable}{L{0.26\linewidth}L{0.20\linewidth}L{0.36\linewidth}}
\toprule
手册参数 & 典型量级 & 对本机的含义 \\
\midrule
地址到输出的读取延迟 & 几十纳秒 & 毫秒级一拍（最快挡约 20 毫秒），余量约五个数量级 \\\tablerowrule
写脉冲最小宽度 & 二十纳秒上下 & 低半拍开窗，窗口宽十毫秒上下 \\\tablerowrule
数据先于写沿的建立时间 & 二十纳秒上下 & 数据在窗口打开前已稳定半拍 \\\tablerowrule
写沿后数据的保持时间 & 接近零 & 窗口先关、数据后撤，天然满足 \\
\bottomrule
\end{longtable}

看手册的顺序因此是：先找这几行参数，再回看自己的时钟周期，只要周期比参数大三四个数量级，时序这一章就可以合上书了。这台机器的全部存储器时序分析到此为止——剩下的都是逻辑问题。

用 74HC189 有一个必须处理的细节：它的数据输出是反相的——读出来的是存储值按位取反的结果（同系列里同相输出的型号是 74HC219，不如 189 常见）。还原的方法不是加一片普通反相器，而是加一片 74HC240。240 是八反相三态缓冲器，两个四位组各有独立的输出使能，把它们并接后由 RO 控制（240 的使能低有效，RO 经一级反相后接入）。这一片器件同时完成两件事：把反相再反相、还原出原值，以及在 RO 无效时让输出呈高阻。“RAM 输出经三态上总线”这条整机约定就藏在这一片 240 里：RO 有效，RAM 选中的字出现在总线上；RO 无效，240 的输出端从总线上断开，RAM 对总线彻底隐身，别的驱动者感觉不到它的存在。

240 的八路缓冲在本模块里刚好用满：两个四位组的使能端（1G、2G，低有效）并接后由 RO 反相驱动，八路输入接两片 189 的反相输出，八路输出上总线。把这片缓冲器的信号也列成表，接线时与上一张引脚表对照使用：

\begin{longtable}{L{0.18\linewidth}L{0.30\linewidth}L{0.34\linewidth}}
\toprule
240 引脚 & 接向 & 说明 \\
\midrule
1G、2G & RO 经反相后并接 & RO 有效时八路同时打开 \\\tablerowrule
输入 8 路 & 两片 189 的 O1–O4 & 收到的是存储值的反相 \\\tablerowrule
输出 8 路 & 数据总线 D0–D7 & 再反相一次，原值上总线 \\
\bottomrule
\end{longtable}

两级反相串联、极性归零，这个结构值得多看一眼：189 的反相输出本是器件的输出特性，240 既校正了输出极性，又提供了三态边界，一个器件同时完成两项功能。若采购到同相输出的 74HC219，结构不变，只把 240 换成同相的 74HC244 即可，其余接线一根不动。

两片 189 的写使能 W 并接，来源是控制线 RI；读路径的开关是 RO。写使能与输出使能的分工因此非常清楚：RI 管“把总线上的值写进 MAR 选中的字”，RO 管“把 MAR 选中的字送到总线上”，两者经过不同的使能路径作用于同一组地址，互不相干。另有一个时序细节值得照做：让写窗口只占时钟的低半拍，即 W 在“RI 有效且 CLK 为低”的区间内打开。这样写窗口在拍末上升沿到来时已经关闭，而 MAR 的地址要在边沿之后、经过触发器延迟才会变化，写窗口关闭永远先于地址变化，不会把数据顺手写进下一个地址。这个细节的代价不过是一级门，换来的是写操作对边沿时刻完全不敏感。

\begin{longtable}{L{0.16\linewidth}L{0.22\linewidth}L{0.24\linewidth}L{0.26\linewidth}}
\toprule
信号 & 接向 & 有效条件 & 作用 \\
\midrule
地址 A3–A0 & MAR 的 Q3–Q0 & 常接 & 选择 16 个字之一 \\\tablerowrule
数据输入 D & 数据总线 & 常接 & 写窗口内由总线驱动者供数 \\\tablerowrule
片选 S & 接地 & 恒有效 & 读写区分交给 W 与 240 \\\tablerowrule
写使能 W & RI 与时钟低半拍相与后取反 & RI 有效的低半拍 & 总线值写入选中字 \\\tablerowrule
反相输出 O & 74HC240 输入 & 读出期间驱动 & 值为存储内容的反相 \\\tablerowrule
240 输出使能 & RO 反相后接入 & RO 有效 & 原值上总线；无效则高阻 \\
\bottomrule
\end{longtable}

第二条路线解决的是另一个问题：程序从哪里来。189 只是存储体，把程序装进去还需要一套输入手段。成本最低的做法是 DIP 开关加一只按钮：4 只开关拨地址，8 只开关拨数据，按钮发出写脉冲。这条路慢，装一个字节要十几秒，但每一个比特都经过手指，地址、数据、写脉冲三个概念被迫分开操作——这恰恰是教学想要的慢。编程电路与运行电路在 189 的三组引脚上交接：地址来源用一片 74HC157（四路二选一选择器）在 MAR 与地址开关之间切换；数据开关经一片 74HC244 送上总线，只在编程模式下使能，189 的数据输入照旧取自总线，无需切换；写脉冲由按钮经 74HC14 施密特整形后直接驱动 189 的 W，绕开 RI。切换的完整次序在本节最后一个主题给出，那里的核心约束只有一句：先停机，再切换。

\begin{longtable}{L{0.24\linewidth}L{0.26\linewidth}L{0.36\linewidth}}
\toprule
路线 & 器件 & 关键问题 \\
\midrule
RAM 本体 & 74HC189 两片 + 74HC240 一片 & 输出反相要靠 240 还原；RI 与 RO 永不同拍 \\\tablerowrule
编程输入 & DIP 开关 12 只、按钮、74HC157、74HC244、74HC14 & 成本低、每字节十几秒；切换前必须先停机，避免与运行电路争用 \\
\bottomrule
\end{longtable}

也有人会问能不能干脆用一片大 SRAM，比如 6116（2K 乘 8）。能，而且读写逻辑与 189 几乎一样，还免去了反相输出的麻烦；代价是 11 根地址线里有 7 根要固定接地或接高，走线反而比两片 189 更琐碎，并且“16 字节的小内存”这个教学尺度被一片 2K 的芯片遮住了——数组大一千倍，能看见的仍然只有 16 个字。189 的价值恰在它的合身：地址脚、字数的规格与整机一一对应，没有一根要假装不存在的线。三条路线里没有对错，只有“想让学习者看见什么”的差别。

把第二条路线再推到极端，还存在一种连 189 都省掉的历史方案：纯开关存储器。七十年代的 Altair 8800 前面板就是一排开关加一排灯，靠手拨开关输入引导程序；再往前，早期计算机的操作员直接在控制台上拨地址和数据。那条路对 16 字节的小机器不是不能用——16 组各 8 只开关加 16 片八输入与门译码，就是一块“硬接线 ROM”——但 128 只开关的体积和 16 片译码器的接线量，换来的只是省掉两片 RAM，教学上是壮举、工程上是倒退。本书把它列为见识而非建议：知道存储器可以原始到就是一排开关，才能理解 189 那 16 个字是多么划算的文明。

DIP 开关路线的真正价值因此不在省钱，而在节奏。用编程器烧 EEPROM 是“整批交付”，错了一个字节也要重新上编程器烧写（28C 系列按字节擦写，重烧该字节即可）；DIP 开关是“逐字节对话”，每按一次按钮都能立刻读回核对。对第一次把程序送进机器的人来说，这种一个字节一次心跳的节奏，本身就是存储器语义的最好教具。

本节涉及的器件汇总如下，备料时照表清点即可；门器件一栏按需备，控制线的反相与门控都从这里出：

\tablechunkintro{1}{2}{74HC189}{74HC14}
\begin{longtable}{L{0.20\linewidth}L{0.10\linewidth}L{0.52\linewidth}}
\toprule
器件 & 数量 & 用途 \\
\midrule
74HC189 & 2 & 16 乘 8 RAM 本体 \\\tablerowrule
74HC240 & 1 & 反相还原加三态上总线 \\\tablerowrule
74HC173 & 1 & MAR \\\tablerowrule
74HC377 & 1 & OUT 寄存器 \\\tablerowrule
74HC157 & 1 & 编程与运行的地址来源切换 \\\tablerowrule
74HC244 & 1 & 编程数据开关上总线 \\\tablerowrule
74HC14 & 1 & 按钮整形，兼作控制线反相 \\
\bottomrule
\end{longtable}

\tablechunkintro{2}{2}{DIP 开关}{0.1$\mu$F 电容}
\begin{longtable}{L{0.20\linewidth}L{0.10\linewidth}L{0.52\linewidth}}
\toprule
器件 & 数量 & 用途 \\
\midrule
DIP 开关 & 12 位 & 编程地址 4 位加数据 8 位 \\\tablerowrule
按钮 & 1 & 编程写脉冲 \\\tablerowrule
LED 加 330 欧电阻 & 8 套 & 二进制显示 \\\tablerowrule
10k$\Omega$ 电阻 & 12 只 & 开关的上拉或下拉 \\\tablerowrule
0.1$\mu$F 电容 & 10 只 & 各片电源去耦 \\
\bottomrule
\end{longtable}

\topic{MAR 是 4 位地址锁存器}
MAR（存储器地址寄存器）只有 4 位，一片 74HC173 正好装下。它必须是锁存器而不是一段导线，原因在于地址要比总线活得久。取指时 PC 只在 T0 这一拍驱动总线，T1 总线就要还给 RAM 传送指令；如果 189 的地址直接取自总线，T1 总线内容一换，地址就丢了，指令会从一个错误的位置读出来。MAR 在 T0 拍末把地址锁存下来，此后无论总线怎么翻，189 的地址引脚始终稳定，直到下一次 MI 有效的拍末才更新。接法上，173 的两个装载使能 E1、E2 并接，经反相后由 MI 驱动；两个输出使能 OE1、OE2 直接接地——MAR 到 189 是一条点对点的专线，不经过共享总线，输出可以常开，不需要三态；清除端接整机复位，保证上电后地址是已知值，而不是一片随机花样。注意 173 的清除端 MR 是高有效，而整机复位链按低有效设计（“总体架构：总线、时钟与复位”一节里 273 的 $\overline{CLR}$、161 的 $\overline{CLR}$ 都是低有效），中间要加一级反相再接复位线，极性接反，这片 MAR 就会常年被钳在清零状态。

整个内存模块的连接关系先画成一张图：MAR 锁地址，157 选地址来源，两
片 189 拼字宽，240 管读出边界。

\begin{pdfdiagram}
  % 左侧总线
  \draw[draw=BookMuted, line width=1.8pt] (0,-1.6) -- (0,2.9);
  \node[hw label, anchor=south] at (0,2.95) {8-bit 总线};
  % MAR 与地址来源选择
  \node[hw state, minimum width=2.2cm] (mar) at (2.4,1.2) {MAR\\74HC173（4 位）};
  \node[hw compute, minimum width=2.2cm] (sel) at (2.4,-0.6) {74HC157\\地址二选一};
  \node[hw control, minimum width=2.2cm] (dip) at (2.4,-2.2) {地址 DIP×4\\（编程）};
  \draw[hw bus] (0,1.2) -- node[midway, above=1pt, hw label, fill=white, inner sep=1pt]{MI} (mar.west);
  \draw[hw bus] (mar.south) -- node[midway, right=1pt, hw label, fill=white, inner sep=1pt]{A 组} (sel.north);
  \draw[hw bus] (dip.north) -- node[pos=0.25, right=1pt, hw label, fill=white, inner sep=1pt]{B 组} (sel.south);
  \draw[hw return] (4.35,-1.55) -- (4.35,-0.75) -- ($(sel.east)+(0,-0.15)$);
  \node[hw label, fill=white, inner sep=1pt] at (4.35,-1.78) {PROG 选择};
  % 两片 189 拼 16×8
  \node[hw memory, minimum width=1.9cm] (ramlo) at (5.6,0.7) {74HC189\\低 4 位};
  \node[hw memory, minimum width=1.9cm] (ramhi) at (8.2,0.7) {74HC189\\高 4 位};
  \draw[hw bus] (sel.east) -- (4.0,-0.6) -- (4.0,0.7) -- (ramlo.west);
  \node[hw label, anchor=west, fill=white, inner sep=1pt] at (4.1,-0.2) {A3--A0};
  \draw[hw bus] (3.7,-0.6) -- (3.7,1.45) -- (8.2,1.45) -- (ramhi.north);
  \fill[BookMuted] (3.7,-0.6) circle (0.05);
  \node[hw label, above=1pt, fill=white, inner sep=1pt] at (6.0,1.45) {地址两片并联};
  % 写数据与写使能
  \draw[hw bus] (0,-1.3) -- (8.2,-1.3) -- (8.2,0.2);
  \draw[hw bus] (5.6,-1.3) -- (5.6,0.2);
  \fill[BookMuted] (5.6,-1.3) circle (0.05);
  \node[hw label, above=1pt, fill=white, inner sep=1pt] at (2.8,-1.3) {写数据（总线）};
  \node[hw label, anchor=east] at (0.15,-2.9) {RI};
  \draw[hw return] (0.3,-2.9) -- (4.4,-2.9) -- (4.4,0.45) -- ($(ramlo.west)+(0,-0.25)$);
  \node[hw label, anchor=west, fill=white, inner sep=1pt] at (4.5,-1.8) {W（低半拍开窗，两片并接）};
  % 读出：189 反相输出经 240 还原并三态上总线
  \node[hw compute, minimum width=2.4cm] (buf) at (6.9,2.6) {74HC240\\反相 + 三态};
  \draw[hw bus] (ramlo.east) -- (6.9,0.7) -- (6.9,2.02);
  \node[hw label, anchor=east, fill=white, inner sep=1pt] at (6.5,-0.05) {反相};
  \draw[hw bus] ($(ramhi.east)+(0,0.15)$) -- (9.5,0.85) -- (9.5,1.8) -- (7.5,1.8) -- (7.5,2.02);
  \draw[hw bus] (buf.west) -- node[midway, above=1pt, hw label, fill=white, inner sep=1pt]{RO：原值上总线} (0,2.6);
  \draw[hw return] (6.9,3.62) -- (buf.north);
  \node[hw label, anchor=west, fill=white, inner sep=1pt] at (7.0,3.35) {RO（经反相）};
\end{pdfdiagram}
\diagramnote{内存模块。MAR 是一片 74HC173，MI 有效的拍末把总线低 4 位锁存为地址；地址经 74HC157 二选一——运行位（A 组）取自 MAR，编程位（B 组）取自 4 只地址 DIP 开关，由 PROG 信号选择——并联送到两片 74HC189 的地址脚。两片 189 各管 4 位：写时总线数据在 RI 的低半拍窗口（W）进入选中字；读时 189 的反相输出经 74HC240 再反相、由 RO 打开三态，原值上总线。RI 与 RO 经不同的使能路径作用于同一组地址，互斥由微码表保证。}

173 的引脚与接法同样列成表。它是四位 D 寄存器，三态输出、双装载使能、高有效清除，一片不多不少正好是一个 MAR：

\begin{longtable}{L{0.16\linewidth}L{0.28\linewidth}L{0.38\linewidth}}
\toprule
173 引脚 & 接向 & 说明 \\
\midrule
D0–D3 & 数据总线低 4 位 & 地址的两个来源都经总线进来 \\\tablerowrule
CP & 整机时钟 & 拍末上升沿锁存 \\\tablerowrule
E1、E2 & 并接，经反相接 MI & MI 有效的拍末装载 \\\tablerowrule
OE1、OE2 & 并接后接地 & 输出常开，直送 189 地址脚 \\\tablerowrule
MR & 整机复位 & 上电清为 0000 \\\tablerowrule
Q0–Q3 & 两片 189 的 A0–A3（经 74HC157） & 点对点专线，不经总线 \\
\bottomrule
\end{longtable}

MAR 的地址有两个来源，但都经同一个动作进入：先由某个驱动者把地址放上总线低 4 位，再由 MI 在拍末锁存。T0 取指时来源是 PC——CO 有效，PC 的 4 位经三态缓冲上总线；访存指令的 T2 来源是 IR 低 4 位——IO 有效，指令里的地址字段上总线。于是内存的全部访问形态可以归纳为一句话：MAR 给地址，RI 与 RO 两个互斥使能决定方向和有无。互斥不是建议而是硬边界：RI 有效的拍 RO 必须无效，否则 RAM 经 240 向总线驱动读出值，而写数据的来源（比如 A 寄存器）同时向总线驱动写入值，两个驱动者争用，总线电平由两路输出级的力气决定，189 收到的可能是两者的混合。本机的微码表里不存在任何一行同时置 RI 和 RO——互斥是在控制字层面被穷举掉的，不靠接线时的运气。

三态边界同样值得说透。RO 无效时，240 输出高阻，RAM 一侧的存储值无论是什么，都不会以任何方式影响总线电平；此时总线由别的驱动者占用，或者在空闲拍无人驱动、由“总体架构：总线、时钟与复位”一节的总线下拉电阻读作 0。RO 有效时，240 是总线上唯一的 RAM 侧驱动者，微码表保证同一拍没有第二个驱动者。读与写的边界因此清晰：读拍是“RO 有效、RI 无效、RAM 驱动总线、某寄存器拍末采样”；写拍是“RI 有效、RO 无效、某寄存器驱动总线、RAM 在写窗口内收数”。下表把全机涉及内存的拍列全。

还有一个结构性的事实顺手点破：MAR 只有一个，RAM 只有一个口。这意味着取指和读写数据永远不可能同拍进行——取指要用 MAR，LDA、STA 也要用 MAR，大家排队。本机每条指令先花两拍取指、再花若干拍执行的多拍结构，一半原因就在这一个 MAR 上；冯·诺依曼结构“程序与数据同处一块内存”的简洁，代价正是访存的串行化。这个代价在 16 字节的规模上无关紧要，但它会在后面每一章的存储器讨论里反复出现。

\begin{longtable}{L{0.20\linewidth}L{0.22\linewidth}L{0.20\linewidth}L{0.24\linewidth}}
\toprule
场合 & 地址来源 & 有效控制线 & 数据方向 \\
\midrule
取指 T1 & MAR（T0 自 PC 锁存） & RO（另有 II、CE） & RAM 到 IR \\\tablerowrule
LDA/ADD/SUB 的 T3 & MAR（T2 自 IR 锁存） & RO（另有 AI 或 BI） & RAM 到 A 或 B \\\tablerowrule
STA 的 T3 & MAR（T2 自 IR 锁存） & RI（另有 AO） & A 到 RAM \\\tablerowrule
其余各拍 & 不变 & 无 & RAM 不访问总线 \\
\bottomrule
\end{longtable}

三态是不是真的好，用一只电阻就能量出来。RO 无效时，把 240 的某个输出脚经 1k$\Omega$ 电阻先后接到 5 V 和地各量一次：真正高阻的脚会被电阻先后拉成 1 和 0；如果电平固执不变，说明还有别的驱动者挂在上面——最常见的嫌疑人是 240 的使能接反、某片 189 的 W 被误接地、或相邻模块的某根输出使能恒有效。这个“电阻拉扯法”对全机每一条三态路径都适用，是“总体架构：总线、时钟与复位”一节总线检查的标准动作之一。

写拍的完整时序走一遍更有体感。以 STA 的 T3 为例，拍首上升沿一过，控制字 AO+RI 稳定：A 寄存器的输出缓冲打开，总线被驱动为 A 的值；MAR 里是上拍锁存的目标地址，早已稳定；时钟转入低半拍后，189 的 W 被拉低，写窗口打开，总线值流进选中字；拍末上升沿到来，时钟转高，W 先随门控电平关闭，控制字随后翻成下一拍的内容，A 的输出缓冲关闭，总线释放。整个写操作里，数据在写窗口打开前已经稳定了半拍，窗口关闭后地址数据才开始变动——建立与保持两条边界都靠“低半拍开窗”这一方法天然满足。

\begin{lstlisting}
// 注释（推进）：按行追踪数据来源、经过的部件和最终写入目标。
// 注释（边界）：只有标出的时钟沿、阶段或异常入口会改变体系结构可见状态。
STA 的 T3 一拍之内（时间从左到右）：
  拍首上升沿   控制字 AO+RI 稳定
  高半拍       A 驱动总线，MAR 地址早已稳定，W 保持高
  转入低半拍   W 拉低，写窗口开，总线值写入 RAM[MAR]
  拍末上升沿   时钟转高，W 先关，控制字再翻，总线释放
\end{lstlisting}

值得注意，MAR 在一条指令的生命里常常要被装两次。取指 T0 装一次（来自 PC），执行期 T2 可能再装一次（来自 IR 的地址字段）——LDA、ADD、SUB、STA 这四条访存指令都是如此。两次装载共用同一个 MI、同一条总线路径，不发生任何冲突，因为它们的用途首尾相接：第一次为了把指令取回来，第二次为了按指令里的地址去取数据。T2 的 IO+MI 与 T0 的 CO+MI 在微码表里是两行不同的控制字，驱动者不同、锁存者相同，这正是“总线轮值、寄存器复用”的标准画面。

最后回答一个容易被跳过的问题：MAR 为什么只需要 4 位。地址空间是 16 字节，4 位二进制正好数完 0 到 15，多一位浪费、少一位不够。这个“刚好”是整机规格联动的结果：PC 是 4 位、MAR 是 4 位、指令的地址字段是 4 位、189 的地址脚是 4 根，四处必须一致，任何一处扩到 5 位都会强迫其余三处跟着动。也正因如此，这台机器连存储与设备事务相关章节讲的地址译码器都不需要——16 个字就是全部存储，选中就是全部译码。将来第一轮扩展（比如 32 字节内存）要动的正是这一组联动的 4 位，“编程、调试与扩展”一节会回到这个话题。

\topic{OUT 显示把寄存器值变成人可读结果}
OUT 寄存器是全机唯一的“只进不出”状态：它在 OUT 指令的 T2 拍末锁存总线值（那一拍 AO 与 OI 同时有效，A 驱动总线，OUT 收数），此后一直驱动显示电路，直到下一条 OUT 指令执行才更新。显示的内容因此是“最近一次 OUT 时 A 的值”，而不是 A 的实时值——这个区别调试时要想清楚：程序继续跑，A 早已变了，数码管上留着的仍是 OUT 那一刻的快照。器件用一片 74HC377（八位 D 触发器，带低有效装载使能）正好：使能端接反相后的 OI，输出常开接显示。377 没有清除端也不要紧，上电显示一个随机值，第一条 OUT 执行后自然被覆盖。

显示的第一档方案是 8 只 LED 各串一只 330 欧电阻，直接挂在 377 的 8 根输出上，亮为 1、灭为 0。LED 有极性：长脚是正极，接 377 的输出端，短脚经电阻接地，接反了不亮也不坏，调过来就是。二进制读数需要心算，但它最忠实——每一位都对应一根真实的导线，没有译码器可以怀疑。第二档方案是十六进制数码管：两位显示 00 到 FF，需要能把 4 位译成包括 A–F 在内的 16 种段码的驱动器，DM9368 这类十六进制译码驱动一体的芯片两片即可；也可以用一片小 EEPROM 自制译码——把 4 位二进制当地址、7 段码当内容烧进去，这是把微码“查表”的思想在显示端再用一次。这里有一个常见坑要点名：74HC47、74HC48 是 BCD 译码器，输入 1010 到 1111 时输出的不是 A–F，而是没有定义的段组合；拿它做十六进制显示，数值超过 9 就开始出怪字。至于十进制显示，需要先做二进制到 BCD 的转换，对本机得不偿失，不做。

限流电阻的取值是算出来的，不是抄来的。红色 LED 的正向压降约 2 V，74HC377 输出高电平带载后约 4.5 V，目标电流取 5 mA——足够亮，又不逼近 HC 器件单脚 25 mA 的输出上限——电阻值就是 (4.5 − 2) / 0.005 = 500 欧，标准件里取 470 欧或 330 欧都成立：330 欧对应电流约 7.6 mA，亮度余量更足，这也是本书一律写 330 欧的原因。数码管的每一段就是一只 LED，同样的算法再算一遍，七段加小数点共八只限流电阻，一只都省不掉；省掉的代价是段亮度随字形变化——显示 1 时两段分一份电流显得亮，显示 8 时七段分显得暗。

选 EEPROM 自制译码这条路的，烧片时直接照抄下面这张十六进制字形表。段码按 g、f、e、d、c、b、a 的顺序排成 7 位，1 表示该段点亮，字形就是常见数码管上的 0 到 F：

\begin{longtable}{L{0.10\linewidth}L{0.16\linewidth}L{0.10\linewidth}L{0.16\linewidth}L{0.10\linewidth}L{0.16\linewidth}}
\toprule
数字 & 段码 & 数字 & 段码 & 数字 & 段码 \\
\midrule
0 & 0111111 & 6 & 1011111 & C & 0111001 \\\tablerowrule
1 & 0000110 & 7 & 0000111 & D & 1011110 \\\tablerowrule
2 & 1011011 & 8 & 1111111 & E & 1111001 \\\tablerowrule
3 & 1001111 & 9 & 1101111 & F & 1110001 \\\tablerowrule
4 & 1100110 & A & 1110111 & & \\\tablerowrule
5 & 1101101 & B & 1111100 & & \\
\bottomrule
\end{longtable}

表里的 B 和 D 用小写字形（b、d）是数码管的通行做法：大写 B 与 8、大写 D 与 0 在七段上无法区分，小写字形虽然委屈了字母，却保住了可读性。烧写时地址就是 4 位输入值，内容就是对应段码，16 行写完即收工——译码在这里不是逻辑设计，而是又一次填表。

最重要的一条边界是：显示只是观察，不能影响总线。OUT 的输出去向只有显示负载，整机里不存在任何一根“OUT 输出回总线”的控制线，所以显示部分无论怎么接——LED、数码管、示波器探头，乃至接错一个脚——都不可能改变总线上的任何电平。这条单方向边界值得在接线完成后当场核实：用万用表通断档量 OUT 的每个输出脚，它们应当只连到显示负载，不连总线。

数码管驱动还要做一个小小的架构选择：静态还是扫描。静态驱动是每位数管一套译码、常亮；扫描驱动是多位共用一套译码、轮流快速点亮，靠视觉暂留看起来同时亮。商用设备爱用扫描，因为省线省芯片；本机必须用静态，理由就在调试场景里——单步时钟和 HLT 停机时，扫描的驱动时钟也停了，数码管会只剩一位亮着甚至全灭，而静态显示在任何状态下都忠实保持 OUT 寄存器的内容。显示电路的存在是为了让机器“随时可看”，一个在停机时罢工的显示器，等于在最需要观察的时刻闭了眼。

上电瞬间 OUT 寄存器里是随机值，显示一片乱码，属正常现象——377 没有清除端，第一条 OUT 指令执行后乱码自然被正确的值覆盖。若执意让显示也从已知值出发，把 377 换成两片 74HC173（输出使能接地、装载使能接反相后的 OI、清除端接整机复位）即可，代价是多占一片的位置；多数搭建者的选择是不在意这片刻的乱码，因为程序的第一条显示指令之前，显示内容本来就没有约定。

OI 的时序与别的装载使能没有两样：OUT 指令的 T2，AO 让 A 驱动总线，OI 让 377 在拍末上升沿把总线值锁存。同一拍里 377 既不驱动总线也不影响任何别的模块，所以 OUT 是全部指令里最安静的一条——总线给它送数据，它什么也不用还。

摆放 OUT 指令的位置是编程风格问题，先在这里立个规矩：凡是想让操作者看到的中间结果或最终结果，显式地用一条 OUT 送出来；凡是要停机的地方，HLT 之前先放一条 OUT，否则停机后显示的是上一次 OUT 的旧值，结果等于没显示。“编程、调试与扩展”一节的两个程序都按这个规矩写，读程序时可以对照。

\begin{longtable}{L{0.22\linewidth}L{0.24\linewidth}L{0.22\linewidth}L{0.20\linewidth}}
\toprule
显示方案 & 器件 & 能显示什么 & 关键问题 \\
\midrule
LED 8 只 & 330 欧电阻 8 只 & 二进制 8 位 & 最忠实，读数要心算 \\\tablerowrule
十六进制数码管两位 & DM9368 两片 & 00–FF & 芯片较老，注意共阴共阳与限流 \\\tablerowrule
EEPROM 自制译码 & 28C16 一片加数码管 & 任意自定义段码 & 查表思想的又一次复用 \\\tablerowrule
74HC47 两片 & BCD 译码 & 只能正确显示 0–9 & 输入 A–F 出怪段，不适合十六进制 \\
\bottomrule
\end{longtable}

\topic{编程模式把程序逐字节拨进 RAM}
编程模式要解决的本质问题，是让人的手临时取代控制单元，在不与运行电路冲突的前提下操作 189。模式开关（一只拨动开关，以下称 PROG）拨到编程位时做三件事。第一，切断时钟：T 状态计数器停在原地，全部控制线随之静止，RI、RO 都为无效——RAM 既不被读也不被写，总线上没有任何运行中的驱动者。第二，74HC157 的选择端切到 DIP 地址开关一侧，189 的地址来源从 MAR 换成手指。第三，使能数据开关侧的 74HC244 和写按钮的通路：8 只数据开关经 244 驱动总线，按钮脉冲经 74HC14 整形后接 189 的 W。此后每装一个字节都是同一组动作：拨好 4 位地址、拨好 8 位数据、按一下写按钮。按钮抖动对这种写法其实无害——重复写同一个地址同一个值，结果与写一次相同；真正要紧的是次序，先把开关全部拨到位再按写，松手之后才允许动开关，否则在地址切换的瞬间，数据可能被写进别的字。

开关的电平方向要先定死再动手。本机的约定是开关拨到 ON 一侧时对应位为 1，实现上有两种接法：开关串在 5 V 与信号线之间、信号线经 10k$\Omega$ 下拉电阻接地，断开为 0、闭合为 1，拨上去就是 1，最合直觉；或者开关一端接地、信号线经 10k$\Omega$ 上拉，电平正好反过来。两种都对，但 12 只开关必须用同一种接法，并且在面板上用贴纸标明哪边是 1、哪只是最高位。高位低位标反是编程模式的头号人为故障：把 D0 当成 D7，写进去的字节正好按位倒序，程序的表现完全无法理喻，而电路没有任何毛病。

74HC157 的选择逻辑一行就能写完：选择端接 PROG 模式信号，运行位时 4 路输出取自 A 组——MAR 的 Q0–Q3；编程位时取自 B 组——4 只地址开关。157 的使能端接地常开，输出直送两片 189 的地址脚。MAR 的输出照常挂在 A 组输入上，不需要断开——选择器本身就隔离了两路来源，这正是用选择器而不用跳线的理由。写按钮一侧，74HC14 加 RC（10k$\Omega$ 与 0.1$\mu$F，时间常数 1 ms）把按钮的机械抖动滤成干净的一次跳变，输出接两片 189 的 W；编程模式下 RI 恒无效，运行电路对 W 没有任何影响。

把两种模式下 189 各引脚的信号来源对照列出来，切换电路的设计就只剩照表施工：

\begin{longtable}{L{0.20\linewidth}L{0.30\linewidth}L{0.32\linewidth}}
\toprule
信号 & 运行模式来源 & 编程模式来源 \\
\midrule
189 地址 A3–A0 & MAR（经 74HC157 A 组） & 4 只地址开关（经 157 B 组） \\\tablerowrule
189 数据输入 & 数据总线（写拍由 AO 驱动） & 数据总线（8 只数据开关经 244 驱动） \\\tablerowrule
189 写使能 W & RI 与时钟低半拍相与后取反 & 写按钮经 74HC14 整形 \\\tablerowrule
240 输出使能 & RO 反相 & 保持无效，RAM 不上总线 \\\tablerowrule
总线其它驱动者 & 各模块按控制字轮值 & 全部静默（时钟已停） \\\tablerowrule
T 状态计数器 & 逐拍计数 & 冻结 \\
\bottomrule
\end{longtable}

模式切换本身要不要担心毛刺？不太要。切换发生在时钟停止的状态下，全机没有任何边沿在走，157 换源、244 开关、按钮通断，都是静态电平的此消彼长；毛刺之所以可怕，是因为它会冒充时钟边沿触发装载，而此刻没有装载对象醒着。唯一真正有破坏力的动作是在运行中把 PROG 拨过去——时钟被拦腰截断的那一拍可能窄到反常，触发器的行为不再可预期。所以纪律只有一条：先停机，后切换；先切回，后复位，再启动。

整个流程先画成一张图，再按步展开。

\begin{pdfdiagram}
  \node[hw control, minimum width=4.6cm] (f1) at (0,0) {停机（HLT 或拨停时钟）};
  \node[hw control, minimum width=4.6cm] (f2) at (0,-1.35) {PROG 拨到编程位};
  \node[hw state, minimum width=4.6cm] (f3) at (0,-2.7) {拨 4 位地址（DIP 开关）};
  \node[hw state, minimum width=4.6cm] (f4) at (0,-4.05) {拨 8 位数据（DIP 开关）};
  \node[hw compute, minimum width=4.6cm] (f5) at (0,-5.4) {按一下写按钮（写脉冲）};
  \node[hw compute, minimum width=4.6cm] (f6) at (0,-6.75) {还有字节要装？};
  \node[hw control, minimum width=4.6cm] (f7) at (0,-8.1) {PROG 拨回运行位};
  \node[hw control, minimum width=4.6cm] (f8) at (0,-9.45) {按复位（PC←0、T←T0、IR 清）};
  \node[hw state, minimum width=4.6cm] (f9) at (0,-10.8) {启动时钟，从 T0 运行};
  \draw[hw bus] (f1.south) -- (f2.north);
  \draw[hw bus] (f2.south) -- (f3.north);
  \draw[hw bus] (f3.south) -- (f4.north);
  \draw[hw bus] (f4.south) -- (f5.north);
  \draw[hw bus] (f5.south) -- (f6.north);
  \draw[hw bus] (f6.south) -- node[midway, right=1pt, hw label, fill=white, inner sep=1pt]{否，装完} (f7.north);
  \draw[hw bus] (f7.south) -- (f8.north);
  \draw[hw bus] (f8.south) -- (f9.north);
  \draw[hw return] (f6.west) -- (-3.6,-6.75) -- (-3.6,-2.7) -- (f3.west);
  \node[hw label, above=1pt, fill=white, inner sep=1pt] at (-2.95,-6.75) {是};
  \node[hw label, rotate=90] at (-3.82,-4.7) {下一字节};
\end{pdfdiagram}
\diagramnote{编程模式流程。先停机再切换：PROG 拨到编程位后时钟冻结、74HC157 切到 DIP 地址、数据开关经 244 上总线、写按钮接管 189 的 W；每个字节都是“拨地址、拨数据、按写按钮”三步的循环。装完切回运行位后，复位不能省——它把 PC、T、IR 清到已知值，程序才从地址 0 的 T0 起跑。}

完整流程七步：

\begin{lstlisting}
// 注释（推进）：按步骤依次执行，前一步的稳定结果是后一步的前提。
// 注释（边界）：每一步都保留可观察读回；不满足预期就停在当前层排错。
1. 停机：     程序自然跑到 HLT，或直接把时钟开关拨到停
2. 切模式：   PROG 拨到编程位（时钟断、157 切到 DIP、按钮通路开）
3. 设地址：   4 只地址开关拨出目标地址（0x0 到 0xF）
4. 设数据：   8 只数据开关拨出要写入的字节
5. 写入：     按一下写按钮，松开；回到第 3 步装下一字节
6. 切回：     全部字节装完，PROG 拨回运行位
7. 复位运行： 按复位（PC 清零、T 计数器回 T0、IR 清空），再启动时钟
\end{lstlisting}

第 7 步的复位不能省。编程期间 T 状态计数器和 PC 停在任意值，IR 里是上次停机时的指令残影；不复位就启动时钟，机器会从某个未知的 T 状态、某个未知的 PC 接着走，第一条执行的指令几乎必然不是程序的起点，后面的行为也就没有讨论意义。复位把 PC、T 计数器、IR 清到已知值，等于把机器放回一条确定的出发线上，程序从地址 0 开始、从 T0 开始，一拍一拍都可以对照“控制单元：指令译码与微码”一节的微操作表推演。

下面是一段 5 字节的写入示例。程序功能是把立即数 6 与存在 0x4 单元的 7 相加，显示结果后停机：

\begin{longtable}{L{0.10\linewidth}L{0.14\linewidth}L{0.18\linewidth}L{0.14\linewidth}L{0.30\linewidth}}
\toprule
次序 & 地址开关 & 数据开关 & 字节 & 含义 \\
\midrule
1 & 0000 & 0101 0110 & \code{0x56} & \code{LDI 0x6}：A ← 6 \\\tablerowrule
2 & 0001 & 0010 0100 & \code{0x24} & \code{ADD 0x4}：A ← A + RAM[0x4] \\\tablerowrule
3 & 0010 & 1110 0000 & \code{0xE0} & \code{OUT}：显示 A \\\tablerowrule
4 & 0011 & 1111 0000 & \code{0xF0} & \code{HLT}：停机 \\\tablerowrule
5 & 0100 & 0000 0111 & \code{0x07} & 数据：加数 7 \\
\bottomrule
\end{longtable}

5 个字节拨完、切回运行、复位、启动时钟，预期结果是显示 0000 1101，即 0x0D，十进制 13：LDI 把 6 装入 A，ADD 从 0x4 单元取出 7 相加，OUT 把 13 锁存进显示，HLT 停住时钟。如果显示的不是这个值，先别急着怀疑运行电路——编程模式下最容易错的是人：拨错的开关、漏按的按钮、忘了复位。逐字节核对与逐拍追踪的完整流程在“编程、调试与扩展”一节展开。

有人会用一种正当的疑惑看待这套手工流程：既然程序固定，为什么不把它烧进 ROM，省去每次上电拨开关？答案是整机规格里“程序与数据同处一块 RAM”这半句——STA 指令要写内存，累加、乘法循环都要靠可写的单元保存中间结果，ROM 撑不起这些程序。编程模式不是权宜之计，是这台机器唯一合理的装程序方式。时间预算也算一笔：熟练后每字节约十秒，装满 16 字节三分钟，再加一遍读回核对，总共不超过十分钟；“编程、调试与扩展”一节的两个程序连代码带数据都装在这 16 字节里，一次拨全。

把常见的人为错误集中在这里，装程序时对照自查。一是地址开关高低位拨反，字节按位倒序写入；二是数据开关拨错一位，程序能跑但结果错，最难怀疑；三是拨完忘记按按钮，该单元仍是旧值或上电随机值；四是按钮按了两次本无妨，但两次之间动了开关，数据进了相邻单元；五是切回运行位后忘了复位，机器从随机的 T 状态和 PC 起跑。这五种错里没有一种是电路故障——先用 DIP 开关装程序的另一重收获正在这里：把人会犯什么错，在最小规模上全部经历一遍。

读回核对的方法也说清。PROG 位下时钟是停的，控制单元帮不上忙，但核对只需要回答一个问题：这个地址里现在是什么。把地址开关拨到待核对的字节，再临时给 240 的使能一个有效电平——用一只带引线的按钮把 RO 反相后的节点接地即可——该字节就会经 240 出现在总线上，由“总体架构：总线、时钟与复位”一节的总线显示 LED 读出，核对完松开。这个动作本质上是手工重演了一遍 LDA 的 T3：地址给好、RO 打开、存储值上总线，只是每一步都来自手指。把 16 个地址逐个读一遍，与想写入的清单逐字节比对，比运行出错后再回头猜要省一个数量级的时间。

本节收尾给一份检查清单，按从上到下的顺序做，任何一项不符都先停下来排因，再往下走。

\begin{longtable}{L{0.22\linewidth}L{0.38\linewidth}L{0.28\linewidth}}
\toprule
检查点 & 方法 & 预期结果 \\
\midrule
189 输出极性 & 预存 \code{0xA5}，RO 有效时量总线 & 总线读得 \code{0xA5}，而非 \code{0x5A} \\\tablerowrule
RO 无效的三态 & RO 无效时量 240 输出脚 & 高阻，总线电平由别人决定 \\\tablerowrule
MAR 锁存 & 单步到 T1，总线内容已换 & 189 地址脚仍保持 T0 锁存的地址 \\\tablerowrule
写窗口边界 & 执行 STA 后读相邻单元 & 只有目标地址被改写 \\\tablerowrule
OUT 显示 & LDA 一个已知值后执行 OUT & LED 逐位等于该值 \\\tablerowrule
位序方向 & 预存 \code{0x01}，读出时看总线 D0 & 是最低位亮，而非最高位 \\\tablerowrule
编程读回 & PROG 下逐地址读出 16 字节 & 与写入清单逐字节一致 \\\tablerowrule
切换纪律 & 运行中误拨 PROG 后复位重跑 & 复位后行为恢复正常 \\\tablerowrule
5 字节程序 & 按上表装入后复位运行 & 显示 \code{0x0D} 后停机 \\
\bottomrule
\end{longtable}

至此，机器有了记忆也有了脸面：RAM 能存能取，MAR 把地址看住，OUT 把结果留下，编程模式让人能把程序一个字节一个字节地请进内存。这一节反复出现的词是边界——RI 与 RO 的互斥边界、三态的高阻边界、显示的单向边界、模式切换的先后边界。“控制单元：指令译码与微码”一节轮到控制单元登场，它的全部工作，就是按一张表逐拍打开和关上这一节装好的每一扇门。

\section{四、控制单元：指令译码与微码}

控制单元是全机的“读表机器”。它不再引入新的数据通路部件，只回答一个问题：当前这一拍，16 根控制线各自是有效还是无效。答案来自一张 128 行乘 16 位的表，行地址由两块状态拼成——IR 里的 opcode 说明“现在执行的是哪条指令”，T 状态计数器说明“现在走到这条指令的第几拍”。

把控制单元的输入输出先登记成表——它的接口之窄，本身就是微码方案的优点：

\begin{longtable}{L{0.16\linewidth}L{0.30\linewidth}L{0.36\linewidth}}
\toprule
方向 & 信号 & 说明 \\
\midrule
入 & 时钟、复位 & 节拍推进与清零，“总体架构：总线、时钟与复位”一节的约定 \\\tablerowrule
入 & 总线 D0–D7 & 仅 T1 供 IR 装入指令 \\\tablerowrule
入 & CF、ZF & 仅 JC、JZ 的 J 选通使用 \\\tablerowrule
出 & 16 根控制线 & 全机所有模块的使能与方式 \\
\bottomrule
\end{longtable}

\topic{控制单元 = IR + T 状态计数器 + 微码 ROM}
IR 在 T1 拍末从总线装入指令字节（II 有效），此后整条指令的执行期间它保持不变——这是它与 A、B 这些“数据寄存器”的本质区别：数据寄存器装的是会被加工的值，IR 装的是要管好几拍的命令。本机用两片 74HC173 拼成 8 位 IR，两片的分工不是随意的高低位，而是按去向划的：高片存 opcode（IR7–IR4），输出使能接地常开，4 根输出线直接送往微码 ROM 的地址脚；低片存地址或立即数字段（IR3–IR0），输出使能由 IO 控制，IO 有效时这 4 位经三态上总线低 4 位。两片共用同一个时钟与装载使能（E1、E2 并接，经反相接 II），清除端接整机复位——与 MAR 一样，MR 是高有效，需经一级反相接低有效的复位线。

这个分工里藏着一条重要的接线纪律：opcode 永远不经过总线。IR 高片到 ROM 地址脚是一条点对点专线，就像 MAR 到 RAM 的那条一样——凡是“控制信息”，在本机里都走专线，总线只运“数据”。好处不止是不争用：取指完成的同一瞬间，opcode 已经摆在 ROM 地址脚上，T2 的控制字不经过任何中转就开始译出。按“总体架构：总线、时钟与复位”一节的总线约定，总线高 4 位在无人驱动时读为 0，所以 IO 有效时总线呈现的是“高 4 位为 0、低 4 位是指令字段”，LDI 装入 A 的自然是零扩展后的 4 位立即数，JMP 装给 PC 的也正好是 4 位地址，不多不少。

两片 173 的接法对照如下，与“内存与输出模块”一节 MAR 那片 173 是同一片型、三种用法：

\begin{longtable}{L{0.16\linewidth}L{0.30\linewidth}L{0.36\linewidth}}
\toprule
引脚 & IR 高片（opcode） & IR 低片（地址/立即数） \\
\midrule
D0–D3 & 总线 D4–D7 & 总线 D0–D3 \\\tablerowrule
CP & 整机时钟 & 整机时钟 \\\tablerowrule
E1、E2 & 并接，经反相接 II & 并接，经反相接 II \\\tablerowrule
OE1、OE2 & 接地，常开 & 并接，经反相接 IO \\\tablerowrule
MR & 整机复位 & 整机复位 \\\tablerowrule
Q0–Q3 & 微码 ROM 地址高 4 位 & 总线低 4 位 \\
\bottomrule
\end{longtable}

注意两片唯一的差别在输出使能：高片常开、低片受控。这个差别在排错时是个趁手的判据——如果 opcode 出现“时有时无”，先查是不是把高片的 OE 也接到了 IO 上；如果 IO 拍总线低 4 位不动，先查低片的 OE 是不是忘了接。

T 状态计数器是一片 74HC161，只用低 3 位输出 Q2、Q1、Q0——它们就是拍号的三位二进制编码，000 是 T0，111 是 T7。计数使能 CEP、CET 接高电平，时钟脚经一级反相接入，每个时钟下降沿加一——T 状态在拍中推进，新控制字到下一个上升沿前半拍就已稳定，装载沿采到的因此是毫不含糊的电平。回绕用同步装载实现：74HC11 的一个三输入与门译出 111，接 161 的同步装载使能 PE（低有效），并行数据输入 D3–D0 全部接地；计数到 T7 的那一拍，装载条件成立，拍末下降沿把 000 装进去，下一拍回到 T0。第 4 位输出 Q3 悬空不用，它永远不会变成 1。

把八个拍与计数器状态、回绕动作的对应列出来，单步调试时对照 Q2–Q0 就知道机器走到哪一拍：

\begin{longtable}{L{0.10\linewidth}L{0.16\linewidth}L{0.24\linewidth}L{0.30\linewidth}}
\toprule
拍 & Q2 Q1 Q0 & ROM 地址低 3 位 & 拍末发生什么 \\
\midrule
T0 & 0 0 0 & 000 & 计数加一 \\\tablerowrule
T1 & 0 0 1 & 001 & 计数加一 \\\tablerowrule
T2 & 0 1 0 & 010 & 计数加一 \\\tablerowrule
T3 & 0 1 1 & 011 & 计数加一 \\\tablerowrule
T4 & 1 0 0 & 100 & 计数加一 \\\tablerowrule
T5 & 1 0 1 & 101 & 计数加一 \\\tablerowrule
T6 & 1 1 0 & 110 & 计数加一 \\\tablerowrule
T7 & 1 1 1 & 111 & PE 有效，装入 000 回 T0 \\
\bottomrule
\end{longtable}

161 的其余引脚一一交代：CP 经一级反相接整机时钟，于是计数在时钟下降沿发生；CEP、CET 并接高电平；PE 接 74HC11 译出的 T7（低有效，故与门输出要反相一次——用 74HC11 加 74HC04 的一门，或干脆用一片 74HC10 三输入与非门，省掉反相器）；D0–D3 接地；MR 接整机复位；Q3 与进位输出 TC 悬空。整片 161 在本机里扮演的角色等价于一个“八状态的模 8 计数器”，只是用同步装载拼出来的模 8，比专用模 8 器件更容易买到。

这里必须强调用同步装载而不用异步清零的理由。若把译出的 T7 接去 161 的异步清零端 MR，T7 这个状态只存在译码延迟加清零延迟那么几十纳秒，随即被抹掉——T7 行的控制字来不及稳定，依赖 T7 的扩展指令就全部出错；更糟的是这个窄脉冲本身是个毛刺源。同步装载让 T7 是和别的拍一样完整的一拍，回绕发生在拍末边沿，干干净净。161 的 MR 另有正经用途：接整机复位，保证上电和每次复位后第一拍是 T0。

把这片 161 和它的回绕译码画成结构图，引脚去向一并标出：

\begin{pdfdiagram}
  % 计数器本体
  \node[hw state, align=center, minimum width=3.2cm] (cnt) at (3.6,0.5) {74HC161\\CEP、CET 接高，D 接地};
  % 时钟与复位
  \node[hw control, align=center] (clk) at (0,2.0) {整机时钟};
  \draw[hw return] (clk.east) -- (1.5,2.0) -- (1.5,0.5) -- (cnt.west);
  \node[hw label, align=center, anchor=east] at (1.4,1.25) {CP 经反相\\下降沿加一};
  \node[hw control, align=center] (rst) at (0,-0.9) {整机复位};
  \draw[hw return] (rst.east) -- node[midway, above=1pt, hw label, fill=white, inner sep=1pt]{MR} (3.6,-0.9) -- (cnt.south);
  % Q 输出：一路去 ROM 地址，一路去 T7 译码
  \node[hw memory, align=center, minimum width=2.8cm] (romd) at (8.6,0.5) {微码 ROM\\地址低 3 位};
  \draw[hw bus] (cnt.east) -- node[pos=0.35, above=1pt, hw label, fill=white, inner sep=1pt]{Q2--Q0（拍号）} (romd.west);
  \node[hw compute, align=center, minimum width=3.2cm] (dec) at (8.6,2.4) {T7 译码\\与门译 111 加反相};
  \draw[hw bus] (cnt.north) -- (3.6,2.4) -- (dec.west);
  \node[hw label, fill=white, inner sep=1pt] at (5.4,2.66) {Q2Q1Q0 = 111};
  % 同步装载回绕
  \draw[hw return] (dec.south) -- (8.6,1.6) -- (4.6,1.6) -- ($(cnt.north)+(1.0,0)$);
  \node[hw label, fill=white, inner sep=1pt] at (6.6,1.86) {PE：拍末装入 000};
  \node[hw label, align=center] at (7.4,-1.3) {同步装载：T7 是完整一拍\\回绕发生在拍末下降沿};
\end{pdfdiagram}
\diagramnote{T 状态计数器结构。74HC161 的 CEP、CET 接高电平，时钟经一级反相接 CP——每个下降沿加一，新拍号到装载沿前半拍已经稳定；Q2--Q0 直接作微码 ROM 地址的低 3 位。回绕用同步装载：三输入与门译出 111（T7），反相后接 PE，并行输入 D3--D0 全接地，T7 拍末的下降沿把 000 装入，下一拍回到 T0——T7 因此是和别的拍一样完整的一拍，不像接异步清零那样只剩几十纳秒的窄脉冲。MR 另有正经用途：接整机复位，保证上电和每次复位后第一拍是 T0。}

微码 ROM 是这张表的物理载体。表有 128 行，行地址 7 位：高 4 位接 IR 的 opcode，低 3 位接 T 计数器的 Q2–Q0，即 ROM 地址 = opcode 乘 8 加 T。每条指令因此占据连续 8 行：第 0、1 行是取指，第 2–7 行是执行拍，本机最长的指令只用到第 4 行。每行 16 位，正好对应 16 根控制线，用两片 8 位 EEPROM（AT28C64、28C16 均可）并联地址实现：两片共用全部 7 根地址线，高片输出控制字的高 8 位（HLT 到 AO），低片输出低 8 位（BI 到 OI），16 根控制线一次出齐，一根不浪费。EEPROM 的片选与输出使能都接地常开——控制字必须随时有效，控制单元没有休息的权利；用不到的高位地址脚全部接地，28C64 有 8K 行，本机只用其中 128 行，余下的空间留给扩展。

三块部件的接线关系画成图，就是一台查表机器。

\begin{pdfdiagram}
  % 两块状态：IR 高片与 T 计数器
  \node[hw state, minimum width=2.6cm] (irh) at (0,1.6) {IR 高片\\opcode（IR7--IR4）};
  \node[hw state, minimum width=2.6cm] (tcnt) at (0,-0.6) {74HC161\\T（Q2--Q0）};
  \draw[hw bus] (-3.0,1.6) -- node[midway, above=1pt, hw label, fill=white, inner sep=1pt]{总线（II）} (irh.west);
  \draw[hw return] (-3.0,-0.6) -- node[midway, above=1pt, hw label, fill=white, inner sep=1pt]{时钟/复位} (tcnt.west);
  % 地址拼接干线
  \draw[hw bus] (irh.east) -- (3.4,1.6);
  \draw[hw bus] (tcnt.east) -- (3.4,-0.6);
  \draw[draw=BookMuted, line width=0.62pt] (3.4,1.6) -- (3.4,-1.0);
  \fill[BookMuted] (3.4,1.6) circle (0.05);
  \fill[BookMuted] (3.4,-0.6) circle (0.05);
  \node[hw label, rotate=90] at (3.14,0.25) {地址 A6--A0};
  % 两片 EEPROM 出 16 根控制线
  \node[hw memory, minimum width=3.0cm] (romh) at (5.8,1.0) {EEPROM 高片\\bit15--8：HLT…AO};
  \node[hw memory, minimum width=3.0cm] (roml) at (5.8,-1.0) {EEPROM 低片\\bit7--0：BI…OI};
  \draw[hw bus] (3.4,1.0) -- (romh.west);
  \draw[hw bus] (3.4,-1.0) -- (roml.west);
  \fill[BookMuted] (3.4,1.0) circle (0.05);
  \fill[BookMuted] (3.4,-1.0) circle (0.05);
  \draw[hw bus] (romh.east) -- ++(0.8,0);
  \draw[hw bus] (roml.east) -- ++(0.8,0);
  % J 的条件选通
  \node[hw control, minimum width=3.4cm] (jgate) at (4.4,-2.9) {J 选通逻辑\\JC 看 CF、JZ 看 ZF、JMP 直通};
  \draw[hw return] (7.7,-1.0) -- (7.7,-2.9) -- (jgate.east);
  \fill[BookTeal] (7.7,-1.0) circle (0.05);
  \node[hw label, above=1pt, fill=white, inner sep=1pt] at (6.9,-2.62) {J};
  \node[hw label, anchor=east] at (-0.2,-2.9) {CF、ZF（Flags）};
  \draw[hw return] (-0.1,-2.9) -- (jgate.west);
  \draw[hw return] ($(irh.west)+(0,-0.2)$) -- (-2.0,1.4) -- (-2.0,-2.7) -- (3.7,-2.7) -- (3.7,-2.37);
  \node[hw label, rotate=90] at (-2.2,-1.75) {c7、c8（译码）};
  \draw[hw return] (jgate.south) -- ++(0,-0.5);
  \node[hw label, anchor=west, fill=white, inner sep=1pt] at (4.5,-3.65) {J\_PC → PC 装载端};
\end{pdfdiagram}
\diagramnote{控制单元。IR 高片的 opcode（4 位）与 74HC161 的 T 状态（Q2--Q0，3 位）拼成 7 位地址，并联送到两片 EEPROM，每片出 8 位，合起来正好 16 根控制线；地址 = opcode×8+T，每条指令占连续 8 行。J 是唯一的例外：它出 ROM 后还要过一级选通——c7、c8 由 opcode 译出，JC 时看 CF、JZ 时看 ZF、JMP 时直通——选通是纯组合逻辑，控制字整拍稳定，拍末采到的是确定电平。}

地址怎么拼、表怎么查，再单画一张示意。

\begin{pdfdiagram}
  % 地址条：高 4 位 opcode + 低 3 位 T
  \node[hw state, minimum width=2.2cm] (opf) at (0.3,1.5) {opcode（4 位）};
  \node[hw state, minimum width=2.2cm] (tf) at (0.3,0.0) {T（3 位）};
  \draw[hw bus] (opf.east) -- (2.6,1.5);
  \draw[hw bus] (tf.east) -- (2.6,0.0);
  \draw[draw=BookMuted, line width=0.62pt] (2.6,1.5) -- (2.6,0.0);
  \fill[BookMuted] (2.6,1.5) circle (0.05);
  \fill[BookMuted] (2.6,0.0) circle (0.05);
  % ROM 与命中行区
  \node[hw memory, minimum width=2.4cm, minimum height=4.6cm] (rom) at (5.2,0) {微码 ROM\\128 行\\×16 位};
  \draw[hw bus] (2.6,0.75) -- (4.0,0.75);
  \fill[BookMuted] (2.6,0.75) circle (0.05);
  \node[hw label] at (5.2,-2.75) {地址 = opcode×8+T};
  \node[hw compute, minimum width=3.6cm, align=left, inner xsep=6pt] (hit) at (9.35,0.9) {opcode = 0x2（ADD）\\行 0x10--0x17\\T0 0x4008　T1 0x1810\\T2 0x4400　T3 0x1080\\T4 0x0242　T5--7 0x0000};
  \draw[hw bus] ($(rom.east)+(0,0.9)$) -- node[midway, above=1pt, hw label, fill=white, inner sep=1pt]{命中} (hit.west);
\end{pdfdiagram}
\diagramnote{微码查找。7 位地址由高 4 位 opcode 与低 3 位 T 拼成（地址 = opcode×8+T），每条指令占据 ROM 中连续 8 行：第 0、1 行是所有指令相同的取指（CO+MI、RO+II+CE），第 2 行起按指令分化，用不满的行填 0x0000。图中展开 ADD（opcode 0x2）的 8 行：T2 送操作数地址、T3 操作数进 B、T4 求和写 A 并记标志，其余三行空转。}

把命中区的 8 行再逐行摊开，行地址、控制字、动作三者一一对应：

\begin{pdfdiagram}
  % 左侧：ROM 中 ADD 占据的区域
  \node[hw memory, align=center, minimum width=3.0cm, minimum height=5.6cm] (rom) at (0,-2.45) {微码 ROM\\ADD 区\\0x10--0x17};
  % 右侧：8 行逐行展开（anchor=west，行宽随文字自适应）
  \node[hw control, align=center, anchor=west] (r0) at (5.1,0)    {0x10：0x4008 = CO+MI（T0 取指）};
  \node[hw control, align=center, anchor=west] (r1) at (5.1,-0.7) {0x11：0x1810 = RO+II+CE（T1 取指）};
  \node[hw compute, align=center, anchor=west] (r2) at (5.1,-1.4) {0x12：0x4400 = IO+MI（T2 地址进 MAR）};
  \node[hw compute, align=center, anchor=west] (r3) at (5.1,-2.1) {0x13：0x1080 = RO+BI（T3 操作数进 B）};
  \node[hw compute, align=center, anchor=west] (r4) at (5.1,-2.8) {0x14：0x0242 = $\Sigma$O+AI+FI（T4 求和写 A）};
  \node[hw state, align=center, anchor=west] (r5) at (5.1,-3.5)   {0x15：0x0000（T5 空拍）};
  \node[hw state, align=center, anchor=west] (r6) at (5.1,-4.2)   {0x16：0x0000（T6 空拍）};
  \node[hw state, align=center, anchor=west] (r7) at (5.1,-4.9)   {0x17：0x0000（T7 空拍）};
  % 干线扇出到各行
  \draw[draw=BookMuted, line width=0.62pt] (rom.east) -- (3.6,-2.45);
  \draw[draw=BookMuted, line width=0.62pt] (3.6,0) -- (3.6,-4.9);
  \fill[BookMuted] (3.6,-2.45) circle (0.05);
  \draw[draw=BookMuted, line width=0.62pt] (3.6,0) -- (5.1,0);
  \draw[draw=BookMuted, line width=0.62pt] (3.6,-0.7) -- (5.1,-0.7);
  \draw[draw=BookMuted, line width=0.62pt] (3.6,-1.4) -- (5.1,-1.4);
  \draw[draw=BookMuted, line width=0.62pt] (3.6,-2.1) -- (5.1,-2.1);
  \draw[draw=BookMuted, line width=0.62pt] (3.6,-2.8) -- (5.1,-2.8);
  \draw[draw=BookMuted, line width=0.62pt] (3.6,-3.5) -- (5.1,-3.5);
  \draw[draw=BookMuted, line width=0.62pt] (3.6,-4.2) -- (5.1,-4.2);
  \draw[draw=BookMuted, line width=0.62pt] (3.6,-4.9) -- (5.1,-4.9);
  \fill[BookMuted] (3.6,0) circle (0.05);
  \fill[BookMuted] (3.6,-0.7) circle (0.05);
  \fill[BookMuted] (3.6,-1.4) circle (0.05);
  \fill[BookMuted] (3.6,-2.1) circle (0.05);
  \fill[BookMuted] (3.6,-2.8) circle (0.05);
  \fill[BookMuted] (3.6,-3.5) circle (0.05);
  \fill[BookMuted] (3.6,-4.2) circle (0.05);
  \fill[BookMuted] (3.6,-4.9) circle (0.05);
\end{pdfdiagram}
\diagramnote{ADD（opcode 0x2）在微码 ROM 中的 8 行逐行展开，行地址 = opcode 乘 8 加 T。T0、T1 是 16 个 opcode 完全相同的公共取指（0x4008 = CO+MI、0x1810 = RO+II+CE）；T2 把指令里的地址字段送进 MAR，T3 把 RAM 读出的操作数装进 B，T4 打开 ALU 三态输出，A 收 A+B 并锁存标志；T5--T7 三行控制字全 0，总线空闲、无装载，等回绕到 T0 取下一条。}

有读者会注意到一个看似浪费的结构：取指两行在每个 opcode 下重复存了 16 份，32 行内容两两相同。为什么不把取指单独译出来、让 ROM 只存执行拍？因为省不下什么——单独译取指需要额外的“T0、T1 一律放行”逻辑，等于在 ROM 之外再造一小块硬连线译码，而 ROM 里的行本来就是免费的：128 行用满是本机的规格，空着也不会变快。把全部规则留在一张表里，换来的是核对时只需面对一个对象。这个“宁可在表里重复、也不在表外加逻辑”的取向，与算术与数据通路相关章节对微码控制的评价一致：复杂度进了内容，结构就退到最简。

控制线的根数也在这里交代清楚。16 根不是算出来的理论下限，而是数据通路里每一扇“门”各配一根的结果：八个寄存器或部件的装载使能（MI、II、AI、BI、OI、FI、J、RI），五路总线输出使能（CO、RO、IO、AO、$\Sigma$O），一个 ALU 修饰（SU），一个计数使能（CE），一个停机（HLT）。少一根，某扇门就只能常开或常关；多一根，也只是表里多一位恒 0。将来扩展指令时要做的第一件事，就是回头数这张门清单——加寻址方式可能要加输出使能，加状态寄存器可能要加装载使能，两片 ROM 的 16 位就是这时开始不够用的。

控制线在控制字里的位序是任意定的，但一旦定下就是全机的约定，烧表、接线、排错都按它来。本机的位序如下：

\tablechunkintro{1}{3}{bit15}{bit9}
\begin{longtable}{L{0.08\linewidth}L{0.10\linewidth}L{0.56\linewidth}L{0.10\linewidth}}
\toprule
位 & 控制线 & 有效时的动作 & 所在片 \\
\midrule
bit15 & HLT & 切断时钟源，全机冻结 & 高片 \\\tablerowrule
bit14 & MI & 拍末 MAR ← 总线低 4 位 & 高片 \\\tablerowrule
bit13 & RI & 写窗口内 RAM ← 总线 & 高片 \\\tablerowrule
bit12 & RO & RAM 选中字经 240 上总线 & 高片 \\\tablerowrule
bit11 & II & 拍末 IR ← 总线 & 高片 \\\tablerowrule
bit10 & IO & IR 低 4 位上总线低 4 位 & 高片 \\\tablerowrule
bit9 & AI & 拍末 A ← 总线 & 高片 \\
\bottomrule
\end{longtable}

\tablechunkintro{2}{3}{bit8}{bit2}
\begin{longtable}{L{0.08\linewidth}L{0.10\linewidth}L{0.56\linewidth}L{0.10\linewidth}}
\toprule
位 & 控制线 & 有效时的动作 & 所在片 \\
\midrule
bit8 & AO & A 上总线 & 高片 \\\tablerowrule
bit7 & BI & 拍末 B ← 总线 & 低片 \\\tablerowrule
bit6 & $\Sigma$O & ALU 结果上总线 & 低片 \\\tablerowrule
bit5 & SU & ALU 改做减法（B 取反加一） & 低片 \\\tablerowrule
bit4 & CE & 拍末 PC ← PC + 1 & 低片 \\\tablerowrule
bit3 & CO & PC 上总线低 4 位 & 低片 \\\tablerowrule
bit2 & J & 拍末 PC ← 总线低 4 位 & 低片 \\
\bottomrule
\end{longtable}

\tablechunkintro{3}{3}{bit1}{bit0}
\begin{longtable}{L{0.08\linewidth}L{0.10\linewidth}L{0.56\linewidth}L{0.10\linewidth}}
\toprule
位 & 控制线 & 有效时的动作 & 所在片 \\
\midrule
bit1 & FI & 拍末 Flags ← ALU 的 ZF、CF & 低片 \\\tablerowrule
bit0 & OI & 拍末 OUT ← 总线 & 低片 \\
\bottomrule
\end{longtable}

全机的节拍规则只有一条：所有装载类动作——名字以 I 结尾的控制线，外加 CE、J、FI——都在时钟上升沿同时发生。T 状态计数器在时钟下降沿推进（161 时钟脚经一级反相），此后 ROM 地址随之改变，新控制字经过触发器延迟加 ROM 读取延迟后稳定，到下一个上升沿前半拍就已就位。输出类控制线（CO、RO、IO、AO、$\Sigma$O）是电平有效的整拍驱动，RI 是拍内低半拍的一段写窗口。控制字在拍内稳定、寄存器在上升沿装载——这十个字就是多周期控制器的全部节奏。

\topic{取指的两拍对所有指令相同}
取指为什么必须是公共前奏，想一想 T0 开始时机器知道什么就明白了：IR 里还是上一条指令的残值，PC 指着下一条指令的地址，而“下一条指令是什么”要取回来才知道。在不知道指令是什么的前提下，唯一正确的动作就是把 PC 指的字节取回来——所以 T0、T1 两拍的控制字对全部 16 个 opcode 完全相同，微码表里每个 opcode 的第 0、1 行内容一模一样。这不是浪费，是“取指必须先于译码”这个因果关系的直接抄写。

T0 的控制字是 CO+MI：PC 经三态缓冲驱动总线低 4 位，拍末 MAR 把这 4 位锁存——下一条指令的地址就此交给内存。T1 的控制字是 RO+II+CE：RAM 把 MAR 选中的字（即指令字节）经 240 驱动上总线，拍末同时发生两件互不相干的事——IR 锁存总线值，指令到手；PC 在 CE 作用下加一，为再下一次取指备好地址。CE 与 II 同拍是安全的：PC 计数和 IR 装载共用同一个边沿，但一个改的是 PC 内部状态，一个改的是 IR，互不影响；而此时 PC 的输出缓冲是关的，总线上唯一的驱动者是 RAM，CO 与 RO 从不在同一拍出现。

先把这两拍的数据流画成通路图，谁驱动总线、谁在拍末装载，可以与波形逐拍对照：

\begin{pdfdiagram}
  % 中央总线
  \node[hw memory, align=center, minimum width=8.8cm, minimum height=0.55cm] (bus) at (4.4,0) {8 位总线（取指走低 4 位）};
  % 四角部件
  \node[hw state, align=center] (pc) at (1.2,1.7) {PC};
  \node[hw state, align=center] (mar) at (7.6,1.7) {MAR};
  \node[hw memory, align=center, minimum width=2.4cm] (ram) at (7.6,-1.7) {RAM\\74HC189 两片};
  \node[hw state, align=center] (ir) at (1.2,-1.7) {IR};
  % T0：PC → 总线 → MAR
  \draw[hw bus] (pc.south) -- node[midway, right=1pt, hw label, fill=white, inner sep=1pt]{T0：CO 驱动} ($(bus.north)+(-3.2,0)$);
  \draw[hw bus] ($(bus.north)+(3.2,0)$) -- node[midway, right=1pt, hw label, fill=white, inner sep=1pt]{拍末 MI 锁存} (mar.south);
  % MAR → RAM 地址专线
  \draw[hw bus] (mar.east) -- (9.8,1.7) -- (9.8,-1.7) -- (ram.east);
  \node[hw label, fill=white, inner sep=1pt, rotate=90] at (10.04,0) {地址专线};
  % T1：RAM → 总线 → IR
  \draw[hw bus] (ram.north) -- node[midway, right=1pt, hw label, fill=white, inner sep=1pt]{T1：RO 驱动} ($(bus.south)+(3.2,0)$);
  \draw[hw bus] ($(bus.south)+(-3.2,0)$) -- node[midway, right=1pt, hw label, fill=white, inner sep=1pt]{拍末 II 锁存} (ir.north);
  % CE：PC 内部加一（不占总线）
  \draw[hw return] (0.0,2.5) -- (1.2,2.5) -- (pc.north);
  \node[hw label, fill=white, inner sep=1pt] at (0.6,2.76) {CE：同拍 PC 内部加一};
\end{pdfdiagram}
\diagramnote{取指两拍的数据流（对照本节波形图）。T0：CO 打开 PC 的三态缓冲，地址经总线低 4 位送出，拍末 MI 把它锁进 MAR——此后地址走 MAR 到 RAM 的点对点专线，不再依赖总线；T1：RO 让 RAM 把 MAR 选中的指令字节经 240 驱动上总线，拍末 II 锁进 IR，同拍 CE 让 PC 内部加一。总线一拍只服务一对收发，所以取指至少两拍；两拍的驱动者分别是 CO 与 RO，从不在同一拍出现。}

把这两拍画成波形，沿与电平的关系一眼可见。

\begin{waveform}[x=0.86cm,y=0.9cm]
  % T 状态翻转的下降沿参考虚线（x=2,6,10）
  \draw[hw return] (2,0.2) -- (2,9.3);
  \draw[hw return] (6,0.2) -- (6,9.3);
  \draw[hw return] (10,0.2) -- (10,9.3);
  % CLK：下降沿在 2/6/10（拍界，T 推进），上升沿在 4/8（装载）
  \draw[BookAccent, line width=0.9pt] (0,9.0) -- (2,9.0) -- (2,8.2) -- (4,8.2) -- (4,9.0) -- (6,9.0) -- (6,8.2) -- (8,8.2) -- (8,9.0) -- (10,9.0) -- (10,8.2) -- (10.5,8.2);
  \node[hw label, anchor=east] at (-0.2,8.6) {CLK};
  % T 状态
  \draw[BookTeal, line width=0.9pt] (0,6.9) rectangle (2,7.7);
  \node at (1,7.3) {复位};
  \draw[BookTeal, line width=0.9pt] (2,6.9) rectangle (6,7.7);
  \node at (4,7.3) {T0};
  \draw[BookTeal, line width=0.9pt] (6,6.9) rectangle (10,7.7);
  \node at (8,7.3) {T1};
  \node[hw label, anchor=east] at (-0.2,7.3) {T 状态};
  % CO / MI（T0 整拍有效）
  \draw[BookAccent, line width=0.9pt] (0,5.8) -- (2,5.8) -- (2,6.6) -- (6,6.6) -- (6,5.8) -- (10.5,5.8);
  \node[hw label, anchor=east] at (-0.2,6.2) {CO};
  \draw[BookAccent, line width=0.9pt] (0,4.8) -- (2,4.8) -- (2,5.6) -- (6,5.6) -- (6,4.8) -- (10.5,4.8);
  \node[hw label, anchor=east] at (-0.2,5.2) {MI};
  % RO / II / CE（T1 整拍有效）
  \draw[BookAccent, line width=0.9pt] (0,3.8) -- (6,3.8) -- (6,4.6) -- (10,4.6) -- (10,3.8) -- (10.5,3.8);
  \node[hw label, anchor=east] at (-0.2,4.2) {RO};
  \draw[BookAccent, line width=0.9pt] (0,2.8) -- (6,2.8) -- (6,3.6) -- (10,3.6) -- (10,2.8) -- (10.5,2.8);
  \node[hw label, anchor=east] at (-0.2,3.2) {II};
  \draw[BookAccent, line width=0.9pt] (0,1.8) -- (6,1.8) -- (6,2.6) -- (10,2.6) -- (10,1.8) -- (10.5,1.8);
  \node[hw label, anchor=east] at (-0.2,2.2) {CE};
  % 总线值
  \draw[BookTeal, line width=0.9pt] (0,0.4) rectangle (2,1.2);
  \node at (1,0.8) {0x00};
  \draw[BookTeal, line width=0.9pt] (2,0.4) rectangle (6,1.2);
  \node at (4,0.8) {0x00（PC）};
  \draw[BookTeal, line width=0.9pt] (6,0.4) rectangle (10,1.2);
  \node at (8,0.8) {0x1F（RAM[0]）};
  \node[hw label, anchor=east] at (-0.2,0.8) {总线};
  % 装载动作（拍内上升沿）
  \node[hw label, anchor=north] at (4,-0.1) {装载沿 MAR←0x00};
  \node[hw label, anchor=north] at (8,-0.85) {装载沿 IR←0x1F，PC←1};
\end{waveform}
\diagramnote{取指两拍（复位后的第一条指令，对照“编程、调试与扩展”一节程序一：RAM[0]=0x1F）。T0：CO 让 PC 驱动总线，装载沿（上升沿）MI 把地址 0x00 锁进 MAR；T1：RO 让 RAM[0] 上总线，装载沿 II 把它锁进 IR、CE 让 PC 加一。三根参考虚线是 T 状态翻转的下降沿：T 计数器在下降沿推进，新控制字到装载沿前半拍已经稳定，所有装载都在拍内上升沿发生。这两拍对全部 16 个 opcode 完全相同。}

\begin{longtable}{L{0.08\linewidth}L{0.14\linewidth}L{0.20\linewidth}L{0.14\linewidth}L{0.30\linewidth}}
\toprule
T 拍 & 控制字 & 有效控制线 & 总线驱动者 & 拍末动作 \\
\midrule
T0 & \code{0x4008} & CO、MI & PC & MAR ← PC \\\tablerowrule
T1 & \code{0x1810} & RO、II、CE & RAM & IR ← RAM[MAR]，PC ← PC+1 \\
\bottomrule
\end{longtable}

控制字读不顺的时候，把它拆回位再数一遍是最快的自检。\code{0x4008} 展开是 0100 0000 0000 1000：bit14 为 1 是 MI，bit3 为 1 是 CO，合起来 CO+MI；\code{0x1810} 展开是 0001 1000 0001 0000：bit12 为 RO、bit11 为 II、bit4 为 CE。任何一行控制字都经得起这样拆——拆出来的控制线集合与微操作表不一致时，错的要么是位序表，要么是你的加法，而位序表只有一张，加法可以重算。

T1 拍末是个值得停下来看一秒的时刻：IR 刚换上新指令，ROM 地址高 4 位跟着换新，从 T2 起控制字开始按指令分化——取指的公共性到此为止。还要说一句执行拍用不满的指令：本机所有指令固定走满 T0 到 T7 八拍，T4 或 T2 就做完事的指令，余下的拍控制字全 0，总线空闲、无装载，机器空转到 T7 再回 T0。固定八拍换来的是控制单元里没有任何“提前结束”逻辑——想省这些空拍，需要加一套按 opcode 让 T 计数器提前回零的预置电路，电路不难，但每多一种拍数，控制字表就多一列要核对的例外。本机选择不省，把简单留给核对。

两个相关的问题一并回答。其一，取指为什么不能并成一拍：T0 做的是“PC 经总线进 MAR”，T1 做的是“RAM 经总线进 IR”，两件事都依赖总线，而总线一拍只能服务一对收发——把两拍并成一拍，等于要求总线在同一拍先送地址、再送指令，这是两个总线周期，不是一个。单端口内存加单总线的结构，注定取指至少两拍。其二，复位之后机器的第一动作长什么样：复位把 PC 清为 0、T 计数器清为 T0、IR 清为 \code{0x00}，ROM 地址是 opcode 0 的第 0 行，控制字正是 CO+MI——机器从地址 0 的取指开始，而 IR 里的 NOP 残值保证了即使有什么使能残存，执行的也只是空操作。设计的每个已知初值都指向同一个出发线。

取指是两拍里最容易核对的一段，建议第一次上电就把它核熟。单步时钟，T0 时看总线 LED 应显示 PC 的值（第一次是 \code{0x0}），拍末 MAR 的四个输出脚应变成同一个值（用表笔或给 MAR 也挂一组 LED 都行，调试期值得挂）；T1 时总线应显示 RAM 该单元的内容，拍末 IR 的高 4 位 LED 应变成该字节的 opcode。三个量、两个沿，全对，说明 PC、MAR、RAM、IR、T 计数器、ROM 取指行这一整条链是通的；任何一环不对，都停在这两拍里排，不要往后走。

\topic{11 条指令的逐条微操作表}
下面这张全表是本章的心脏。读表约定：每格列出该拍有效的控制线，“—”表示该拍控制字为 \code{0x0000}，无动作；所有装载类动作在该拍末尾的上升沿完成。JC、JZ 两行里的 J 带星号，含义随后说明。

\tablechunkintro{1}{2}{NOP}{JMP}
\begin{longtable}{L{0.12\linewidth}L{0.22\linewidth}L{0.22\linewidth}L{0.26\linewidth}L{0.06\linewidth}}
\toprule
指令 & T2 & T3 & T4 & T5–T7 \\
\midrule
NOP & — & — & — & — \\\tablerowrule
LDA & IO+MI & RO+AI & — & — \\\tablerowrule
ADD & IO+MI & RO+BI & $\Sigma$O+AI+FI & — \\\tablerowrule
SUB & IO+MI & RO+BI & SU+$\Sigma$O+AI+FI & — \\\tablerowrule
STA & IO+MI & AO+RI & — & — \\\tablerowrule
LDI & IO+AI & — & — & — \\\tablerowrule
JMP & IO+J & — & — & — \\
\bottomrule
\end{longtable}

\tablechunkintro{2}{2}{JC}{HLT}
\begin{longtable}{L{0.12\linewidth}L{0.22\linewidth}L{0.22\linewidth}L{0.26\linewidth}L{0.06\linewidth}}
\toprule
指令 & T2 & T3 & T4 & T5–T7 \\
\midrule
JC & IO+J* & — & — & — \\\tablerowrule
JZ & IO+J* & — & — & — \\\tablerowrule
OUT & AO+OI & — & — & — \\\tablerowrule
HLT & HLT & — & — & — \\
\bottomrule
\end{longtable}

逐条看一遍这张表在说什么。NOP 取指后空转六拍，唯一的作用是占时间。LDA 用 T2 把指令里的地址字段送进 MAR，T3 把该地址的字读进 A。ADD 比 LDA 多一拍：操作数先进 B（T3），T4 才打开 ALU 的三态输出，把 A+B 送回 A，同拍 FI 把新的 ZF、CF 记进 Flags——记住标志的办法是在它还新鲜的那拍就把它锁存。SUB 与 ADD 同构，T4 多一根 SU：B 经 74HC86 取反、进位输入置 1，加法器算出的就是 A−B。STA 是唯一写内存的指令，T3 的 AO+RI 让 A 驱动总线、RAM 在低半拍收数。LDI 是最快的装载，IR 低 4 位直接进 A，连 MAR 都不惊动。三条跳转里，JMP 无条件把指令字段装进 PC；JC、JZ 外形与 JMP 相同，差别就在那个星号上。OUT 把 A 抄给显示。HLT 在 T2 拉起 HLT 线，“总体架构：总线、时钟与复位”一节的时钟与门随即关闭，T 计数器冻结，全机停在原地——退出停机的唯一办法是复位：复位清 T 计数器和 IR，HLT 线随之撤销，时钟恢复。

把每条指令的执行要领浓缩成一栏，与上表对照着记：

\tablechunkintro{1}{2}{NOP}{JMP}
\begin{longtable}{L{0.12\linewidth}L{0.72\linewidth}}
\toprule
指令 & 执行期要领 \\
\midrule
NOP & 取指之后空转六拍，只占时间 \\\tablerowrule
LDA & 先送地址（T2），再读数进 A（T3） \\\tablerowrule
ADD & 操作数先进 B（T3），再整加写 A 并记标志（T4） \\\tablerowrule
SUB & 同 ADD，T4 加 SU：B 取反加一即减 \\\tablerowrule
STA & 唯一写内存：AO 与 RI 同拍（T3） \\\tablerowrule
LDI & 最快的装载：IR 低 4 位直接进 A（T2） \\\tablerowrule
JMP & 无条件：IR 低 4 位进 PC（T2） \\
\bottomrule
\end{longtable}

\tablechunkintro{2}{2}{JC}{HLT}
\begin{longtable}{L{0.12\linewidth}L{0.72\linewidth}}
\toprule
指令 & 执行期要领 \\
\midrule
JC & 同 JMP 的外形，CF 为 1 才装载 \\\tablerowrule
JZ & 同 JMP 的外形，ZF 为 1 才装载 \\\tablerowrule
OUT & 把 A 抄给显示，总线只送不还（T2） \\\tablerowrule
HLT & 关掉时钟源，全机冻结，唯复位可退出（T2） \\
\bottomrule
\end{longtable}

星号的含义是条件装载。微码 ROM 在 JC、JZ 的 T2 行照常输出 J 请求，但这根线到 PC 的装载端之前要过一级标志选通：仅当对应标志成立时才放行。用 IR 高 4 位译出 c7（当前指令是 JC）和 c8（当前指令是 JZ）两条线，选通逻辑如下：

\begin{lstlisting}
c7 = (opcode == 0111)        ; IR7..IR4 = 0x7，由 74HC21 译出
c8 = (opcode == 1000)        ; IR7..IR4 = 0x8，同上
J_PC = J AND ( (c7 AND CF) OR (c8 AND ZF) OR NOT(c7 OR c8) )
\end{lstlisting}

逐项检查这个式子：JMP（0x6）时 c7、c8 都为 0，第三项为 1，J 原样放行，无条件跳；JC 时 c7 为 1，式子退化成 J AND CF，进位标志说了算；JZ 同理退化成 J AND ZF；其它指令的 T2 行里 J 本来就是 0，选通无影响。器件账：反相器用 74HC04，两个四输入译码用一片 74HC21，尾部的与或用 74HC08 和 74HC32 各一片，合计四片，全部接在 ROM 的 J 输出脚与 PC 装载端之间。选通是纯组合逻辑，控制字整拍稳定，标志早在前面的 ADD 或 SUB 里就锁存好了，拍末边沿采样到的 J 是毫不含糊的电平。

把这片选通网络画成门级图，四片 74HC 的分工如下。

\begin{pdfdiagram}[x=0.92cm,y=1cm]
  \node[hw state, minimum width=1.8cm] (ir) at (0.8,4.4) {IR 高 4 位\\IR7..IR4};
  \node[hw compute, minimum width=3.0cm] (dec) at (4.0,4.4) {74HC04 → 74HC21\\c7：0111（JC）\\c8：1000（JZ）};
  \node[hw compute, minimum width=2.1cm] (a1) at (7.3,6.0) {74HC08\\c7·CF};
  \node[hw compute, minimum width=2.1cm] (n1) at (7.3,4.4) {74HC32 + 04\\$\overline{c7+c8}$};
  \node[hw compute, minimum width=2.1cm] (a2) at (7.3,2.8) {74HC08\\c8·ZF};
  \node[hw compute, minimum width=1.9cm] (or3) at (10.3,4.4) {74HC32\\三路或};
  \node[hw compute, minimum width=1.9cm] (fan) at (12.9,4.4) {74HC08\\J·K};

  \draw[hw bus] (ir) -- (dec);
  \draw[hw return] ($(dec.east)+(0,0.35)$) -- (5.8,4.75) -- (5.8,6.25) -- (a1.west);
  \draw[hw return] (5.8,4.75) -- (n1.west);
  \draw[hw return] ($(dec.east)+(0,-0.35)$) -- (6.1,4.05) -- (6.1,2.55) -- (a2.west);
  \draw[hw return] (6.1,4.05) -- (n1.west);

  \node[hw label, anchor=east] at (0.7,6.0) {CF};
  \draw[hw return] (0.8,6.0) -- (a1.west);
  \node[hw label, anchor=east] at (0.7,2.8) {ZF};
  \draw[hw return] (0.8,2.8) -- (a2.west);

  \draw[hw return] (a1.east) -- (9.0,6.0) -- (9.0,4.75) -- (or3.west);
  \draw[hw return] (n1.east) -- (or3.west);
  \draw[hw return] (a2.east) -- (9.0,2.8) -- (9.0,4.05) -- (or3.west);
  \draw[hw return] (or3) -- (fan);

  \node[hw label] at (12.9,1.6) {J（ROM 输出）};
  \draw[hw return] (12.9,2.0) -- (fan.south);
  \draw[hw return] (fan.east) -- (14.4,4.4);
  \node[hw label, anchor=west] at (14.55,4.4) {J\_PC};
\end{pdfdiagram}
\diagramnote{J 的条件选通逻辑。IR 高 4 位经 74HC04 按位取反、74HC21 四输入译码，分出 c7（opcode=0111，JC）与 c8（opcode=1000，JZ）两条指令线；74HC08 把 c7、c8 分别与 CF、ZF 相与，74HC32 加 74HC04 给出第三项（c7、c8 都为 0，对应 JMP）；三路经 74HC32 相或后再与 J 相与，得到送进 PC 装载端的 J\_PC。JMP 走第三项无条件放行，JC、JZ 由各自标志裁决，其余指令的 T2 行 J 本就为 0，选通无影响。}

条件不成立时的“顺序执行”值得拆开看一眼，因为它其实是免费的：JC 的 T1 里 CE 已经让 PC 加过一，PC 指向的就是下一条指令；T2 的 J 被标志拦下，PC 保持不动，回绕后的 T0 照常把这个已加一的值送去 MAR——不跳转不需要任何专门动作，它只是“装载没有发生”。跳转目标的语义也借这里定死：装进 PC 的是指令低 4 位，一个绝对地址，所以本机的跳转只能落在 0 到 15 这 16 个字节里，没有相对跳转、没有偏移量——程序多大，跳转范围就多大，这就是 4 位 PC 的世界观。

最后交代空闲拍的总线。T5 到 T7 的控制字全 0，没有任何驱动者占用总线，总线电平由“总体架构：总线、时钟与复位”一节的总线下拉电阻给出确定值——读作 \code{0x00}。这个约定不是装饰：没有它，空闲拍浮空的总线会在 0 与 1 之间漂，虽然无人装载、逻辑无害，但总线 LED 会无规律地闪烁，单步调试时平白添一层噪声。让所有无人使用的观测量都有确定读数，是这台机器“每一步可观察”的最后一块拼图。

把微操作表逐列翻译成 16 位控制字，再按“地址 = opcode 乘 8 加 T”排列，就是要烧进两片 EEPROM 的全部内容。128 行里非零的只有下面这些——外加每个 opcode 都相同的两行取指——其余一律 \code{0x0000}，包括未使用的 opcode \code{0x9} 到 \code{0xD}：它们的执行拍全空，行为与 NOP 无异，程序里误用了不会破坏状态，但也不会有任何作为。

\tablechunkintro{1}{3}{\code{op*8+0}}{\code{0x14}}
\begin{longtable}{L{0.14\linewidth}L{0.14\linewidth}L{0.56\linewidth}}
\toprule
ROM 地址 & 控制字 & 含义 \\
\midrule
\code{op*8+0} & \code{0x4008} & T0 取指：CO+MI，16 个 opcode 相同 \\\tablerowrule
\code{op*8+1} & \code{0x1810} & T1 取指：RO+II+CE，16 个 opcode 相同 \\\tablerowrule
\code{0x0A} & \code{0x4400} & LDA T2：IO+MI \\\tablerowrule
\code{0x0B} & \code{0x1200} & LDA T3：RO+AI \\\tablerowrule
\code{0x12} & \code{0x4400} & ADD T2：IO+MI \\\tablerowrule
\code{0x13} & \code{0x1080} & ADD T3：RO+BI \\\tablerowrule
\code{0x14} & \code{0x0242} & ADD T4：$\Sigma$O+AI+FI \\
\bottomrule
\end{longtable}

\tablechunkintro{2}{3}{\code{0x1A}}{\code{0x32}}
\begin{longtable}{L{0.14\linewidth}L{0.14\linewidth}L{0.56\linewidth}}
\toprule
ROM 地址 & 控制字 & 含义 \\
\midrule
\code{0x1A} & \code{0x4400} & SUB T2：IO+MI \\\tablerowrule
\code{0x1B} & \code{0x1080} & SUB T3：RO+BI \\\tablerowrule
\code{0x1C} & \code{0x0262} & SUB T4：SU+$\Sigma$O+AI+FI \\\tablerowrule
\code{0x22} & \code{0x4400} & STA T2：IO+MI \\\tablerowrule
\code{0x23} & \code{0x2100} & STA T3：AO+RI \\\tablerowrule
\code{0x2A} & \code{0x0600} & LDI T2：IO+AI \\\tablerowrule
\code{0x32} & \code{0x0404} & JMP T2：IO+J \\
\bottomrule
\end{longtable}

\tablechunkintro{3}{3}{\code{0x3A}}{\code{0x7A}}
\begin{longtable}{L{0.14\linewidth}L{0.14\linewidth}L{0.56\linewidth}}
\toprule
ROM 地址 & 控制字 & 含义 \\
\midrule
\code{0x3A} & \code{0x0404} & JC T2：IO+J，J 经 CF 选通 \\\tablerowrule
\code{0x42} & \code{0x0404} & JZ T2：IO+J，J 经 ZF 选通 \\\tablerowrule
\code{0x72} & \code{0x0101} & OUT T2：AO+OI \\\tablerowrule
\code{0x7A} & \code{0x8000} & HLT T2：HLT \\
\bottomrule
\end{longtable}

两表的一致性要逐列核对，方法很机械：从微操作表的每一格出发，按位序表把控制线翻成位、加出控制字，再到 ROM 表里按 opcode 乘 8 加 T 找到那一行，两者必须相同。以 ADD 的 T4 为例：$\Sigma$O 是 bit6、AI 是 bit9、FI 是 bit1，控制字为 \code{0x0040+0x0200+0x0002 = 0x0242}；ADD 的 opcode 是 \code{0x2}，地址为 \code{0x2*8+4 = 0x14}——ROM 表第 \code{0x14} 行正是 \code{0x0242}。其余 17 个非零行照此全部核一遍，这是烧片前不可省略的一步，错一位就是一条指令的隐疾。

最后把 ADD 的完整六拍走一遍，作为两张表的联合读法（设 A=\code{0x06}，RAM[\code{0x4}]=\code{0x07}，该指令存在地址 \code{0x1}；清单是纯文本，把控制线 $\Sigma$O 写作 SO）：

\begin{lstlisting}
// 注释（推进）：每个 u/T 行代表一个控制节拍，同一行内的控制信号并行生效。
// 注释（边界）：next 决定下一微地址；等待或错误时不得提前进入写回。
T0  CO+MI        总线 <- PC(0x1)          拍末 MAR <- 0x1
T1  RO+II+CE     总线 <- RAM[0x1]=0x24    拍末 IR <- 0x24, PC <- 0x2
T2  IO+MI        总线 <- IR低4位=0x4      拍末 MAR <- 0x4
T3  RO+BI        总线 <- RAM[0x4]=0x07    拍末 B <- 0x07
T4  SO+AI+FI     总线 <- A+B=0x0D         拍末 A <- 0x0D, Flags <- {ZF=0,CF=0}
T5..T7  --       总线空闲，无装载，等回绕到 T0 取下一条
\end{lstlisting}

\topic{微码 ROM 的两种实现与门阵列对照}
EEPROM 路线的全部制造工作就是上面两张表：算好 128 个 16 位控制字，高字节烧进高片、低字节烧进低片，插上板子，控制单元就完工。它的实现特性是“规则写在内容里”——改指令行为就是改表重烧：想让 LDA 顺便清标志，把 \code{0x0B} 行加上 FI 位即可；想添一条新指令，把空着的 opcode \code{0x9} 的 8 行填上即可，面包板上一根线不动。烧写用通用编程器最省事；按“内存与输出模块”一节编程模式的思路用开关加按钮自编程也行，EEPROM 在 5 V 下的字节写入时序与 RAM 相仿，只是每字节要留足器件手册要求的写入时间，典型是几毫秒。手工填 128 行容易出错，用一段几行的脚本生成更稳妥：

\begin{lstlisting}
# 注释（推进）：每个 u/T 行代表一个控制节拍，同一行内的控制信号并行生效。
# 注释（边界）：next 决定下一微地址；等待或错误时不得提前进入写回。
W = [0x0000] * 128
for op in range(16):
    W[op*8+0] = 0x4008     # T0: CO+MI（所有指令相同）
    W[op*8+1] = 0x1810     # T1: RO+II+CE
def put(op, t, word): W[op*8+t] = word
put(0x1,2,0x4400); put(0x1,3,0x1200)                 # LDA
put(0x2,2,0x4400); put(0x2,3,0x1080); put(0x2,4,0x0242)  # ADD
put(0x3,2,0x4400); put(0x3,3,0x1080); put(0x3,4,0x0262)  # SUB
put(0x4,2,0x4400); put(0x4,3,0x2100)                 # STA
put(0x5,2,0x0600)                                    # LDI
put(0x6,2,0x0404)                                    # JMP
put(0x7,2,0x0404); put(0x8,2,0x0404)                 # JC/JZ（J 经标志选通）
put(0xE,2,0x0101)                                    # OUT
put(0xF,2,0x8000)                                    # HLT
# 高字节写高片，低字节写低片；未置的行保持 0x0000
\end{lstlisting}

烧完不等于烧对，上板之前先把 ROM 内容读回来核一遍：编程器都有读出校验功能，把两片各 128 行读出，与源表逐字节比对，全对再插。跳过这一步直接上板的后果，是把“烧写错误”混进“接线错误”里一起排——两类故障在面包板上的症状一模一样，而区分它们只需要一次读回。养成“凡可写必先读回”的习惯，在“内存与输出模块”一节的 DIP 开关和这里的 EEPROM 上是同一条纪律。把从脚本到上板的流程画成一条链，读回校验在链条里的位置就清楚了。

\begin{pdfdiagram}
  \node[hw compute, align=center, minimum width=2.5cm] (gen) at (0,0) {脚本生成\\128 个 16 位控制字};
  \node[hw state, align=center, minimum width=2.5cm] (img) at (3.3,0) {二进制映像\\高 / 低字节分开};
  \node[hw memory, align=center, minimum width=2.5cm] (burn) at (6.6,0) {编程器烧写\\EEPROM 高片 + 低片};
  \node[hw state, align=center, minimum width=2.2cm] (board) at (9.7,0) {插上板子};
  \node[hw compute, align=center, minimum width=2.8cm] (ver) at (6.6,-1.9) {读回两片各 128 行\\与源表逐字节比对};
  \draw[hw bus] (gen) -- node[midway, above=1pt, hw label, fill=white, inner sep=1pt]{按表填字} (img);
  \draw[hw bus] (img) -- node[midway, above=1pt, hw label, fill=white, inner sep=1pt]{高字节进高片} (burn);
  \draw[hw return] (burn.south) -- node[midway, right=1pt, hw label]{烧完先读回} (ver.north);
  \draw[hw bus] (ver.east) -- (9.7,-1.9) -- node[pos=0.7, right=1pt, hw label]{全对再上板} (board.south);
  \draw[hw return] (ver.west) -- node[pos=0.4, below=1pt, hw label]{有错改表重烧} (3.3,-1.9) -- (img.south);
\end{pdfdiagram}
\diagramnote{微码 ROM 的烧写流程。脚本按微操作表生成 128 个 16 位控制字，高字节烧进高片、低字节烧进低片；烧完先用编程器把两片各 128 行读回，与源表逐字节比对，全对再上板。跳过读回直接上板，会把烧写错误混进接线错误里一起排查——两类故障症状一样，区分它们只需要一次读回。}

条件跳转在 EEPROM 路线上还有第二种做法，值得一提：把 CF、ZF 直接接成 ROM 的第 8、9 根地址线，表从 128 行扩成 512 行，JC、JZ 的 T2 行按标志的四种组合各写一份——成立的位置写 \code{0x0404}，不成立的位置写 \code{0x0000}。这样 J 不再需要片外选通，四片门省掉了，代价是表大四倍、烧写时要多一层小心。这是“存储空间换组合逻辑”的典型交易，28C64 的容量完全付得起。

门阵列路线把同一张表用组合逻辑“算”出来：一片 74HC154（4 线到 16 线译码器）把 opcode 译成 16 条指令线，一片 74HC138（3 线到 8 线译码器）把 T 译成 8 条拍线，每根控制线就是若干“指令线乘拍线”的或。两片的输出都是低有效，实际搭的时候用与非的视角更顺手；下面的式子按高有效写出逻辑关系：

\begin{lstlisting}
MI = T0 + T2*(LDA+ADD+SUB+STA)
RO = T1 + T3*(LDA+ADD+SUB)
II = T1
CE = T1
CO = T0
IO = T2*(LDA+ADD+SUB+STA+LDI+JMP+JC+JZ)
AI = T2*LDI + T3*LDA + T4*(ADD+SUB)
BI = T3*(ADD+SUB)
SO = T4*(ADD+SUB)        ; SO 即正文所说的加法器结果输出使能
SU = T4*SUB
FI = T4*(ADD+SUB)
AO = T2*OUT + T3*STA
RI = T3*STA
J  = T2*(JMP + JC*CF + JZ*ZF)
OI = T2*OUT
HLT = T2*opF             ; opF 即 HLT 指令的译码线
\end{lstlisting}

注意 J 那一行：条件跳转的标志选通在门阵列里被天然吸收进译码式，不再需要 EEPROM 方案那级片外选通——这是门阵列少有的扳回一分的地方。其余各行读法相同：每根控制线追溯下去，不过是几条指令线、几条拍线、最多两根标志的组合。代价在器件数量和可改性：译码两片加与或非门约十片，总共十二片上下，占用整整一块面包板；更实质的代价是改一条指令的行为就要动接线，改上三五次之后，那团线的正确性就只剩当事人敢打包票。

\begin{longtable}{L{0.22\linewidth}L{0.34\linewidth}L{0.34\linewidth}}
\toprule
维度 & EEPROM 微码 & 74HC 门阵列 \\
\midrule
规则写在哪里 & ROM 的内容，即数据 & 门的接线，即结构 \\\tablerowrule
改一条指令 & 重烧几行，硬件不动 & 改接线 \\\tablerowrule
器件数 & 两片，加选通门四片 & 十二片上下 \\\tablerowrule
控制字延迟 & 触发器加 ROM 读取，约两百纳秒 & 触发器加几级门，几十纳秒 \\\tablerowrule
扩展条件跳转 & 加选通门，或把标志接成地址线扩表 & 直接写进译码式 \\\tablerowrule
调试方式 & 把 ROM 内容读回，与表逐行比对 & 逐级量门输出，对照译码式 \\\tablerowrule
适合阶段 & 指令集还在调整期 & 指令集冻结、想省掉 ROM 时 \\
\bottomrule
\end{longtable}

两条路线与算术与数据通路相关章节、处理器核心相关章节的对照是同一个故事的实物版。算术与数据通路相关章节比较过硬连线、状态机与微码三种控制器组织方式，处理器核心相关章节在最小 CPU 上演示过同一张状态表的两种实现：一边是“状态触发器加次态逻辑加输出译码”，一边是“每行一条微指令的控制存储器”。本节的 EEPROM 与门阵列正是那一组取舍落在面包板上的样子——两边每一拍输出的控制线必须一根不差，差别只在规则存放在内容里还是接线里。务实的建议是用 EEPROM 把机器先跑通，表错了重烧就是；等指令集稳定，再把门阵列版当作一次“把存储器译码展开成逻辑”的练习，两边对照着量同一根控制线，是理解那两章取舍表最直接的方式。速度差异在这里不成议题：两百纳秒对几十纳秒，在十赫兹量级的教学时钟面前都是四个数量级的余量。

两条路线各有自己的高发故障，排法也不同。EEPROM 路线的高发故障是位序类：高低两片插反（控制字高低字节对调，症状是全机行为荒诞但隐约像样）、位序接反（每片内部 D0 到 D7 颠倒，症状是某几类控制线集体错位）、地址线序错（opcode 与 T 的拼接顺序颠倒，取指能对、执行全错）。这类故障的共同点是“错得整齐”，把 ROM 内容读回与表比对一次即现形。门阵列路线的高发故障是接线类：某根拍线漏接（对应拍所有动作消失）、154 与 138 的低有效输出被当成高有效用（所有动作集体反相）、某级与门输入错线（单根控制线行为异常）。排法是顺一根控制线的译码式逐级量，式子在手，最多三四级就到底。

收尾仍是检查清单，按序做：

\begin{longtable}{L{0.20\linewidth}L{0.38\linewidth}L{0.30\linewidth}}
\toprule
检查点 & 方法 & 预期结果 \\
\midrule
复位状态 & 按复位后量 T 计数器与 IR & T=0，IR=\code{0x00}；微码 ROM 组合输出第 0 行 \code{0x4008}，仅 CO、MI 有效，PC 驱动总线 \\\tablerowrule
取指两拍 & 单步任意指令的 T0、T1 & T0 为 CO+MI，T1 为 RO+II+CE \\\tablerowrule
opcode 专线 & 单步到 T2，量 ROM 地址高 4 位 & 等于刚取回字节的 opcode \\\tablerowrule
回绕干净 & 连续单步 16 拍看 Q2–Q0 & T0 到 T7 循环两遍，无跳态 \\\tablerowrule
地址拼接 & 装入 ADD 后单步到 T4，量 ROM 地址脚 & 地址 \code{0x14}，控制字 \code{0x0242} \\\tablerowrule
互斥边界 & 逐拍列出总线驱动者 & 任何一拍只有一个驱动者 \\\tablerowrule
条件跳转 & CF 为 0、1 各执行一次 JC & 前者顺序执行，后者跳到目标 \\\tablerowrule
空闲拍 & 单步到 T5–T7，看总线 LED & 总线读 \code{0x00}，无驱动者 \\\tablerowrule
停机与退出 & 执行 HLT 后量时钟，再按复位 & 时钟先冻结，复位后从 T0 重跑 \\\tablerowrule
两表一致 & 微操作表逐格译成控制字 & 与 ROM 表逐行相同 \\
\bottomrule
\end{longtable}

至此控制单元也立起来了：一张表、两块状态、十六根线，机器从此会自己按拍发号施令。“编程、调试与扩展”一节把前两节的部件和这一节的时序合在一起，用两个完整程序做一次从头到尾的演练。

\section{五、编程、调试与扩展}

\topic{程序一：累加循环的完整机器码与逐拍追踪（trace）}程序一的任务是从 1 加到 4，结果为 10。
一个累加循环需要三样东西：累加和、计数器、结束判断。
先做一个容易忽视的决定：累加和 sum 与计数器 i 都必须住在 RAM 里，而不能住在 A 里。
原因是循环中途要用 A 做比较——把 i 减去上限——A 的旧值会被覆盖，凡是要跨指令保存的量都必须写回 RAM。
这是单累加器机器的第一课：A 是工作台，RAM 才是仓库。
第二个决定来自 16 字节预算的约束。
若把初始化也写进程序（两条 LDI 加两条 STA 共 4 字节），再加常数 1、上限 4、i、sum 四个数据字节，总共 20 字节，装不下。
解法是把初值交给编程模式：“内存与输出模块”一节讲过，拨码可以把每个字节预先写好，i 拨成 1、sum 拨成 0，程序从地址 0 起只保留循环体本身。
这也解释了为什么程序一里一条 LDI 都用不上：LDI 是为运行中赋初值准备的，而这里的初值在上电前就已经在 RAM 里了。
程序与数据同处一块 RAM，改一个数据字节就改变了程序行为——本章习题会让读者亲手验证这一点。
布局上沿用本章约定：程序从地址 0 向上长，数据从地址 15 向下长，中间是空闲带；程序一旦超长就会撞上数据，这正是真实系统里代码段与数据段的最小模型。
循环结构采用“先加、再显示、后判断”：先把 i 加进 sum，用 OUT 显示部分和；再把 i 与上限 4 比较。
比较不需要专门的指令：SUB 13 计算 i 减 4，结果为零时 ALU 把 ZF 置 1，紧跟的 JZ 11 读到这个标志就跳到 HLT。
判断没通过时才真正计数：i 加 1、写回、JMP 0 回循环头。
跳转目标 11 是数着编码表数出来的：数据区从 12 开始，HLT 放在 11。
完整编码如下，高 4 位 opcode、低 4 位地址或立即数，与章首的 ISA 表一一对应。

\tablechunkintro{1}{3}{\code{0}}{\code{6}}
\begin{longtable}{L{0.10\linewidth}L{0.12\linewidth}L{0.16\linewidth}L{0.50\linewidth}}
\toprule
地址 & 内容 & 助记符 & 注释 \\
\midrule
\code{0} & \code{0x1F} & \code{LDA 15} & A ← sum；循环从这里开始 \\\tablerowrule
\code{1} & \code{0x2E} & \code{ADD 14} & A ← sum + i \\\tablerowrule
\code{2} & \code{0x4F} & \code{STA 15} & sum ← sum + i，写回 RAM \\\tablerowrule
\code{3} & \code{0xE0} & \code{OUT} & 显示当前部分和 \\\tablerowrule
\code{4} & \code{0x1E} & \code{LDA 14} & A ← i \\\tablerowrule
\code{5} & \code{0x3D} & \code{SUB 13} & A ← i − 4；i 到顶时结果为 0，ZF=1 \\\tablerowrule
\code{6} & \code{0x8B} & \code{JZ 11} & ZF=1 即跳到 11 号停机 \\
\bottomrule
\end{longtable}

\tablechunkintro{2}{3}{\code{7}}{\code{13}}
\begin{longtable}{L{0.10\linewidth}L{0.12\linewidth}L{0.16\linewidth}L{0.50\linewidth}}
\toprule
地址 & 内容 & 助记符 & 注释 \\
\midrule
\code{7} & \code{0x1E} & \code{LDA 14} & 没跳走：A ← i（刚才的差已无用了） \\\tablerowrule
\code{8} & \code{0x2C} & \code{ADD 12} & A ← i + 1 \\\tablerowrule
\code{9} & \code{0x4E} & \code{STA 14} & i ← i + 1 \\\tablerowrule
\code{10} & \code{0x60} & \code{JMP 0} & 回循环头 \\\tablerowrule
\code{11} & \code{0xF0} & \code{HLT} & 停机，时钟冻结 \\\tablerowrule
\code{12} & \code{0x01} & 常数 1 & 计数增量 \\\tablerowrule
\code{13} & \code{0x04} & 常数 4 & 累加上限 \\
\bottomrule
\end{longtable}

\tablechunkintro{3}{3}{\code{14}}{\code{15}}
\begin{longtable}{L{0.10\linewidth}L{0.12\linewidth}L{0.16\linewidth}L{0.50\linewidth}}
\toprule
地址 & 内容 & 助记符 & 注释 \\
\midrule
\code{14} & \code{0x01} & 变量 i & 编程模式预置 1，运行中被改写 \\\tablerowrule
\code{15} & \code{0x00} & 变量 sum & 编程模式预置 0，运行中被改写 \\
\bottomrule
\end{longtable}

以地址 6 为例读一遍编码：JZ 的 opcode 是 \code{0x8}，目标 11 即 \code{0xB}，合成 \code{0x8B}，拨码时这个字节是二进制 \code{10001011}。
把 16 个字节拨进机器的顺序是：上电、切到编程模式、从地址 0 到 15 逐字节操作——先拨地址开关，再拨数据开关，按一次写按钮；全部拨完切回运行模式，按复位，再启动时钟。
拨码时最容易犯的错是地址或数据拨串了一位，所以拨完把 16 个字节通读一遍再运行——读一遍只要一分钟，能省下半天的追查。
四圈循环的关键数据演化如下。
“A ← i − 4”一列是 SUB 13 的结果，ZF 由它决定；注意 ZF=1 全程只在第四圈出现一次。

\begin{longtable}{L{0.08\linewidth}L{0.11\linewidth}L{0.12\linewidth}L{0.12\linewidth}L{0.15\linewidth}L{0.08\linewidth}L{0.20\linewidth}}
\toprule
圈数 & i（入圈） & sum（加后） & OUT 显示 & A ← i − 4 & ZF & 去向 \\
\midrule
1 & 1 & 1 & 1 & \code{0xFD}（$-3$） & 0 & 继续，i 变 2 \\\tablerowrule
2 & 2 & 3 & 3 & \code{0xFE}（$-2$） & 0 & 继续，i 变 3 \\\tablerowrule
3 & 3 & 6 & 6 & \code{0xFF}（$-1$） & 0 & 继续，i 变 4 \\\tablerowrule
4 & 4 & 10 & 10 & \code{0x00} & 1 & JZ 11，停机 \\
\bottomrule
\end{longtable}

OUT 依次显示 1、3、6、10——部分和按升序出现，最后一个显示值就是答案 10。
把 OUT 放在循环体内而不是循环外，是为了让每一次累加都可见：只看最终结果，就无法区分“算对了”与“胡乱跳到一个恰好等于 10 的值”。
第四圈还有一个细节值得记下：SUB 13 算出 \code{0x00} 的同时 CF=1，因为 4 加 4 的取反再加 1 等于 256，向第 9 位进了位。
JZ 只看 ZF 不看 CF，所以这个进位无害。顺带一提，本程序里 i 到 4 时 CF 与 ZF 同时置 1，把判断换成 JC 恰好与 JZ 在同一点跳出，行为看不出差别；两条指令的语义差别要在 i 可能越过上限、CF 与 ZF 不再同置时才显现——选条件跳转前，先看清楚标志的定义。
再看时间维度。
复位后 T 计数器从 T0 起走；本机的 T 计数器没有提前清零逻辑，T0 到 T7 走满一圈才回绕，所以每条指令恰好占 8 拍。
指令提前执行完，剩余的拍就是空拍：所有控制线无效、总线高阻。
第一条指令 LDA 15 的完整 8 拍记录如下（复位后 PC=0；A 的旧值无关紧要，T3 会被覆盖）。

\begin{longtable}{L{0.06\linewidth}L{0.15\linewidth}L{0.06\linewidth}L{0.09\linewidth}L{0.10\linewidth}L{0.09\linewidth}L{0.31\linewidth}}
\toprule
拍 & 有效控制线 & PC & IR & 总线 & A & 说明 \\
\midrule
T0 & \code{CO+MI} & 0 & 旧值 & \code{0x00} & 旧值 & PC 把地址 0 送上总线，时钟沿 MAR ← 0，取指开始 \\\tablerowrule
T1 & \code{RO+II+CE} & 1 & \code{0x1F} & \code{0x1F} & 旧值 & RAM[0] 出到总线，IR ← \code{0x1F}，PC 同时加 1 \\\tablerowrule
T2 & \code{IO+MI} & 1 & \code{0x1F} & \code{0x0F} & 旧值 & IR 低 4 位送总线，MAR ← 15：操作数在 RAM[15] \\\tablerowrule
T3 & \code{RO+AI} & 1 & \code{0x1F} & \code{0x00} & \code{0x00} & RAM[15]（sum 的初值 0）出到总线，A ← 0 \\\tablerowrule
T4 & （无） & 1 & \code{0x1F} & 高阻 & \code{0x00} & LDA 已完成，空拍 \\\tablerowrule
T5 & （无） & 1 & \code{0x1F} & 高阻 & \code{0x00} & 空拍 \\\tablerowrule
T6 & （无） & 1 & \code{0x1F} & 高阻 & \code{0x00} & 空拍 \\\tablerowrule
T7 & （无） & 1 & \code{0x1F} & 高阻 & \code{0x00} & 空拍；下一拍回绕到 T0，开始取下一条 \\
\bottomrule
\end{longtable}

第 9 拍起进入第二条指令 ADD 14：T0 送地址 1，T1 取回 \code{0x2E}，T2 把操作数地址 14 送 MAR，T3 是 \code{RO+BI}——RAM[14]（此时 i=1）装入 B，T4 是 $\Sigma$O+\code{AI+FI}——A 得到 A+B=1，同时锁存 ZF 与 CF。
把这两拍与“控制单元：指令译码与微码”一节的微码表对照，编码、微码、时序三者就闭合了。
把 ADD 14 的关键五拍画成波形（第一圈：A=0x00，i=RAM[14]=0x01），装
载沿与总线值的对齐关系一眼可见。

\begin{waveform}[x=0.4cm,y=0.85cm]
  % T 状态翻转的下降沿参考虚线（x=4,8,12,16,20）
  \draw[hw return] (4,0.2) -- (4,9.7);
  \draw[hw return] (8,0.2) -- (8,9.7);
  \draw[hw return] (12,0.2) -- (12,9.7);
  \draw[hw return] (16,0.2) -- (16,9.7);
  \draw[hw return] (20,0.2) -- (20,9.7);
  % CLK：下降沿在 0/4/8/12/16/20/24（拍界，T 推进），上升沿在 2/6/10/14/18/22（装载）
  \draw[BookAccent, line width=0.9pt] (0,8.6) -- (2,8.6) -- (2,9.4) -- (4,9.4) -- (4,8.6) -- (6,8.6) -- (6,9.4) -- (8,9.4) -- (8,8.6) -- (10,8.6) -- (10,9.4) -- (12,9.4) -- (12,8.6) -- (14,8.6) -- (14,9.4) -- (16,9.4) -- (16,8.6) -- (18,8.6) -- (18,9.4) -- (20,9.4) -- (20,8.6) -- (22,8.6) -- (22,9.4) -- (24,9.4);
  \node[hw label, anchor=east] at (-0.2,9.0) {CLK};
  % T 状态
  \draw[BookTeal, line width=0.9pt] (0,7.0) rectangle (4,7.8);
  \node at (2,7.4) {T0};
  \draw[BookTeal, line width=0.9pt] (4,7.0) rectangle (8,7.8);
  \node at (6,7.4) {T1};
  \draw[BookTeal, line width=0.9pt] (8,7.0) rectangle (12,7.8);
  \node at (10,7.4) {T2};
  \draw[BookTeal, line width=0.9pt] (12,7.0) rectangle (16,7.8);
  \node at (14,7.4) {T3};
  \draw[BookTeal, line width=0.9pt] (16,7.0) rectangle (20,7.8);
  \node at (18,7.4) {T4};
  \draw[BookTeal, line width=0.9pt] (20,7.0) rectangle (24,7.8);
  \node at (22,7.4) {T5--T7};
  \node[hw label, anchor=east] at (-0.2,7.4) {T 状态};
  % 有效控制线（拍内稳定）
  \draw[BookTeal, line width=0.9pt] (0,5.4) rectangle (4,6.2);
  \node at (2,5.8) {CO+MI};
  \draw[BookTeal, line width=0.9pt] (4,5.4) rectangle (8,6.2);
  \node at (6,5.8) {RO+II+CE};
  \draw[BookTeal, line width=0.9pt] (8,5.4) rectangle (12,6.2);
  \node at (10,5.8) {IO+MI};
  \draw[BookTeal, line width=0.9pt] (12,5.4) rectangle (16,6.2);
  \node at (14,5.8) {RO+BI};
  \draw[BookTeal, line width=0.9pt] (16,5.4) rectangle (20,6.2);
  \node at (18,5.8) {$\Sigma$O+AI+FI};
  \draw[BookTeal, line width=0.9pt] (20,5.4) rectangle (24,6.2);
  \node at (22,5.8) {（无）};
  \node[hw label, anchor=east] at (-0.2,5.8) {控制线};
  % 总线值
  \draw[BookTeal, line width=0.9pt] (0,3.8) rectangle (4,4.6);
  \node at (2,4.2) {0x01};
  \draw[BookTeal, line width=0.9pt] (4,3.8) rectangle (8,4.6);
  \node at (6,4.2) {0x2E};
  \draw[BookTeal, line width=0.9pt] (8,3.8) rectangle (12,4.6);
  \node at (10,4.2) {0x0E};
  \draw[BookTeal, line width=0.9pt] (12,3.8) rectangle (16,4.6);
  \node at (14,4.2) {0x01};
  \draw[BookTeal, line width=0.9pt] (16,3.8) rectangle (20,4.6);
  \node at (18,4.2) {0x01};
  \draw[BookTeal, line width=0.9pt] (20,3.8) rectangle (24,4.6);
  \node at (22,4.2) {0x00};
  \node[hw label, anchor=east] at (-0.2,4.2) {总线};
  % A、B 寄存器（装载沿为拍内上升沿）
  \draw[BookTeal, line width=0.9pt] (0,2.2) rectangle (18,3.0);
  \node at (9,2.6) {0x00};
  \draw[BookTeal, line width=0.9pt] (18,2.2) rectangle (24,3.0);
  \node at (21,2.6) {0x01};
  \node[hw label, anchor=east] at (-0.2,2.6) {A};
  \draw[BookTeal, line width=0.9pt] (0,0.6) rectangle (14,1.4);
  \node at (7,1.0) {—};
  \draw[BookTeal, line width=0.9pt] (14,0.6) rectangle (24,1.4);
  \node at (19,1.0) {0x01};
  \node[hw label, anchor=east] at (-0.2,1.0) {B};
  % 装载动作（拍内上升沿）
  \node[hw label, anchor=north] at (14,-0.15) {装载沿 B←0x01};
  \node[hw label, anchor=north] at (18,-0.9) {装载沿 A←0x01，FI 记标志};
\end{waveform}
\diagramnote{程序一第二条指令 ADD 14 的 T0--T4（第一圈：A=0x00，i=RAM[14]=0x01）。T0、T1 是公共取指：总线依次是 PC 的 0x01 和指令字节 0x2E；T2 总线是 IR 低 4 位 0x0E，拍末 MAR←14；T3 总线是 RAM[14] 的 0x01，拍末 B←0x01；T4 $\Sigma$O 把 A+B=0x01 放上总线，拍末 AI 收回 A、FI 锁存标志。参考虚线是 T 状态翻转的下降沿：B 在 T3 的装载沿（拍内上升沿）更新，A 在 T4 的装载沿更新，控制字与总线值在装载沿前半拍已经稳定。}

读追踪表还要记住一个事实：空拍里总线高阻，LED 显示的是下拉后的 \code{0x00}，那不算任何模块在驱动总线。
判断“谁在驱动”，只看有效控制线里的“出”使能。
最后算一笔时间账。
每条指令固定 8 拍，程序一共执行 41 条指令：前三圈各 11 条，第四圈在 JZ 处跳走、只执行 8 条（LDA、ADD、STA、OUT、LDA、SUB、JZ 之后直达 HLT）。
逐条统计如下。

\begin{longtable}{L{0.14\linewidth}L{0.44\linewidth}L{0.18\linewidth}}
\toprule
指令 & 执行次数 & 拍数小计 \\
\midrule
\code{LDA} & 11（前三圈各 3，末圈 2） & 88 \\\tablerowrule
\code{ADD} & 7（前三圈各 2，末圈 1） & 56 \\\tablerowrule
\code{STA} & 7（前三圈各 2，末圈 1） & 56 \\\tablerowrule
\code{OUT} & 4 & 32 \\\tablerowrule
\code{SUB} & 4 & 32 \\\tablerowrule
\code{JZ} & 4 & 32 \\\tablerowrule
\code{JMP} & 3 & 24 \\\tablerowrule
\code{HLT} & 1 & 3（T2 冻结） \\\tablerowrule
合计 & 41 & 323 \\
\bottomrule
\end{longtable}

323 拍的含义很具体：1 Hz 的单步节奏下要五分半钟才能跑完，10 Hz 的 555 自动时钟下约 32 秒。
想亲眼看清每一拍，就把时钟切到单步按 323 次——这个数听起来大，按起来其实不慢，而且每一次都能在 LED 上核对总线值。

上电之前，先在纸上把逐圈表自己走一遍：从 i=1、sum=0 起，每圈写下 A、i、sum、ZF 四个值，直到 JZ 跳走。
纸上结果与机器显示一致，说明人与机器理解的是同一个程序；不一致，先怀疑纸上的自己——机器只是忠实执行你拨进去的字节。
桌面演练能查出两类错误：拨错的字节，和想错的循环，两类错误在上电前区分不开。
程序一还有一个天生的边界：8 位机器的累加和最大 255。
把上限从 4 改成 9（从 1 加到 9 得 45）依然安全，再往上就要先估算结果是否溢出——能显示的最大答案由数据宽度决定，这一点程序二同样躲不开。
顺带一提，JC 在本章两个程序里都没有出场：它适合“溢出即停”这类保护性判断，留给读者在改编程序时试用。

\topic{程序二：用循环加法做 3×4 乘法}ISA 里没有乘法指令，而乘法并没有消失——它变成了“计数加法”：3×4 就是把 3 自加 4 次。
只有加法器的机器上这是唯一办法；后来有乘法指令的处理器，在电路或微码里做的本质上也是这件事的加速版。
程序二与程序一结构同构：循环体仍是“加—显示—减—判断”，只是被加的量换成存在数据区的被乘数 3，判断对象换成递减到 0 的计数器。
数据区安排四个字节：地址 12 放常数 1（递减步长），地址 13 放被乘数 3，地址 14 是计数器 n、预置 4，地址 15 是积 p、预置 0。
与程序一不同，这次把判断放在圈尾：先把 3 加进 p，再让 n 减 1，减到 0 就停。
为什么放圈尾？这套 ISA 没有“与 0 比较”的指令，要测 n 是否为 0，只能借递减那一下的 ZF——把判断紧跟在 SUB 之后，减法就身兼“计数”与“测零”两职，省掉一条专门的比较指令。
圈尾判断还有一个关键细节：\code{SUB 12} 与 \code{JZ 9} 之间隔着一条 \code{STA 14}。
JZ 读的是标志寄存器，而 STA 不做任何 ALU 操作、\code{FI} 无效，ZF 仍是 SUB 算出的结果——标志不会被“路过”的指令破坏，这一点在“寄存器与 ALU 模块”一节装载标志的电路里就已经保证了。
完整编码如下。

\tablechunkintro{1}{3}{\code{0}}{\code{6}}
\begin{longtable}{L{0.10\linewidth}L{0.12\linewidth}L{0.16\linewidth}L{0.50\linewidth}}
\toprule
地址 & 内容 & 助记符 & 注释 \\
\midrule
\code{0} & \code{0x1F} & \code{LDA 15} & A ← p；循环从这里开始 \\\tablerowrule
\code{1} & \code{0x2D} & \code{ADD 13} & A ← p + 3 \\\tablerowrule
\code{2} & \code{0x4F} & \code{STA 15} & p ← p + 3 \\\tablerowrule
\code{3} & \code{0xE0} & \code{OUT} & 显示部分积 \\\tablerowrule
\code{4} & \code{0x1E} & \code{LDA 14} & A ← n \\\tablerowrule
\code{5} & \code{0x3C} & \code{SUB 12} & A ← n − 1；结果为 0 时 ZF=1 \\\tablerowrule
\code{6} & \code{0x4E} & \code{STA 14} & n ← n − 1；不影响标志 \\
\bottomrule
\end{longtable}

\tablechunkintro{2}{3}{\code{7}}{\code{13}}
\begin{longtable}{L{0.10\linewidth}L{0.12\linewidth}L{0.16\linewidth}L{0.50\linewidth}}
\toprule
地址 & 内容 & 助记符 & 注释 \\
\midrule
\code{7} & \code{0x89} & \code{JZ 9} & n 减到 0 就跳到 9 号停机 \\\tablerowrule
\code{8} & \code{0x60} & \code{JMP 0} & 否则回循环头 \\\tablerowrule
\code{9} & \code{0xF0} & \code{HLT} & 停机 \\\tablerowrule
\code{10} & \code{0x00} & \code{NOP} & 填充字节，正常流程到不了 \\\tablerowrule
\code{11} & \code{0x00} & \code{NOP} & 填充字节，正常流程到不了 \\\tablerowrule
\code{12} & \code{0x01} & 常数 1 & 递减步长 \\\tablerowrule
\code{13} & \code{0x03} & 被乘数 3 & 要算几乘几，改这里与计数器 \\
\bottomrule
\end{longtable}

\tablechunkintro{3}{3}{\code{14}}{\code{15}}
\begin{longtable}{L{0.10\linewidth}L{0.12\linewidth}L{0.16\linewidth}L{0.50\linewidth}}
\toprule
地址 & 内容 & 助记符 & 注释 \\
\midrule
\code{14} & \code{0x04} & 计数器 n & 编程模式预置 4，运行中被改写 \\\tablerowrule
\code{15} & \code{0x00} & 积 p & 编程模式预置 0，运行中被改写 \\
\bottomrule
\end{longtable}

程序对不对，用循环不变量一句话就能说清：每次到达地址 0 时，恒有 p + 3n = 12。
进入循环时 p=0、n=4，0+3×4=12，不变量成立。
每转一圈，p 增加 3、n 减少 1，p + 3n 的值不变，不变量保持。
退出循环时 n=0，于是 p=12，正是 3×4。
不变量把“算多少次”与“每次加多少”锁在一起，是计数加法类程序通用的论证方法。
四圈的数据演化如下。

\begin{longtable}{L{0.08\linewidth}L{0.11\linewidth}L{0.12\linewidth}L{0.12\linewidth}L{0.11\linewidth}L{0.08\linewidth}L{0.24\linewidth}}
\toprule
圈数 & n（入圈） & p（加后） & OUT 显示 & n − 1 & ZF & 去向 \\
\midrule
1 & 4 & 3 & 3 & 3 & 0 & 继续 \\\tablerowrule
2 & 3 & 6 & 6 & 2 & 0 & 继续 \\\tablerowrule
3 & 2 & 9 & 9 & 1 & 0 & 继续 \\\tablerowrule
4 & 1 & 12 & 12 & 0 & 1 & JZ 9，停机 \\
\bottomrule
\end{longtable}

OUT 依次显示 3、6、9、12，共转 4 圈——放在循环里的 OUT 等于把乘法表的一行打了出来；若只要答案，把 OUT 挪到 HLT 前即可，但那就少了一次观察计数过程的机会。
拍数同样好算：前三圈各 9 条指令，末圈在 JZ 处跳走、连 HLT 共 9 条，合计 36 条；前 35 条各 8 拍，HLT 在 T2 冻结只走 3 拍，共 $35\times8+3=283$ 拍。
把“乘法 = 计数加法”再推进一步：除法就是计数减法——反复减去除数、数减了几次；乘方是嵌套的计数加法。
这台 16 字节的小机器装不下那些程序，但道理完全一样：没有某条指令，不等于没有那种运算，只是要多花几条指令、几拍时间。
想算 5×6 也不必改程序：把地址 13 拨成 5、地址 14 拨成 6，机器就算另一道题——乘法与程序无关，只与数据有关。

程序二还有两个细节值得多说一句。
一是被乘数为什么放在数据区，而不是用 LDI 内嵌进程序：内嵌也能算，但那样每换一道题就要改程序字节；被乘数放数据区，同一份程序对任何乘法题都成立。
二是 10、11 两个字节为什么填 NOP：编程模式下每个字节都要拨出确定值，NOP（\code{0x00}）是最安全的填充。
正常流程不会到达它们；万一调试时乱按时钟把 PC 送过了 9 号，NOP 至少不会意外改写 RAM——当然，PC 继续滑进 12 号之后执行的就是“数据当指令”，所以填充只是对齐手段，不是保险。
三是乘积的位宽：p 同样只有 8 位，3×4=12 很安全，16×16=256 就溢出了。
挑乘法题要按结果位宽挑，超过 255 的乘法需要两个字节存积，那是下一台机器的程序。

\topic{故障排查案例三则}程序跑不起来时，故障几乎总在三类地方：总线、控制线、复位。
下面三则案例按“症状、定位、根因、修复”各讲一遍；它们也是本节末尾整机检查清单里最容易栽跟头的三项。

案例一：总线争用。
\textbf{症状}某一拍总线 LED 显示一个谁也不认识的值——既非 RAM 内容也非任何寄存器内容，个别 LED 亮度发暗；整机电流明显偏大，摸相关芯片微微发热。
\textbf{定位}把时钟停在出事的拍，列出该拍微码里全部“出”使能（\code{CO}、\code{RO}、\code{AO}、$\Sigma$O、\code{IO}），用逻辑笔逐根点：任何一拍只能有一根为高；发现两根同时为高，就找到了争用的双方。
\textbf{根因}使能互斥被破坏：微码 ROM 某一行多烧了一位“出”使能，或者两根使能线在面包板上插错了孔，或者某片三态输出的使能端极性接反、输出常开。
\textbf{修复}改微码或改接线；此后把“任何时刻总线只有一个驱动者”当作硬约定——每改动一根使能线，就把 11 条指令的全部微码行重新核一遍，上电先单步空走 8 拍，看总线有无意外驱动。

案例二：\code{RI}、\code{RO} 接反。
\textbf{症状}读变写、写变读：LDA 之后 A 拿到的是总线上的随机值；更糟的是 RAM 被意外改写——执行完 STA 15，15 号单元没变、别的单元却变了，程序越跑越乱，因为程序代码正在被意外改写。
\textbf{定位}拨入一条 LDA 15 加 HLT，单步到 T3，用两根逻辑笔同时点 \code{RI} 与 \code{RO}：这一拍应当 \code{RO}=1、\code{RI}=0，若正好相反就实锤了。
另有一个旁证：跑完任何程序后，把几个没写过的 RAM 单元倒出来看，若内容与拨入时不同，说明写使能在不该有效的时候有效过。
\textbf{根因}微码 ROM 到 RAM 的两根控制线在排线中相邻，最容易整对调；也可能是烧 ROM 时两行控制字填串了。
\textbf{修复}断电，对调两根线（或重烧 ROM），再做一次全指令抽查：每条指令跑完，核对相关 RAM 单元未被意外改写。
顺手把 \code{RI} 的时序也核一遍：\code{RI} 必须在 MAR 地址稳定之后才有效，且只在需要写的那一拍有效。

案例三：复位未清 T 状态。
\textbf{症状}按复位后 PC 确实回到 0，但第一条指令执行得莫名其妙——像是从半中间开始，少了一个取指拍，或多出一个不该有的写 RAM 动作；每次上电犯的错还不一样。
\textbf{定位}复位后连续单步，观察有效控制线：第一拍必须是 \code{CO+MI}。
若看到的第一拍是 \code{IO+MI} 或别的组合，说明 T 计数器没有从 T0 起走。
\textbf{根因}复位线只接了 PC 的清零端，没接 T 计数器（74HC161）的清零端；T 停在复位前的那一拍——比如 T3——取指序列整个错位：该取指时在执行，该执行时在取指。
\textbf{修复}复位必须覆盖所有时序状态：PC、T 计数器、IR、标志寄存器同时清零；A、B、OUT 这些数据通路寄存器可以保留，反而便于上电观察。
改完后的检查：复位后单步 8 拍，控制线序列必须逐拍与微码表一致，且每次上电结果都一样。

三类故障的速查对照如下，排查时先对症状、再用首选手段确认，不要盲目换芯片。

\begin{longtable}{L{0.16\linewidth}L{0.40\linewidth}L{0.32\linewidth}}
\toprule
故障类型 & 典型症状 & 首选定位手段 \\
\midrule
总线争用 & 某拍总线显示谁也不认识的中间值，电流偏大 & 停在该拍，逐根点全部“出”使能 \\\tablerowrule
控制线接反 & 读变写、写变读，RAM 被意外改写 & 单步到相关拍，同时点 \code{RI} 与 \code{RO} \\\tablerowrule
复位不全 & 复位后第一拍不是 \code{CO+MI}，每次上电错得不一样 & 复位后连续单步，核对控制线序列 \\
\bottomrule
\end{longtable}

三则案例对应三条搭机纪律：使能互斥逐拍核、相邻控制线双人核对、复位覆盖全部时序状态。
它们也有共同点：定位手段都是“单步加逻辑笔”，没有一件工具超出实验盒的范围。
机器小到没有一处不可观察，是它作为教具最大的本事。

三则案例之外还有一条排查总原则：先怀疑最后动过的地方，一次只改一处，改完立刻重跑最小复现程序。
同时改两处，等于亲手毁掉对照组。
这几条听起来像软件调试，因为调试的本质不分软硬件：控制变量、最小复现、对照实验。
另有一类不在三则之列的隐性问题值得点名——接触不良：面包板孔内的簧片会疲劳，症状是“轻碰板子，程序行为就变”，单步和逻辑笔都抓不住它。
预防胜于定位：走线横平竖直、导线插到底、不用过长的飞线，能躲开一大半这类问题。

\topic{扩展方向：更大、更快、更多指令}这台机器故意做得很小，但每个方向都留着扩的口子。
下表按“要动哪些模块”把四条最常见的扩展路径排出来，随后逐条细说。

\begin{longtable}{L{0.18\linewidth}L{0.36\linewidth}L{0.34\linewidth}}
\toprule
扩展方向 & 要动的模块 & 主要代价 \\
\midrule
RAM 扩到 256 字节 & MAR 加宽到 8 位、RAM 芯片、编程模式拨码 & 低 4 位放不下地址，须改两字节指令格式，微码全部重排 \\\tablerowrule
加立即数指令（ADI） & 只动微码 ROM：T2 发 \code{IO+BI}，T3 发 $\Sigma$O+\code{AI+FI} & 最便宜的扩展：不动数据通路，只烧两行控制字 \\\tablerowrule
加 CALL/RET 与栈 & 新增 SP 计数器模块、RAM 划出栈区、微码序列拉长 & 最贵的扩展：PC 要新增经总线进出 RAM 的路径 \\\tablerowrule
提高时钟频率 & 不加模块，重算关键路径：总线 → ALU → A 建立时间 & 面包板引线电容压低实际上限，频率须留一倍余量 \\\tablerowrule
微码 ROM 换成 RAM & 微码存储改静态 RAM，另加上电加载电路 & 得到可写控制存储；加载电路本身成为新的故障源 \\
\bottomrule
\end{longtable}

扩 RAM 是最诱人的方向，也是连锁反应最多的：地址从 4 位变 8 位，MAR 要加一片寄存器，RAM 要换大芯片，而最大的改动在指令格式——一条指令一个字节，低 4 位最多指到 15。
解决办法是学真实机器：指令占两个字节，第一字节 opcode、第二字节地址，取指之后多插一拍把地址字节读进来；于是每条指令的微码序列都要重排，本章所有编码表作废。
这也解释了为什么本章把地址空间限定在 16 字节：小，才能让每张表完整印在一页里。
加指令则展示了微码结构的甜头。
以立即数加法 ADI（A ← A + 立即数）为例：数据通路什么也不用加——IR 低 4 位本来就能上总线（\code{IO}），B 本来就能装载（\code{BI}），T2 发 \code{IO+BI}、T3 发 $\Sigma$O+\code{AI+FI}，两行控制字烧进 ROM 就完了。
硬连线控制里加一条指令要重画译码逻辑，微码控制里只是往表里多写两行——“微码便宜、硬线贵”，这是值得亲手烧一次 ROM 的理由。
CALL 与 RET 是另一个极端：它们需要栈，而栈需要新硬件。
栈指针 SP 是一个能加一减一的计数器模块；CALL 要把 PC 的当前值经总线写进 RAM[SP]、SP 减一、再把目标地址装进 PC，RET 反过来把 RAM[SP] 读回 PC。
每条都要 6 到 8 个执行拍，PC 还得新增“把值送上总线”之外的路径——在 16 字节的机器上做栈得不偿失，但它标出了下一台机器的方向。
提高时钟不需要加任何模块，只需要算账。
本机最长的一拍是 ADD 的 T4：操作数在 T3 已装进 B，T4 这一拍要走过 74HC283 的行波进位（8 位逐级传递）、总线传播和 A 寄存器的建立时间，全部加起来取倒数，就是时钟频率的理论上限。
面包板引线长、分布电容大，实际上限比器件手册算出来的低得多，所以预算规则是：算出上限再留一倍余量，555 的阻容按这个值配。
把这段关键路径的延迟组成画成分段条，各级谁占多少、余量留在哪里就一目了然。

\begin{pdfdiagram}
  % 关键路径延迟分段 + 一倍余量
  \node[hw label, anchor=east] at (-0.15,0.4) {ADD 的 T4};
  \draw[BookTeal, line width=0.9pt] (0,0) rectangle (5.0,0.8);
  \draw[BookTeal, line width=0.9pt] (2.8,0) -- (2.8,0.8);
  \draw[BookTeal, line width=0.9pt] (3.9,0) -- (3.9,0.8);
  \node[hw label] at (1.4,0.4) {74HC283 行波进位};
  \node[hw label] at (3.35,0.4) {总线传播};
  \node[hw label] at (4.45,0.4) {A 建立时间};
  \draw[BookMuted, dashed, line width=0.7pt] (5.0,0) rectangle (10.0,0.8);
  \node[hw label] at (7.5,0.4) {余量：上限再留一倍};
  \draw[BookMuted, line width=0.6pt, <->] (0,1.15) -- node[hw label, above]{关键路径总延迟，取倒数得频率理论上限} (5.0,1.15);
  \draw[BookMuted, line width=0.6pt, <->] (0,-0.35) -- node[hw label, below]{实际一拍：555 阻容按“上限加一倍余量”配} (10.0,-0.35);
  \node[hw label, text width=10cm] at (5,-1.4) {面包板引线长、分布电容大，实际能跑的上限比按器件手册算出的低，余量就是留给这些分布参数的；要抬高上限，先把行波进位换成超前进位。};
\end{pdfdiagram}
\diagramnote{整机性能预算：ADD 的 T4 是全机最长的一拍，关键路径由 74HC283 的行波进位（8 位逐级传递）、总线传播和 A 寄存器建立时间三段串联而成，总延迟取倒数就是时钟频率的理论上限。实际配置再留一倍余量——分段条右半的虚线部分，555 的阻容按这个加了余量的周期来配。}
最后一个方向是把微码 ROM 换成 RAM，得到可写控制存储：上电时由加载电路把微码灌进静态 RAM，之后运行时也能改指令定义——今天 OUT 是显示，明天可以改成“显示并响蜂鸣器”。
代价是加载电路本身成为新的故障源，而且微码错了，机器会以最难调试的方式出错：每条指令都不对。
四条路径没有先后，但有一条共同纪律：每扩一处，先回章首把整机规格表被改动的行改掉，再重做本节末尾的整机检查。

有读者会问中断和进位链：它们其实已经藏在表里。
进位链藏在“提高时钟”那一行——74HC283 的行波进位正是关键路径的主角，换超前进位器件能直接抬高频率上限。
中断则藏在 CALL/RET 那一行：中断的本质是“硬件替你执行一次 CALL”，同样要保存 PC 现场，同样需要栈；没有栈的机器谈中断，就像没有 RAM 的机器谈子程序。
所以扩展的顺序不是随机的：先有栈，才谈得上子程序；再有富余的地址空间，才谈得上中断向量。

把四条路径按目标画成决策图，每条路径内部的先后就写在链上：

\begin{pdfdiagram}
  % 目标一：加指令，按成本递增
  \node[hw state, align=center] (ga) at (0.2,2.0) {更多指令};
  \node[hw compute, align=center, minimum width=2.6cm] (s1) at (3.7,2.0) {ADI\\新增两行控制字};
  \node[hw compute, align=center, minimum width=2.4cm] (s2) at (7.1,2.0) {两字节指令格式\\微码全表重排};
  \node[hw compute, align=center, minimum width=2.4cm] (s3) at (10.3,2.0) {CALL/RET\\新增 SP 模块};
  \draw[hw bus] (ga.east) -- node[midway, above=1pt, hw label, fill=white, inner sep=1pt]{先} (s1.west);
  \draw[hw bus] (s1.east) -- node[midway, above=1pt, hw label, fill=white, inner sep=1pt]{再} (s2.west);
  \draw[hw bus] (s2.east) -- node[midway, above=1pt, hw label, fill=white, inner sep=1pt]{最后} (s3.west);
  % 目标二：扩 RAM，多个约束汇合到同一格式改动
  \node[hw state, align=center] (gb) at (0.2,0.1) {更大 RAM};
  \node[hw compute, align=center, minimum width=3.0cm] (t1) at (3.7,0.1) {MAR 加宽到 8 位\\RAM 换大芯片};
  \node[hw compute, align=center, minimum width=2.4cm] (t2) at (7.1,0.1) {低 4 位无法容纳地址\\因此需要两字节格式};
  \draw[hw bus] (gb.east) -- (t1.west);
  \draw[hw bus] (t1.east) -- (t2.west);
  \draw[hw return] (t2.north) -- (s2.south);
  \node[hw label, fill=white, inner sep=1pt] at (8.0,1.05) {格式改动汇合};
  % 目标三：提时钟，不增加模块，只核算时序
  \node[hw state, align=center] (gc) at (0.2,-1.8) {更高时钟};
  \node[hw compute, align=center, minimum width=2.6cm] (u1) at (3.7,-1.8) {不加模块\\重算关键路径};
  \node[hw compute, align=center, minimum width=3.0cm] (u2) at (7.1,-1.8) {上限留一倍余量\\555 按值配阻容};
  \draw[hw bus] (gc.east) -- (u1.west);
  \draw[hw bus] (u1.east) -- (u2.west);
\end{pdfdiagram}
\diagramnote{扩展路径决策（对应正文扩展方向表）。三个目标各有相应的实现路径：加指令按成本递增排——ADI 只需新增两行控制字，成本最低；两字节指令格式要重排全表微码；CALL/RET 与栈需要新增 SP 计数器模块，成本最高。扩 RAM 连锁最多：MAR 加宽、RAM 换大之外，低 4 位无法容纳地址，因此同样需要两字节格式。提高时钟不加模块，只重算“总线到 ALU 到 A 建立时间”这条关键路径，理论上限再留一倍余量。所以扩展的顺序不是随机的：先有栈才谈得上子程序，再有富余的地址空间才谈得上中断向量。}

\topic{整机上电检查清单}整机第一次上电，不要直接跑程序。
按“静态、单模块、单指令、程序”四级往上走，每一级都把故障隔离在尽可能小的范围：静态不过谈不上模块，模块不过谈不上指令，指令不过谈不上程序。
四级顺序本质上也是二分法：每通过一级，故障可能被锁定的范围就缩小一半以上。

\begin{longtable}{L{0.13\linewidth}L{0.42\linewidth}L{0.33\linewidth}}
\toprule
层级 & 检查内容 & 预期现象 \\
\midrule
静态（不上时钟） & 目检焊点与走线；量电源与地之间的电压；复位有效时看 PC、T、IR、标志是否全为已知值；复位后 T=0、IR 为 \code{0x00}，微码 ROM 组合输出第 0 行 \code{0x4008}，仅 CO、MI 有效，PC 驱动总线 & 总线 LED 全灭（PC=0 驱动总线，读 \code{0x00}）；整机电流在几十毫安量级；无芯片发热 \\\tablerowrule
单模块 & 单步下逐一激励各模块：PC 计数与装载、A/B/IR 装载、RAM 单字节读写、ALU 加减与 ZF/CF、OUT 显示 & 各模块孤立动作与本章“寄存器与 ALU 模块”一节、“内存与输出模块”一节的分模块检查结果一致 \\\tablerowrule
单指令 & 拨入 \code{LDA 15} 与 \code{HLT}，单步走完两圈 8 拍；再抽验 \code{ADD}、\code{STA}、\code{JZ}、\code{OUT} 各一条 & 每拍有效控制线与“控制单元：指令译码与微码”一节微码表逐拍一致；任何一拍总线只有一个驱动者 \\\tablerowrule
程序 & 拨入程序一、程序二的完整 16 字节，自动时钟跑到底 & OUT 依次显示 1、3、6、10（程序一）与 3、6、9、12（程序二）；HLT 后时钟冻结、状态稳定可读 \\
\bottomrule
\end{longtable}

任何一级不过，就停在那一级修好再往上走——跨级调试是最常见的浪费时间方式。
这张清单值得打印出来贴在面包板旁边，每次上电按顺序走一遍。
它还有一条使用规则：每次改动接线或重烧微码之后，无论改动多小，都从静态级重新走起——改动之后直奔程序级，是最常见的跳级错误。
两个程序的输出序列 1、3、6、10 与 3、6、9、12 是这台整机的最终签名：它们对，说明从 555 的 RC 到微码 ROM 的每一位都对。
到这里，从“数字硬件基础”部分的第一个开关到本章的一台可编程计算机，整条链路闭合。

\topic{专题：整机联调日志——从一次真实 bring-up 看排错顺序}下面这份
日志来自一次真实的整机 bring-up：各模块在分测阶段都单独过了一遍，第
一次连成整机跑程序，从上午十点到傍晚六点，一共踩了六个坑。日志按时
间顺序原样保留，每一步只记四件事——现象、第一怀疑、测量点、结论与
改动。

\tablechunkintro{1}{3}{① 上电静态检查}{③ 复位后 PC 不归零}
\begin{longtable}{L{0.09\linewidth}L{0.22\linewidth}L{0.17\linewidth}L{0.19\linewidth}L{0.21\linewidth}}
\toprule
步骤 & 现象 & 第一怀疑 & 测量点 & 结论与改动 \\
\midrule
① 上电静态检查 & 尚未通电：蜂鸣档量电源轨，5 V 与 GND 直通 & 有跳线插错孔，或某片芯片插反 & 逐段拔跳线缩小范围；目检全部芯片缺口方向 & 一根跳线把电源轨与地轨短在相邻孔；另有一片 74HC273 缺口朝反。挪线、掉头；上电后整机电流几十毫安，无芯片发热，逐片去耦电容在位 \\\tablerowrule
② 时钟不起振 & 555 输出脚 LED 常亮不闪，预期 1--50Hz 连续方波 & 555 损坏，或芯片插反 & LED 逐点看：输出脚半亮不灭，阈值脚也半亮——半亮说明在快速开关，只是肉眼跟不上 & 定时电容错一档：10$\mu$F 插成 0.01$\mu$F，频率高了三个数量级，“不起振”其实是“快得看不见”。换回 10$\mu$F，电位器拧到 1Hz 起振 \\\tablerowrule
③ 复位后 PC 不归零 & 按复位按钮，PC 不归零，反而开始计数；一松手，PC、T、IR 全被清掉——相位整个反了 & PC 那片计数器的清零脚虚接 & 复位链末级输出脚的常态电平：低有效清零端平时应为高，实测平时为低 & 按钮消抖级的上下拉接反，复位链整体反相：平时等于按住复位，按下等于松手。改正后按一下全机归零，松手从 T0 起跑 \\
\bottomrule
\end{longtable}

\tablechunkintro{2}{3}{④ 取指取回 \code{0xFF}}{⑤ ADD 结果错 1}
\begin{longtable}{L{0.09\linewidth}L{0.22\linewidth}L{0.17\linewidth}L{0.19\linewidth}L{0.21\linewidth}}
\toprule
步骤 & 现象 & 第一怀疑 & 测量点 & 结论与改动 \\
\midrule
④ 取指取回 \code{0xFF} & 复位后单步：T0 总线 LED 全亮（应显示 PC 的 \code{0x0}），T1 取回 \code{0xFF}，机器随即停在 HLT & 程序字节拨错，或微码 ROM 取指行烧错 & 逐拍看总线：T0 拍就不是一个驱动者该有的读数；再量 74HC240 的使能脚——悬空，被器件读成有效 & RO 反相器到 240 使能脚一根线漏插，RAM 全程驱动总线，T0 与 PC 争用，两边输出级对抗把总线顶成全亮。补线后 T0 显示 \code{0x0}，T1 取回程序第一字节 \\\tablerowrule
⑤ ADD 结果错 1 & 拨入“内存与输出模块”一节的 5 字节程序：\code{LDI 6}、\code{ADD 0x4}（存 7），应显示 \code{0x0D}，实测 \code{0x0E}，恰好大 1；SUB 却正确 & 微码 ROM 里 ADD 的 T4 行烧错一位 & 加法器输出那排常驻观察 LED：A=6、B=7、SU=0 时已显示 \code{0x0E}，错在 ALU 内部，与微码无关；再量最低位 283 的 CI 脚：悬空 & SU 的第九个分叉漏接：8 个异或门都接了，进位输入 CI 没接，悬空被读成 1，加法变 A+B+1；减法 CI 本来就要 1，故碰巧正确。补线后显示 \code{0x0D} \\
\bottomrule
\end{longtable}

\tablechunkintro{3}{3}{⑥ 程序跑通}{⑥ 程序跑通}
\begin{longtable}{L{0.09\linewidth}L{0.22\linewidth}L{0.17\linewidth}L{0.19\linewidth}L{0.21\linewidth}}
\toprule
步骤 & 现象 & 第一怀疑 & 测量点 & 结论与改动 \\
\midrule
⑥ 程序跑通 & 拨入程序一全部 16 字节，自动时钟跑到底：OUT 依次显示 1、3、6、10，HLT 后时钟冻结、读数稳定 & 无须怀疑，只须确认 & 先单步走完第一条指令的 8 拍，逐拍对微码表；再拨自动挡 & 签名序列 1、3、6、10 全部对上。这一步不改任何线，只做回归：从静态级把四级检查重走一遍，全部通过 \\
\bottomrule
\end{longtable}

排错动作可以写成一个循环，贴在面包板旁边比任何口诀都顶用：

\begin{lstlisting}
// 注释（推进）：先固定现象和测量点，再做一次单变量改动。
// 注释（边界）：改动后回到上一边界复测，确认故障没有迁移到下一级。
while 行为与预期不符:
    1. 停在当前层级，写下现象与第一怀疑
    2. 选一个测量点，先量再改
    3. 一次只改一处，改完退回上一级重查
    4. 本级全部通过，才允许走向下一级
\end{lstlisting}

把这个循环画成决策树，“怀疑提出假设、测量负责否决”的走法就固定下来：

\begin{pdfdiagram}
  % 主干
  \node[hw state, align=center, minimum width=3.4cm] (sym) at (4.4,2.4) {症状：行为与预期不符};
  \node[hw compute, align=center, minimum width=3.4cm] (sus) at (4.4,1.1) {写下现象与第一怀疑};
  \node[hw control, align=center, minimum width=3.4cm] (meas) at (4.4,-0.2) {选测量点：\\一刀分在嫌疑范围中间};
  \node[hw compute, align=center, minimum width=3.4cm] (judge) at (4.4,-1.6) {读数与怀疑相符？};
  % 左右分支
  \node[hw state, align=center, minimum width=2.6cm] (rej) at (0.5,-3.0) {怀疑作废\\换假设再量};
  \node[hw state, align=center, minimum width=2.6cm] (fix) at (8.3,-3.0) {定位根因\\一次只改一处};
  \node[hw control, align=center, minimum width=2.6cm] (ret) at (8.3,-4.6) {退回上一级重查\\全过才走向下一级};
  \draw[hw bus] (sym.south) -- (sus.north);
  \draw[hw bus] (sus.south) -- (meas.north);
  \draw[hw bus] (meas.south) -- (judge.north);
  \draw[hw bus] (judge.west) -- node[midway, above=1pt, hw label, fill=white, inner sep=1pt]{不符} (0.5,-1.6) -- (rej.north);
  \draw[hw bus] (judge.east) -- node[midway, above=1pt, hw label, fill=white, inner sep=1pt]{相符} (8.3,-1.6) -- (fix.north);
  \draw[hw bus] (fix.south) -- (ret.north);
  % 否决回路：回到“换假设”
  \draw[hw return] (rej.west) -- (-1.3,-3.0) -- (-1.3,1.1) -- (sus.west);
  \node[hw label, rotate=90] at (-1.58,-1.0) {被否决的怀疑留在纸上};
\end{pdfdiagram}
\diagramnote{故障排查决策树（对应联调日志六步的走法）。症状写下后提出第一怀疑，怀疑的职责是指定第一个测量点——选在与各嫌疑都无关、一刀分在范围中间的位置量：读数与怀疑不符，怀疑就地作废、换假设再量；相符才动手，一次只改一处；改完退回上一级重查，本级全过才允许走向下一级。日志里六个坑平均八十分钟一个，靠的就是这条看起来慢的纪律。}

第一条纪律是一次只改一处，理由是纯因果的：一次改两处，机器好了，你
不知道是哪处起的作用；更常见的是机器没好，而你从此分不清“原本的故
障”和“新引入的故障”。第④步最能说明问题：取指取回 \code{0xFF} 的
时候，手边至少有三个诱人的改动——重拨程序、重烧微码、补 240 那根
线。若是把前两个“顺手”做掉再补线，机器也能跑通，但这份日志就此失
去结论，下次遇到同样的症状还得从头再来。忍住不改、先量使能脚，才把
故障定位到那根漏插的线上。所以每次动手之前先问一句：这次改动针对哪一
条测量？答不上来，就说明还没到改的时候，先回去量。

第二条纪律是改完退回上一级。本章的四级检查顺序——静态、单模块、单
指令、程序——把整机切成四层，任何一处改动之后，从被改动的那一层重
走，而不是从刚才停下的地方继续。第②步换完电容，不是直接再跑程序，
而是退回单模块级，把时钟的两路输出重新核一遍；第④步补完线，退回单
指令级重核 T0、T1 两拍，而不是直接拨自动挡。退回的代价是几分钟，跳
级的代价是故障在上一层复活时，你得再花一整轮从头定位。六个坑平均八十分钟一个，靠的就是这条看起来慢的纪律；在面包板上，一天之内从通电
到跑通，已经算是一场顺利的 bring-up。

第三条要说的是“第一怀疑”的用法。六个坑里，第一怀疑只命中了一次，
这不是记录者笨拙，而是硬件故障的常态：第一怀疑来自模式匹配，而模式
匹配在跨模块的故障面前经常指错方向。它的真正用途不是命中，而是决定
第一个测量点——怀疑 555 坏了，于是先量输出脚；量出来是半亮，与怀疑
不符，怀疑就地作废，测量领着你走向电容。怀疑可以被推翻，测量不会说
谎；日志里的每一行，都是“怀疑提出假设、测量负责否决”的一次循环。
把怀疑写下来还有一层好处：被否决的怀疑留在纸上，两个小时后你就不会
把它当成新想法再试一遍。

六条记录之外，表里还有一条暗线值得挑明：测量点的选法。好的测量点要
一刀剁在故障可能范围的中间——怀疑程序拨错还是微码烧错，就量一个与
两者都无关的点：240 的使能脚。它正常，ROM 与程序的嫌疑同时洗清；它
不正常，两边都不用再拆。第⑤步的测量点尤其典型：加法器输出那排常驻
LED 把“微码错”与“ALU 错”一刀分开，一次读数否决一半可能。选测量
点的本领，说到底是二分法的本领：每一次测量都将剩余可能性排除一半，
六次二分足以从六十四个候选点中定位一个——而一块面包板上的嫌疑，通常
凑不齐这个数。

每一条“结论与改动”栏的末尾还有一个新读数：1Hz 起振、T0 显示
\code{0x0}、OUT 显示 \code{0x0D}。这不是行文习惯，是纪律的最后一环：
改动之后必须有一个能证明改动生效的测量，否则这次改动不算完成。修机
器最常见的自欺，是改完线看一眼“好像好了”就往下走，三天后故障换一
副面孔回来，没人记得它上次是怎么“好”的。新读数把每次修复限定在一
个可复查的状态上：它既是对本次改动的确认，也是下一次退回上一级时的
起跑线。

最后一层是记录本身。bring-up 到第三个小时，人会开始原地转圈：这根
线量过没有？这组读数是改动前还是改动后的？纸面上的日志把记忆问题变
成查表问题，也把一下午的六个坑变成下次搭建的六条经验——第二次见到
“总线全亮”，你会直接去量 240 的使能脚，而不是再怀疑一遍程序。日
志的形式不拘：纸本、纯文本、手机照片都行，要紧的只有三点——四栏齐
全、按时间排列、每次改动都配一个新读数。满足这三点，铅笔头和键盘没
有区别。这台机器造好之后，真正陪你走向下一台机器的，不是面包板，是
这本日志。

\topic{专题：面包板常见的十类硬件故障}
下表按发生频率排出十类，每类给三栏：典型症状、最快定位法、一行实测
经验。症状与故障不是一一对应的——“时好时坏”可能是接触不良，也可
能是悬空脚——所以表里最值钱的是定位法那一栏：它回答“怎样用一次测
量把这一类排除掉”。

\begin{longtable}{L{0.15\linewidth}L{0.25\linewidth}L{0.24\linewidth}L{0.24\linewidth}}
\toprule
故障类别 & 典型症状 & 最快定位法 & 实测经验 \\
\midrule
接触不良 & 时好时坏，晃动板子或按压芯片时现象跟着变 & 分区轻晃、镊子逐个压芯片，找到敏感区再逐孔重插 & 面包板插孔有寿命，“软故障”查到最后，八成是它 \\\tablerowrule
芯片插反 & 上电即发热，该芯片无输出 & 立即断电，摸温度找发烫的片，目检缺口朝向 & 全板缺口统一朝向的纪律，能防掉其中大半 \\\tablerowrule
电源轨半板未接 & 一半模块正常，另一半整体不工作 & 不量电源入口，逐片量芯片的 VCC 脚 & 长面包板的电源轨在中段断开，两半要分别跨接 \\\tablerowrule
未用输入悬空 & 输出随机翻转，手靠近时读数会变 & 对照手册把每片没用到的输入脚数一遍 & 悬空脚是天线，全部按功能接高或接低，一个不留 \\\tablerowrule
三态使能互斥遗漏 & 总线停在中间电平，LED 全亮，整机电流变大 & 逐拍核对有效控制线，“出”使能必须只有一个 & 争用几分钟后输出级会微热，趁热摸能找到当事片 \\\tablerowrule
消抖不足 & 按一次按钮，机器走好几拍 & 看 T 计数器在一次按键里的跳变次数 & 触点抖动几毫秒，RC 取 10ms 才盖得住 \\\tablerowrule
控制线交叉接反 & 功能成对错位：该读时写、该装 A 装了 B & 把可疑的两根对调验证；平时逐根对照控制线清单 & 同色跳线并排最容易拿错，按功能分色省一半误会 \\\tablerowrule
下拉电阻遗漏 & 空闲拍总线读随机值，而不是 \code{0x00} & 让所有驱动者无效，读总线空闲值 & 100k$\Omega$ 一排 8 只，少一只，那一位就漂 \\\tablerowrule
LED 限流电阻错档 & 显示过暗、刺眼，或干脆烧灯 & 量单路电流，与 330$\Omega$ 档的计算值对照 & 33$\Omega$ 与 330$\Omega$ 只差一环色环，亮度差十倍 \\\tablerowrule
去耦电容忘接 & 低速正常，提高时钟后偶发出错 & 故障只在提速后出现时，先数电容再查别的 & 0.1$\mu$F 贴着电源脚放，电流突变全靠它就地吸收 \\
\bottomrule
\end{longtable}

这张表的顺序值得多说一句：按频率排，是为了让排查从概率最高的地方开
始。遇到“机器不工作”，先做的三件事永远在表的前三行——捏一捏芯片、
摸一摸温度、量一量各片的 VCC 脚——全程不过两分钟，却能拦下一大半
的故障。排在后面的几类各有鲜明的指纹：全亮是争用、多拍是消抖、成对
错位是线序、空闲随机是下拉、过暗烧灯是限流、提速才错是去耦，指纹对
上了就直接用对应的定位法，不必从头扫表。真正要警惕的是那些指纹不鲜
明的症状，它们多半属于第一行。

最后要说的是：这张表是攒出来的，不是背出来的。每解决一个故障，就把
它登记成一行——类别、症状、怎么找到的——攒过二十行，你会得到一张
属于自己的表，顺序和这张未必相同：手的习惯不同，故障的分布就不同，
常用长跳线的人，“接触不良”那一行会特别长。别人的表给的是起步的顺
序，自己的表给的才是定位的直觉。等到发现自己的表连续半年没有新增一
行，那不是故障消失了，是这张表毕业了——可以换 PCB 的表开始攒了。

这十类故障在 PCB 的世界里换了一副面孔，但一类不少地都在：接触不良
变成虚焊与冷焊点，芯片插反变成丝印方向看错，半板未接变成电源平面分
割不当；悬空输入还是悬空输入，三态争用还是总线争用，消抖照样要消抖，
去耦只会比面包板更讲究。差别只在定位工具：面包板上用镊子和手指，PCB
上用放大镜和示波器；而“先电源、后时钟、再信号，先静态、后动态”的
顺序不换。把这张表在面包板上练熟，等于提前攒下了 PCB 时代的定位直
觉——面包板因此是最便宜的故障训练场，这里每一类故障的学费，不过一
根跳线。

\topic{专题：把教学机的规格写成一页纸}把整台机器的全部约定压缩进一页纸，一行一条，每条都能用一次测量核对。
规格不是散文，是可核对的清单——“时钟要快”不是规格，“0.7--48Hz”
才是；“内存够用”不是规格，“16 字节”才是。

\tablechunkintro{1}{2}{总线}{控制线}
\begin{longtable}{L{0.16\linewidth}L{0.72\linewidth}}
\toprule
项目 & 一句话规格 \\
\midrule
总线 & 8 位三态共享，任何时刻只有一个驱动者；各线 100k$\Omega$ 下拉，空闲拍读 \code{0x00} \\\tablerowrule
地址空间 & 16 字节，4 位地址；程序与数据同处一块 RAM，不需要译码器 \\\tablerowrule
寄存器 & A、B、IR、OUT 为 8 位；MAR、PC 为 4 位；Flags 存 CF、ZF 两位 \\\tablerowrule
ALU & 只做 A+B 与 A$-$B；减法是 B 取反加一，SU 一根线兼管取反与进位输入 \\\tablerowrule
ISA & 11 条指令（含 NOP），每条一字节，高 4 位 opcode、低 4 位地址或立即数：NOP、LDA、ADD、SUB、STA、LDI、JMP、JC、JZ、OUT、HLT \\\tablerowrule
T 状态 & 每条指令固定走满 T0--T7 八拍；T0、T1 公共取指，用剩的拍控制字全 0 \\\tablerowrule
控制线 & 16 根：HLT、MI、RI、RO、II、IO、AI、AO、BI、$\Sigma$O、SU、CE、CO、J、FI、OI；“出”使能互斥，“入”使能不限 \\
\bottomrule
\end{longtable}

\tablechunkintro{2}{2}{控制实现}{自检签名}
\begin{longtable}{L{0.16\linewidth}L{0.72\linewidth}}
\toprule
项目 & 一句话规格 \\
\midrule
控制实现 & 微码 ROM 共 128 行乘 16 位，地址为 opcode 乘 8 加 T；JC、JZ 的 J 经 CF、ZF 选通 \\\tablerowrule
复位行为 & 异步拉入、同步释放；清 PC、T、IR、Flags；RAM 内容不受复位影响 \\\tablerowrule
时钟 & 自动挡 0.7--48Hz（555 振荡，C=10$\mu$F）与单步挡两路，面板开关选择；HLT 有效时钳停时钟；调试区间 1--10Hz \\\tablerowrule
编程模式 & PROG 开关切到 DIP 拨码：4 位地址、8 位数据、写按钮，逐字节装入；先停机后切换，切回复位再启动 \\\tablerowrule
输出 & OUT 寄存器锁存 A，8 只 LED 二进制直读；停机后读数保持稳定 \\\tablerowrule
自检签名 & 程序一输出 1、3、6、10，程序二输出 3、6、9、12；序列对，则整机对 \\
\bottomrule
\end{longtable}

写这页纸的方法只有三条：一行只写一件事；每件事都给出数值或名字；凡
给不出数值的行，就是设计还没做完的行。“时钟”那一行若写不出
0.7--48Hz，说明 555 的阻容账还没算；“复位行为”若写不出清哪几个寄
存器，说明复位链还没画完。规格先于搭建存在，它的漏洞在写的那一刻就
会暴露，而不是等到调试的第三个小时——这是写规格最实际的回报。还有
一条隐含的要求：不许写形容词。“快”“够”“稳定”这类词在规格里没
有位置，它们各自对应的数值才有位置。

写规格与读规格，是同一种功夫的两个方向。读一片陌生芯片的数据手册，
本质上是读别人机器的一页规格书：先找位宽、地址空间、控制极性、时序
窗口这几行，找到就能上手，找不到就知道该补什么实验。反过来，能把自
己的机器写成一页纸的人，读手册时自然会按同样的栏目对号——栏目是通
用的，只是别人的机器比你这台大几个数量级。日后改动这台机器时，顺序
也由这页纸定死：先改纸上对应的那一行，再动面包板上的线，最后按四级
检查回归。规格是设计者留给后来的自己的约定：它不在乎你隔了多久回来，
只要求每一条都还经得起测量。

这种“读规格”的功夫，本章其实已经练过一次。“内存与输出模块”一节读 74HC189 数据手
册时，满页参数里只取了三行：读路径的延迟、写路径的窗口、各信号的先
后约束——因为对本机而言，手册的其余各行都能用“一拍是毫秒级”这一
句话豁免。会写规格的人，才知道手册里哪些行是要紧的：手册不过是芯片
厂写给你的规格书，它的一页纸版本，就藏在那三行里。写与读在这里是同
一件事的两端：写的一方负责把约定压缩成可核对的行，读的一方负责把行
还原成可动手的约定。

这页纸的第三种用法是当改动的靶子：任何扩展念头，写在纸面上就是划掉
某一行、写上另一行。想把 RAM 扩到 256 字节，“地址空间”那一行从 16
改成 256，紧接着会发现 ISA、T 状态、编程模式三行都跟着失效——连锁
反应在纸面上先走一遍，比在线路上走出来便宜得多。规格页越短，这种推
演越彻底，这正是一页纸的价值：它让“改一处牵动几处”这件事，在动手
之前就能数清楚。

\topic{专题：从这台机器走向自己的设计}把决策画成树，根节点只有一个问题——
这台机器接下来要教你什么？向下分三个杈：指令集、时间轴、软件。

\begin{longtable}{L{0.13\linewidth}L{0.37\linewidth}L{0.15\linewidth}L{0.23\linewidth}}
\toprule
学习目标 & 推荐扩展顺序 & 预估工作量 & 这条路教会你什么 \\
\midrule
学 ISA 设计 & 先加 ADI（只烧两行微码），再改两字节指令格式（微码全表重排），最后 CALL/RET 与栈（新增 SP 模块） & 一个晚上、一个周末、一周以上，逐档倍增 & 一条指令的成本分布在哪里：改表便宜、改格式贵、加通路最贵 \\\tablerowrule
学总线时序 & 先拧高时钟、重算关键路径，再换更快的存储或加等待状态，最后用示波器核对建立与保持时间 & 几个晚上，外加一台示波器 & 频率上限是算出来的，不是试出来的；余量要留给分布参数 \\\tablerowrule
学系统编程 & 先加串口输出（结果不再靠数 LED），再扩 RAM 到 256 字节，最后写一个 monitor 让机器自己装程序 & 数周，三条路里最重 & 硬件的尽头是软件：让程序代替拨码开关伺候程序 \\
\bottomrule
\end{longtable}

三条路不互斥，但不要并行。这台机器的全部调试手段——四级检查、一页
规格、逐拍单步——都建立在“一次只改一处”的前提上；同时扩两处，故
障签名立刻失去意义，你会在自己最熟悉的机器上，体会陌生机器的调试难
度。正确的节奏是：选一个目标，走完它的一整条顺序，回归检查、更新规
格页，然后才允许想下一个。每走完一轮，机器变大一点，而它的可观察性
不许减少——这是扩展的及格线，也是这台机器作为实验平台的本分。

工作量那一栏也值得看一眼：三条路的量级差着不止一档，这不是编排失误，
而是三个学科本身的分量。ISA 的第一课便宜，是因为微码把“改指令”变
成了“改表”；时序课贵一些，因为它要求添置仪器、学会测量；系统编程
课最贵，因为它实际上是两门课——先扩硬件，再写软件。按自己能投入的
时间选路，比仅凭兴趣勉强坚持更有助于完成学习路径。

回到根节点那个问题，它有两种常见的错误答法，值得提前识破。一是“先
加最酷的”：中断、串口、流水线，哪个名词响亮先上哪个，结果每台机器
都半途而废；二是“先加最大的”：RAM、指令、时钟一次全改，第一周就淹
没在自己制造的故障里。两种答法的共同毛病，是把扩展当成了收藏而不是
学习。正确的答法只有一个句式：“我想弄懂 X，所以先改 Y。”这个句子
填不出来，就说明还没选好路——留在原地再玩一个星期这台机器，比贸然
开始实施更有收获。

所以，把教学计算机项目中的这台机器当作什么，决定了它还能给你什么。当作一个完成
品，它的使命在上电那一刻就结束了；当作一个最小可改动的实验平台，它
才刚开始工作。它够小，小到任何改动都能在一页规格上找到对应行；它够
完整，完整到 ISA、时序、软件三条路都走得通；它够透明，透明到每次改
动是对是错，总线 LED 当晚就告诉你。第一台机器是照着本章搭的；第二
台机器不必是——从划掉那页规格书的某一行、写上你自己的一行开始，那
就是你的设计的开头。




























\section{六、把教学机边界映射回现代整机}

面包板教学机把 PC、寄存器、ALU、存储、控制字和输出锁存器放在可直接观察的位置。现代处理器扩大了规模并加入流水线、Cache、虚拟内存和设备队列，但边界没有消失：教学机的总线授权对应片上互连仲裁，RAM 读周期对应 Cache/DRAM 请求，输出寄存器对应 MMIO 设备状态，微码步进对应控制状态推进。理解现代机器的关键不是记住更多模块，而是继续问同一组问题：谁发起、谁拥有状态、何时采样、何时完成、失败怎样返回。

\begin{longtable}{L{0.20\linewidth}L{0.30\linewidth}L{0.40\linewidth}}
\toprule
教学机可见边界 & 现代整机对应物 & 不变的检查问题 \\
\midrule
PC 与取指 ROM & 分支预测前端、I-Cache、页表和主存 & 下一 PC 从何而来，取指错误归属哪条指令 \\\tablerowrule
共享 8-bit 总线 & crossbar/NoC、协议桥和 PCIe & 哪个发起者获授权，背压时谁保持负载 \\\tablerowrule
RAM 读写脉冲 & Cache line、DDR burst、块 I/O 与 DMA & 请求何时接受，数据何时可见，错误怎样闭合 \\\tablerowrule
输出锁存器 & MMIO 寄存器、doorbell、设备队列 & 写返回是否等于设备副作用完成 \\\tablerowrule
微码计数器 & 流水控制、队列状态机和固件阶段 & 当前状态以哪个边沿推进，复位从何处恢复 \\\tablerowrule
LED/示波器测点 & trace buffer、性能计数器、AER/RAS 日志 & 证据来自哪个时钟域，是否保留首次故障 \\
\bottomrule
\end{longtable}

\section{七、一次程序首次运行的全链路时间线}

考虑冷启动后首次运行一个读取文件、处理数据并显示结果的程序。电源控制器先使各电压轨进入允许范围，时钟和复位网络建立已知状态；固件从受保护入口取指，训练 DRAM，枚举 PCIe 与 NVMe，发布内存和中断拓扑，再把控制权交给装载器与内核。到这一刻为止，硬件完成的是“让软件拥有可靠的执行和设备发现环境”，还没有执行目标程序。

进程启动时，装载器解析 ELF 段和重定位，建立虚拟地址空间，把尚未实际读入的代码页记录为文件映射。CPU 首次取指触发缺页：页表遍历发现映射未驻留，异常入口保存精确 PC，内核把文件偏移转换为块请求，NVMe 驱动写 submission queue 和 doorbell。SSD 控制器从 NAND 取得页面，经 PCIe DMA 写主存，写 completion queue 并发送 MSI-X；内核验证完成状态、更新页表并失效必要的 TLB 项，返回异常指令。相同 PC 再次取指，这次通过 TLB、Cache 和 DRAM 获得机器码。

程序读取输入文件时，页缓存命中可能让系统调用不再访问设备；未命中则重走块层和 DMA。处理循环在 CPU 流水线中执行，load/store 经 Cache 与一致性互连交换数据。若结果交给 GPU，CPU 先发布资源和命令，GPU 经设备地址转换读取输入，在 VRAM 或共享内存中计算，写结果并更新 fence；显示引擎只在翻页边界采用完成帧。若程序还写日志，code{write()} 返回通常只说明数据进入内核缓存，只有文件系统提交、设备 flush 与介质持久化边界全部满足后，掉电重启才应看到记录。

\begin{longtable}{L{0.14\linewidth}L{0.24\linewidth}L{0.31\linewidth}L{0.21\linewidth}}
\toprule
阶段 & 发起者与目标 & 完成边界 & 失败证据 \\
\midrule
上电启动 & 电源/固件到 CPU、DRAM、互连 & 内存可用、设备已枚举、内核入口有效 & 复位原因、训练码、固件阶段 \\\tablerowrule
程序装载 & 装载器到页表与文件映射 & 入口地址和初始栈成为进程状态 & ELF/重定位错误、映射冲突 \\\tablerowrule
首次取指 & CPU 到 TLB、页表、存储栈 & 缺页处理后同一 PC 可精确重试 & fault 地址、PTE、块请求 ID \\\tablerowrule
设备读取 & 驱动到 NVMe/SSD 与主存 & CQ 成功且 DMA 数据对 CPU 可见 & SQ/CQ 指针、状态码、IOMMU fault \\\tablerowrule
GPU 处理 & CPU 队列到 GPU 和显存 & 结果写入先于 fence 完成可见 & 命令读指针、GPU fault、fence \\\tablerowrule
显示输出 & 显示引擎到帧缓冲和链路 & 新帧被采用并扫描输出 & 翻页序号、扫描位置、链路状态 \\\tablerowrule
持久化日志 & 文件系统到控制器和介质 & flush/提交协议规定的持久化点 & 日志序号、设备错误、掉电恢复记录 \\
\bottomrule
\end{longtable}

\section{八、用唯一事务身份连接跨层证据}

跨层排错最困难的地方是每层使用不同名字：CPU 有 PC 和异常号，内核有进程、页和 bio，请求队列有 tag，PCIe 有 requester/tag，NVMe 有 queue ID/command ID，GPU 有 context/queue/fence。不能要求所有层使用同一个位宽相同的全局 ID，但可以在进入下一层时记录映射关系和时间戳，使一个上层请求能够沿日志追到下层完成。

时间戳本身也需要校准。CPU 核、设备、网卡和示波器可能使用不同频率与起点；先用共同事件建立偏移和漂移关系，再比较先后。若只按日志打印顺序排列，缓冲和异步刷新可能让较晚发生的事件先出现。可靠 trace 应同时保存事件来源、局部序号、时钟域、请求身份和状态变化，而不只保存一句自然语言说明。

观测还会改变系统。开启详细总线 trace 可能增加缓冲压力，串口打印可能延长中断，示波器探头会增加电容。排错记录应注明观测配置，并优先选择已有硬件计数器、只读状态和触发式缓冲。需要插入探针时，先估算负载和时序影响，再用开启/关闭观测的对照实验确认现象没有被测量手段掩盖。

\section{九、从症状定位最早失效边界}

“程序卡住”只是最终症状。定位应从最后一个已确认完成的边界向前后逐层缩小范围：若 PC 仍退休指令，处理器没有完全停机；若 NVMe SQ 尾指针增长而设备头指针不动，问题位于 doorbell、队列可见性或设备取命令；若 CQ 已更新但中断计数不变，检查 MSI-X 路由和屏蔽；若 fence 已完成而画面未变，转向翻页与扫描域。每一步都用状态转换缩小范围，不从软件表象直接跳到更换硬件。

恢复动作也必须与故障域相称。重试适合瞬态传输错误；功能级复位适合单设备控制状态失步；链路复训适合物理连接暂时失锁；整机复位只在共享状态已无法隔离时使用。复位前先停止新流量并保留证据，复位后重建地址、翻译、队列和中断状态，再用一笔已知事务验证恢复。只看到服务重新工作而不验证旧 DMA、迟到完成和持久化状态，不能证明恢复闭环。

\begin{keyidea}
全书所有机制最终都可以归入七个问题：输入是什么，状态在哪里，数据走哪条路，控制为何这样选择，哪个边界使结果可见，失败如何返回，以及用什么证据证明恢复。能沿一次真实程序逐层回答这七问，就已经建立了从器件到整机的软件可见硬件通路。
\end{keyidea}

\section*{章末练习}

\begin{chapterexercise}{综合练习：总线互斥检查至时钟两路}
\begin{exercisesubproblem}{（a）总线互斥检查}
写出 T0、T1 两拍全部有效的“出”使能与“入”使能，说明为什么不构成争用；再指出若微码把 \code{RO} 与 $\Sigma$O 同拍置 1 会发生什么
\exercisecheck{学习目标}{总线驱动者唯一}
\end{exercisesubproblem}
\begin{exercisesubproblem}{（b）时钟两路}
说明自动时钟（555）与单步按钮之间切换的时机要求，以及按钮为什么必须消抖——不消抖会发生什么
\exercisecheck{学习目标}{一拍不多一拍不少}
\end{exercisesubproblem}
\end{chapterexercise}

\begin{chapterexercise}{综合练习：微码读法至累加程序改编}
\begin{exercisesubproblem}{（a）微码读法}
微码 ROM 地址为 \{opcode, T\}：查 \code{ADD}（opcode \code{0x2}）在 T3、T4 两拍的有效控制线，并写出这两个 ROM 单元的地址（二进制）
\exercisecheck{学习目标}{会按 opcode 拼 T 查表}
\end{exercisesubproblem}
\begin{exercisesubproblem}{（b）累加程序改编}
只改程序一数据区的一个字节，把“从 1 加到 4”改成“从 1 加到 2”（结果为 3）：指出地址与新值，并写出新的输出序列
\exercisecheck{学习目标}{程序与数据同 RAM}
\end{exercisesubproblem}
\end{chapterexercise}

\begin{chapterexercise}{综合练习：故障反推至扩展设计}
\begin{exercisesubproblem}{（a）故障反推}
症状：复位后单步，第一拍有效控制线是 \code{IO+MI} 而不是 \code{CO+MI}，之后执行全错。推断故障点并给出修复
\exercisecheck{学习目标}{复位覆盖全部时序状态}
\end{exercisesubproblem}
\begin{exercisesubproblem}{（b）扩展设计}
把 RAM 扩到 32 字节（5 位地址）：列出至少三处必须同步修改的地方，并说明哪一处改动最大
\exercisecheck{学习目标}{地址位宽牵动全局}
\end{exercisesubproblem}
\end{chapterexercise}

\begin{chapterexercise}{边界映射}
把教学机输出锁存器映射到现代设备事务
\exercisecheck{检查要点}{MMIO 写接受不等于设备完成}
\end{chapterexercise}

\begin{chapterexercise}{首次取指}
写出文件映射代码页首次取指的完整缺页 trace
\exercisecheck{检查要点}{同一 PC 在状态修复后精确重试}
\end{chapterexercise}

\begin{chapterexercise}{GPU 显示}
区分命令提交、GPU fence 和显示扫描完成
\exercisecheck{检查要点}{三个边界不能合并}
\end{chapterexercise}

\begin{chapterexercise}{持久化}
解释 \code{write()} 返回与掉电后可见的差别
\exercisecheck{检查要点}{还要经过文件系统和设备持久化协议}
\end{chapterexercise}

\begin{chapterexercise}{证据关联}
为 NVMe 读列出五层事务身份及映射记录
\exercisecheck{检查要点}{不要求同一 ID，但必须可追踪}
\end{chapterexercise}

\begin{chapterexercise}{故障定位}
CQ 已更新但线程仍等待，给出由近到远的测点
\exercisecheck{检查要点}{先查中断、完成消费和唤醒路径}
\end{chapterexercise}

\begin{chapterexercise}{恢复验收}
设计设备局部复位后的最小验证序列
\exercisecheck{检查要点}{重建映射/队列并拒绝迟到完成}
\end{chapterexercise}

\section*{参考解答与要点}

\begin{chapteranswer}{综合练习：总线互斥检查至时钟两路}
\begin{answersubproblem}{（a）总线互斥检查}
T0 的“出”使能只有 \code{CO}（PC 驱动总线），“入”使能只有 \code{MI}；T1 的“出”使能只有 \code{RO}，“入”使能是 \code{II} 与 \code{CE}——\code{CE} 让 PC 内部计数，并不驱动总线。两拍都只有一个驱动者，故无争用。若 \code{RO} 与 $\Sigma$O 同拍为 1，RAM 与 ALU 同时驱动总线，线上电平由两边输出级对抗决定，LED 显示既非 RAM 值也非 ALU 值的中间态，总线电流猛增，时间一长会损坏输出级
\answercheck{必须掌握的要点}{“出”使能必须互斥，“入”使能不限}
\end{answersubproblem}
\begin{answersubproblem}{（b）时钟两路}
只在停机或按住复位按钮时换挡——运行中切换会在时钟上切出窄毛刺；换完从复位后的 T0 重新单步几拍确认再继续。机械触点抖动持续数毫秒，不消抖则按一次按钮产生多个时钟沿，T 计数器连跳数拍、微操作序列被跳步，而且这种错误无法稳定复现。用 RS 触发器或施密特触发器整形，保证一次按键恰好给出一个完整时钟周期
\answercheck{必须掌握的要点}{单步的价值在于拍数确定}
\end{answersubproblem}
\end{chapteranswer}

\begin{chapteranswer}{综合练习：微码读法至累加程序改编}
\begin{answersubproblem}{（a）微码读法}
\code{ADD} 的 T3：\code{RO+BI}，把 RAM[addr] 经总线装入 B，ROM 地址为 \code{0010011}（高 4 位 opcode \code{0010}，低 3 位 T \code{011}）；T4：$\Sigma$O+\code{AI+FI}，A ← A+B 并锁存 ZF、CF，ROM 地址为 \code{0010100}。两行其余控制线全为 0
\answercheck{必须掌握的要点}{微码 ROM 是“指令 × 拍”的查找表}
\end{answersubproblem}
\begin{answersubproblem}{（b）累加程序改编}
把地址 13 的上限常数由 \code{0x04} 改为 \code{0x02}，其余 15 字节原样。循环在 i=2 时 \code{SUB 13} 结果为 0、ZF=1 而退出，OUT 依次显示 1、3，最终和为 3。数据字节改了，程序行为就改了
\answercheck{必须掌握的要点}{循环边界是数据不是代码}
\end{answersubproblem}
\end{chapteranswer}

\begin{chapteranswer}{综合练习：故障反推至扩展设计}
\begin{answersubproblem}{（a）故障反推}
T 计数器没接复位：PC 被清到 0，T 却停在复位前的 T2（\code{IO+MI} 正是访存类指令 T2 的组合），第一条指令从 T2 开始错位执行，取指序列整个乱掉。修复：把复位信号同时接到 T 计数器（74HC161）的清零端；改后复位再单步，控制线序列必须从 \code{CO+MI} 起逐拍与微码表一致
\answercheck{必须掌握的要点}{复位要清 PC、T、IR、标志}
\end{answersubproblem}
\begin{answersubproblem}{（b）扩展设计}
至少三处：①MAR 由 4 位加宽到 5 位；②RAM 换成 32×8，编程模式拨码覆盖第 5 位；③指令格式：低 4 位放不下 5 位地址——要么改两字节指令（取指后加一拍读地址字节），要么约定程序只用低 16 字节、高 16 字节纯作数据。改动最大的是③：微码 T 序列与本章全部编码表都要跟着重排
\answercheck{必须掌握的要点}{地址位宽是全局约定，牵一发动全身}
\end{answersubproblem}
\end{chapteranswer}

\begin{chapteranswer}{边界映射}
输出锁存器对应 MMIO 寄存器或 doorbell；总线采样沿只证明写被目标接口接受，设备内部状态机可能稍后才完成动作。
\answercheck{必须保留的检查点}{接受、完成和副作用分开}
\end{chapteranswer}

\begin{chapteranswer}{首次取指}
CPU 取指触发缺页，内核按文件偏移发块请求，设备 DMA 并完成，内核更新 PTE/TLB 后返回，原 PC 再次取指。
\answercheck{必须保留的检查点}{异常指令不提前提交}
\end{chapteranswer}

\begin{chapteranswer}{GPU 显示}
doorbell 只发布命令，fence 表示 GPU 结果达到规定可见点，显示完成还需翻页采用并扫描。
\answercheck{必须保留的检查点}{每个域读取独立状态}
\end{chapteranswer}

\begin{chapteranswer}{持久化}
\code{write()} 可在页缓存接收后返回；持久化还依赖日志/元数据顺序、设备 flush 和介质确认。
\answercheck{必须保留的检查点}{完成语义由接口层级定义}
\end{chapteranswer}

\begin{chapteranswer}{证据关联}
可记录进程/页或 bio、块层 tag、NVMe queue/command ID、PCIe requester/tag 与 DMA 地址，并在层间保存映射及时间戳。
\answercheck{必须保留的检查点}{关联关系比统一编号更重要}
\end{chapteranswer}

\begin{chapteranswer}{故障定位}
核对 CQ phase/status、驱动头指针、MSI-X 向量与屏蔽、中断控制器计数、软中断/完成回调和等待队列唤醒。
\answercheck{必须保留的检查点}{CQ 到达后路径仍未结束}
\end{chapteranswer}

\begin{chapteranswer}{恢复验收}
停止提交并保存状态，隔离 DMA/中断，复位后重建 BAR/IOMMU/队列，递增 generation，发一笔已知请求并验证数据、完成和旧事件隔离。
\answercheck{必须保留的检查点}{服务恢复且旧状态不可污染新事务}
\end{chapteranswer}
